\documentclass[11pt]{article}
\usepackage[margin=1.05in]{geometry}
\usepackage{lmodern}
\usepackage{amsmath,amssymb,amsthm,mathtools}
\usepackage{xcolor}
\usepackage[hidelinks]{hyperref}
\usepackage{microtype,mathrsfs}
\newtheorem{theorem}{Theorem}
\newtheorem{proposition}{Proposition}[section]
\newtheorem{lemma}[proposition]{Lemma}
\newtheorem{corollary}[proposition]{Corollary}
\newcommand{\E}{\mathbb E}
\newcommand{\R}{\mathbb R}
\newcommand{\Z}{\mathbb Z}
\newcommand{\Cov}{\operatorname{Cov}}
\newcommand{\Loc}{\operatorname{Loc}}
\newcommand{\diag}{\operatorname{diag}}
\DeclarePairedDelimiter{\norm}{\lVert}{\rVert}
\title{Fractional spread-out Ising models at their upper critical dimension}
\author{Jiwoon Park}
\date{\today}
\begin{document}
\maketitle

\begin{abstract}
We give a renormalisation-group argument for the critical two-point
function of the spread-out Ising model with kinetic operator
\((-\Delta_R)^{d/4}\), for integer \(1\le d\le4\).
An auxiliary Hilbert-space realization supplies positive finite-range
pieces for the Hessian-absorbed fractional covariance. Their natural
postmass power is two. Affine and quotient-regulator estimates yield a
uniform contraction on an enlarged activity ball; the source weights
are stopped at coalescence. Mass tuning, an exact susceptibility
identity, and a summable renewal envelope identify the endpoint with
the physical critical point. The critical free correlation is
asymptotic to a positive constant times \(|x|^{-d/2}\).
\end{abstract}

\section*{Fractional model and main statement}
Let \(d\in\{1,2,3,4\}\), let \(R\) be a sufficiently large integer, and set
\[
 D_{R,d}(x)=\frac{\boldsymbol1_{\{0<|x|_\infty\le R\}}}
                  {(2R+1)^d-1},\qquad
 L_R=-\Delta_R=I-D_{R,d},\qquad
 \alpha=\frac d4,\qquad s=2\alpha=\frac d2.
\]
The fractional power is defined by its Fourier multiplier
\(\widehat{L_R^\alpha}(k)=(1-\widehat D_{R,d}(k))^\alpha\)
on \([-\pi,\pi]^d\). The ferromagnetic model intended here has weight
\begin{equation}\label{frac:weight}
 \exp\{-\tfrac\tau2\langle\sigma,L_R^\alpha\sigma\rangle\}.
\end{equation}
Thus \(L_R^\alpha\) is the kinetic operator: the nonnegative pair
interaction is \(-L_R^\alpha(x)\) for \(x\ne0\). Its diagonal contributes
only a spin-independent constant. On finite sets we use the restriction
of this infinite-volume pair interaction, and on tori its periodization.
Write \(\tau_c=\tau_c(d,R)\) for the spontaneous-magnetization critical
inverse temperature and \(\langle\cdot\rangle^f_\tau\) for the
infinite-volume free state. Throughout, \(|x|\) is the Euclidean norm.

\begin{theorem}[Critical fractional two-point function]
\label{frac:main}
For each \(d\in\{1,2,3,4\}\) and all sufficiently large fixed \(R\),
the model \eqref{frac:weight} has critical free two-point asymptotic
\begin{equation}\label{frac:target}
 \langle\sigma_0\sigma_x\rangle^f_{\tau_c}
       \sim C_{d,R}|x|^{-d/2},\qquad C_{d,R}>0,
       \qquad |x|\longrightarrow\infty.
\end{equation}
\end{theorem}

\paragraph{Organization and RG conventions.}
Part~\ref{part:fractional} constructs the fractional kernel, positive
covariance pieces, regulated map, tuned massive trajectories and source
limits. The finite-range RG extraction and circle-product algebra is
the one specified in Part~\ref{part:fourdim}; the new dimension-dependent
norm estimates are proved explicitly below. The local kinetic coupling
is irrelevant for \(d<4\) and is evolved forward. Only the mass is
physically tuned in those dimensions. Part~\ref{part:fourdim} retains
the four-dimensional specialization, where mass and marginal stiffness
are tuned together. The warning at the top applies to this proof draft.

\part{Fractional construction and flow}\label{part:fractional}

\section{The fractional operator and its ferromagnetic pair interaction}
\label{frac:operator}

Let \(d\in\{1,2,3,4\}\), \(\alpha=d/4\), \(s=2\alpha=d/2\), and
\[
 D_{R,d}(x)=\frac{\boldsymbol1_{\{0<|x|_\infty\le R\}}}
                         {(2R+1)^d-1},\qquad
 L_R=I-D_{R,d}=-\Delta_R.
\]
The fractional power is defined by the spectral calculus:
\(\widehat{L_R^\alpha}(k)=(1-\widehat D_{R,d}(k))^\alpha\).
The ferromagnetic convention is the spin weight
\(\exp[-\tau\langle\sigma,L_R^\alpha\sigma\rangle/2]\).
Thus the pair coupling is the \emph{negative off-diagonal} kernel of
\(L_R^\alpha\). A positive sign in front of this quadratic form
would define a different model.

For \(0<\alpha<1\), put
\[
 c_n=(-1)^{n+1}\binom\alpha n
 =\frac{\alpha\Gamma(n-\alpha)}
        {\Gamma(1-\alpha)\Gamma(n+1)}>0.
\]
The binomial series converges absolutely on the spectral interval,
including its endpoints, since \(\sum_{n\ge1}c_n=1\). Hence
\begin{equation}\label{frac:binomial}
 J_R^{\rm raw}=I-L_R^\alpha
       =\sum_{n\ge1}c_nD_{R,d}^{*n}
\end{equation}
is a nonnegative, symmetric probability kernel. Define
\[
 \kappa_R=L_R^\alpha(0)=1-J_R^{\rm raw}(0),\qquad
 K_R(x)=\boldsymbol1_{\{x\ne0\}}J_R^{\rm raw}(x),\qquad
 D_R^{(\alpha)}=K_R/\kappa_R.
\]
For \(\alpha=1\) use \(J_R^{\rm raw}=D_{R,d}\) and
\(\kappa_R=1\). In every case \(D_R^{(\alpha)}\) is a
nonnegative, zero-diagonal probability kernel and
\begin{equation}\label{frac:normalization}
 \exp[-\tfrac\tau2\langle\sigma,L_R^\alpha\sigma\rangle]
 =\text{constant}\;
   \exp[\tfrac{\tau\kappa_R}{2}
                     \langle\sigma,D_R^{(\alpha)}\sigma\rangle].
\end{equation}
In particular the normalized inverse temperature is
\(\widetilde\beta=\kappa_R\tau\), and the critical inverse
temperatures satisfy
\(\tau_c=\widetilde\beta_c/\kappa_R\).
Also \(\|D_{R,d}^{*n}\|_\infty\le\|D_{R,d}\|_\infty\)
for \(n\ge1\), so
\[
 0\le1-\kappa_R\le C_dR^{-d},\qquad
 \|K_R\|_\infty\le C_dR^{-d}.
\]
These identities explain precisely which pair interaction is meant
by replacing the kinetic operator with \((-\Delta_R)^{d/4}\).

\section{The Gaussian asymptotic and the long-range obstruction}
\label{frac:gaussian}

Set
\[
 a_{R,d}=\frac12\sum_x x_1^2D_{R,d}(x)
 =\frac{(2R+1)^d}{(2R+1)^d-1}\frac{R(R+1)}6.
\]
Hypercubic symmetry and finite range give
\begin{equation}\label{frac:symbol}
 1-\widehat D_{R,d}(k)=a_{R,d}|k|^2+O_R(|k|^4),\qquad
 \widehat{L_R^\alpha}(k)
   =a_{R,d}^{\alpha}|k|^s
           \bigl(1+O_R(|k|^2)\bigr).
\end{equation}
Here \(a_{R,d}\asymp R^2\). Since \(0<s<d\), the massless
Gaussian covariance \(G_{R,0}=L_R^{-\alpha}\) is well-defined
as a Fourier integral.

\begin{proposition}[Fractional Gaussian two-point function]
\label{frac:green-prop}
For each fixed \(R\) and \(d\in\{1,2,3,4\}\),
\begin{equation}\label{frac:green}
 G_{R,0}(x)\sim
 \frac{\Gamma((d-s)/2)}{2^s\pi^{d/2}\Gamma(s/2)}
 a_{R,d}^{-\alpha}|x|^{-(d-s)}
 =\frac{a_{R,d}^{-d/4}}{(2\pi)^{d/2}}|x|^{-d/2}.
\end{equation}
For \(d<4\), the pair interaction has the asymptotic
\begin{equation}\label{frac:tail}
 K_R(x)\sim c_{d,s}a_{R,d}^{\alpha}|x|^{-d-s},\qquad
 c_{d,s}=\frac{s2^{s-1}\Gamma((d+s)/2)}
                    {\pi^{d/2}\Gamma(1-s/2)}>0.
\end{equation}
Consequently \(\sum_x|x|^tK_R(x)<\infty\) precisely for
\(0\le t<s\). In particular the interaction has neither an
exponential moment nor a finite second moment when \(d<4\).
\end{proposition}

\begin{proof}
Let \(p_t(x)=(e^{-tL_R})(0,x)\), the continuous-time walk kernel.
The spectral integral identities are
\[
 L_R^{-\alpha}=\frac1{\Gamma(\alpha)}
       \int_0^\infty t^{\alpha-1}e^{-tL_R}\,dt,
\]
and, for \(0<\alpha<1\),
\[
 K_R(x)=\frac{\alpha}{\Gamma(1-\alpha)}
       \int_0^\infty p_t(x)t^{-1-\alpha}\,dt
       \qquad(x\ne0).
\]
The finite-range walk has diffusion coefficient \(a_{R,d}\).
Fourier inversion of \(e^{-t(1-\widehat D_{R,d})}\), splitting
into a neighborhood of zero and its complement, gives the local
central limit kernel
\[
 (4\pi a_{R,d}t)^{-d/2}\exp[-|x|^2/(4a_{R,d}t)],
\]
in the regime \(t\asymp|x|^2\),
with the usual Gaussian and bounded-jump off-diagonal bounds
outside that regime. These bounds give integrable majorants in
the two displayed integrals after the change \(t=|x|^2u\).
Integration of the Gaussian kernel gives respectively
\[
 \frac{\Gamma(d/2-\alpha)}
       {4^\alpha\pi^{d/2}\Gamma(\alpha)}a_{R,d}^{-\alpha},
 \qquad
 \frac{\alpha4^\alpha\Gamma(d/2+\alpha)}
       {\pi^{d/2}\Gamma(1-\alpha)}a_{R,d}^{\alpha}.
\]
These are the constants in \eqref{frac:green} and
\eqref{frac:tail}. The moment assertion follows by summing
\(|x|^{t-d-s}\) over lattice shells.
\end{proof}

The Gaussian calculation identifies the exponent expected in the
Ising generalization, but it does not itself prove that the
interacting critical two-point function has that exponent. The
quartic is marginal because the field dimension is
\((d-s)/2=d/4\), so that \(4(d-s)/2=d\).
The local sixth-order term has irrelevant eigenvalue \(L^{-s}\).
For \(d<4\), an ordinary local gradient-square term has eigenvalue
\(L^{s-2}<1\); it is not the marginal kinetic coordinate of the
four-dimensional argument.

\section{A finite-range reference split for the fractional model}
\label{frac:split}

Using the full long-range kernel as the high-temperature reference
would lose both its exponential tree norm and the bounded-degree
block graph. There is an exact split that avoids that particular
problem. For \(0<\alpha<1\), choose an even integer \(N\ge2\)
and define
\[
 Q_N=\sum_{1\le n<N}c_nD_{R,d}^{*n},\qquad
 P_N=\sum_{n\ge N}c_nD_{R,d}^{*n}.
\]
Then \(J_R^{\rm raw}=Q_N+P_N\), \(Q_N\ge0\) has range at most
\((N-1)R\), and its row sum is the fixed number
\(q_N=\sum_{n<N}c_n<1\). The diagonal of \(Q_N\) is
\(O(R^{-d})\). It may be left in the reference kernel, since it
only adds a constant to the Ising energy; a mean-field comparison
that retains this diagonal incurs an additional \(O(R^{-d})\)
self-edge error. Transferring it to the Gaussian part as a bare
\(O(R^{-d})\delta_0\) term would require a separate check of the
initial spatial derivative bounds.

\begin{proposition}[Positive Gaussian remainder]
\label{frac:positive-tail}
For even \(N\), \(P_N\) is positive semidefinite as a convolution
operator, as well as pointwise nonnegative. If
\(\varepsilon_R>0\), then \(P_N+\varepsilon_R I\) is strictly
positive definite. Adding \(\varepsilon_R I\) to the total pair
interaction is spin-independent. Moreover,
\begin{equation}\label{frac:tail-resolvent}
 P_N=\frac{\sin(\pi\alpha)}\pi
   \int_0^1 u^{N-\alpha-1}(1-u)^\alpha
              D_{R,d}^{*N}(I-uD_{R,d})^{-1}\,du.
\end{equation}
The integral is an operator-norm integral.
\end{proposition}

\begin{proof}
For \(f(t)=1-(1-t)^\alpha\), every derivative \(f^{(n)}(t)\)
is positive on \(t<1\). The remainder after degree \(N-1\) is
\[
 f(t)-\sum_{n<N}c_nt^n
 =\frac{t^N}{(N-1)!}
       \int_0^1(1-v)^{N-1}f^{(N)}(vt)\,dv\ge0
\]
for even \(N\) and \(-1\le t<1\); continuity covers \(t=1\).
Apply this to the spectrum of \(D_{R,d}\). Alternatively, the
beta integral for \(c_n\) is
\[
 c_n=\frac{\sin(\pi\alpha)}\pi
       \int_0^1u^{n-\alpha-1}(1-u)^\alpha\,du.
\]
Summing its geometric series gives \eqref{frac:tail-resolvent}.
The factor \(D_{R,d}^{*N}\) is positive semidefinite for even
\(N\); \((I-uD_{R,d})^{-1}\) is positive definite for
\(0<u<1\), and the commuting product is positive. Its operator
norm is bounded by \((1-u)^{-1}\), so integrability follows
from \(\alpha>0\). The statement about strict positivity is
immediate.
\end{proof}

One may choose \(\varepsilon_R=e^{-R}\). The same integral gives
\(|\widehat P_N(k)|\le C_{\alpha,N}|\widehat D_{R,d}(k)|^N\).
Thus arbitrarily high fixed orders of high-frequency smoothing
can be imposed by choosing \(N\) before \(R\). This split retains
a finite-range reference system with row sum below one. Section~\ref{complete:covariance} constructs the finite-range decomposition after exact Hessian absorption.

The normalized representation used below is
\[
 D^+=\frac{J_R^{\rm raw}+\varepsilon_R I}{1+\varepsilon_R},
 \qquad P=\frac{P_N+\varepsilon_R I}{1+\varepsilon_R},
 \qquad Q=\frac{Q_N}{1+\varepsilon_R}.
\]
Thus \(D^+=P+Q\), its row sum is one, and its inverse temperature
is \(\beta=(1+\varepsilon_R)\tau\). This diagonal addition and
temperature change leave the spin law exactly unchanged.
The reference row sum is below one in a fixed neighborhood of
\(\beta=1\). Its diagonal can be kept with the self-edge correction
noted above; all reference field changes remain finite-range.

\section{Subordination and the positive-mass tail}
\label{frac:subordination}

For \(0<\alpha<1\), \(\mu>0\), and \(t>0\), the Stieltjes
identity is
\begin{equation}\label{frac:stieltjes}
 \frac1{\mu+t^\alpha}
 =\int_0^\infty\frac{\rho_\mu(r)}{t+r}\,dr,\qquad
 \rho_\mu(r)=\frac{\sin(\pi\alpha)}\pi
   \frac{r^\alpha}
   {r^{2\alpha}+2\mu r^\alpha\cos(\pi\alpha)+\mu^2}.
\end{equation}
It follows either by the Cauchy integral formula along the negative
real axis or by the Stieltjes inversion formula. The limiting
\(\mu=0\) density is \((\sin(\pi\alpha)/\pi)r^{-\alpha}\).
Consequently
\[
 (\mu+L_R^\alpha)^{-1}
       =\int_0^\infty\rho_\mu(r)(L_R+r)^{-1}\,dr.
\]
Positive finite-range pieces of the local resolvent on the right
can be integrated separately; their ranges are unchanged and
positive semidefiniteness is preserved. The same observation
applies to \eqref{frac:tail-resolvent}, whose extra numerator is
finite-range and positive semidefinite. Spatial smoothing and
initial-scale bounds still have to be verified for the actual
covariance used in the Ising representation.

A positive auxiliary mass does \emph{not} give exponential decay
when \(d<4\). At fixed \(R\) and \(\mu>0\),
\begin{equation}\label{frac:massive-tail}
 (\mu+L_R^\alpha)^{-1}(x)
       \sim \mu^{-2}K_R(x)
       \asymp_{R,\mu}|x|^{-d-s}.
\end{equation}
For example, this follows by expanding its Fourier multiplier at
zero: its first nonconstant nonsmooth term is
\(-\mu^{-2}a_{R,d}^\alpha|k|^s\), whose inverse transform is
\(\mu^{-2}\) times the leading off-diagonal kernel in
\eqref{frac:tail}. The remaining homogeneous terms have strictly
higher order, and the multiplier is smooth away from zero.
Thus a finite-range covariance estimate with scaling
\(L^{-(d-s)j}\) cannot have arbitrarily high postmass powers
\((1+\mu L^{sj})^{-p}\): summing pieces whose range reaches
\(|x|\) would otherwise contradict \eqref{frac:massive-tail}.
The maximal tail-compatible power is \(p=2\). Similarly, an
argument identifying a susceptibility limit must use summable
algebraic domination, rather than periodic exponential decay.

The exact linear field change and Hessian absorption give a denominator
of the form
\[
 T=B M+A(K+z\mathcal D),
\]
where \(A\) contains \(P_N\) and is long-range. The scalar
operator \(T\) is not finite-range, even if \(K\) is replaced by a
moment-matched finite-range kernel. The auxiliary-space realization in
Section~\ref{complete:covariance} resolves this issue: locality is
imposed on a larger operator before compressing its covariance.

\section{An auxiliary-space construction of the absorbed fractional covariance}
\label{complete:covariance}
This section closes the covariance-decomposition step for
\(d\in\{1,2,3\}\). Put \(\alpha=d/4\), \(s=2\alpha\),
\(D=D_{R,d}\), and \(\ell(k)=1-\widehat D(k)\).
The local kinetic coordinate is not absorbed separately; its initial
value is zero and its subsequent evolution is irrelevant. We retain
all of the exponentially local Hessian correction in the covariance.
Use the even-tail split of Section~\ref{frac:split}, with a fixed
even integer \(N\) chosen sufficiently large below. Write
\[
 \zeta=\frac\beta{1+\varepsilon_R},\quad
 A=\zeta(P_N+\varepsilon_R I),\quad B=I-\zeta Q_N,
 \qquad \varepsilon_R=e^{-R}.
\]
If \(p=\sum P\), \(q=\sum Q\), set
\(c=(1-\mu-\beta q)/(\beta p)\), so that \((B-cA)\boldsymbol1
=\mu\boldsymbol1\). Suppose the residual Hessian kernel \(K\)
is real, hypercubic symmetric, has zero row sum, and satisfies
\begin{equation}\label{complete:K-assumption}
 \sum_x |K(x)|e^{\eta|x|/R}\le Cg,\qquad g=R^{-d}.
\end{equation}
Work in \(|\beta-1|+\mu\le C_0g\). All constants may depend
on the fixed even integer \(N\), the displayed constants, and the
fixed scale factor, but not on \(R,\mu,\beta\).

\begin{lemma}[Scalar positive-polynomial inverse decomposition]
\label{complete:scalar-polynomials}
Fix an integer \(L>1\). There are nonnegative polynomials \(q_j\)
on \([0,4]\), of degree at most \(L^{j+1}\), such that for \(x>0\),
\[
 \sum_{j\ge0}q_j(x)=x^{-1},\qquad
 q_j(x)\le C_pL^{2j}(1+L^{2j}x)^{-p}\quad(p\ge1).
\]
For a fixed \(c_*>0\),
\(q_j(x)\ge cL^{2j}\) if \(0\le x\le c_*L^{-2j}\).
The initial polynomial is bounded below by a positive constant on
all of \([0,4]\). After a fixed spectral rescaling the same assertions
hold on any fixed bounded interval.
\end{lemma}
\begin{proof}
Choose a nonnegative, even, nonzero smooth function \(\rho\) supported
in \([-1/2,1/2]\), and put
\[
 \widehat\rho(u)=\int\rho(v)e^{iuv}\,dv,\qquad
 f(u)=|\widehat\rho(u)|^2,\qquad
 C_f=\int_0^\infty u f(u)\,du>0.
\]
The function \(f\) is even, nonnegative and rapidly decreasing,
\(f(0)>0\), and its inverse Fourier transform has support
in \([-1,1]\). For \(t>0\) set
\[
 P_t(\theta)=\sum_{n\in\mathbb Z}f(t(\theta+2\pi n)).
\]
Poisson summation shows that \(P_t\) is an even trigonometric
polynomial of degree at most \(t\). Hence it is a polynomial in
\(x=2-2\cos\theta\) of that degree, and is nonnegative. For
\(0<t<1\) it is the positive constant
\(t^{-1}\int\rho(v)^2\,dv\). For \(0<\theta\le\pi\), Tonelli's
theorem and the elementary partial-fraction identity for \(\sin\)
give
\[
 \int_0^\infty tP_t(\theta)\,dt
 = C_f\sum_{n\in\mathbb Z}(\theta+2\pi n)^{-2}
 =\frac{C_f}{4\sin^2(\theta/2)}=\frac{C_f}{x}.
\]
Define \(q_0=C_f^{-1}\int_0^L tP_t\,dt\) and
\(q_j=C_f^{-1}\int_{L^j}^{L^{j+1}}tP_t\,dt\) for \(j\ge1\).
These are polynomials of the asserted degrees and sum to \(x^{-1}\).
Rapid decrease gives, for \(t\ge1\),
\(P_t(\theta)\le C_p(1+t^2x)^{-p}\), since
\(\theta^2\asymp x\) on \([0,\pi]\). Integrating on each
fixed-ratio interval proves the upper bounds; \(q_0\) is bounded
and \(x\le4\). Finally \(f(u)\ge c>0\) for sufficiently small
\(|u|\), so \(P_t(\theta)\ge c\) if
\(x\le c_*L^{-2j}\), with \(c_*\) small enough that
\(t|\theta|\) is in that neighborhood for all
\(L^j\le t\le L^{j+1}\). Integration gives the lower bound.
The integral over \(0<t<1\) gives the global lower bound for \(q_0\).
\end{proof}

\begin{proposition}[Positive fractional pieces after Hessian absorption]
\label{complete:covariance-prop}
For any prescribed finite spatial derivative order \(k_*\), choose
\(N\) sufficiently large and then \(R\) sufficiently large. The exact
covariance
\begin{equation}\label{complete:target-cov}
 \mathscr C=\frac{A}{B(B-cA)+AK}
\end{equation}
has positive semidefinite finite-range pieces \(\mathscr C_j\),
\(j\ge0\), of range at most \(CRL^j\), with
\begin{equation}\label{complete:cov-bounds}
 \|\nabla^k\mathscr C_j\|_\infty
 \le C R^{-d-k}L^{-(d-s+k)j}
              (1+\mu L^{sj})^{-2},\qquad 0\le k\le k_*.
\end{equation}
The same bound holds for the Fourier integral with weight \(|k|^k\).
The pieces are continuous in parameters whenever \(K\) is continuous
in a fixed weaker exponential kernel norm. The normalized pieces
\(C_j=\lambda R^s\mathscr C_j\) satisfy a bubble lower bound
\(\sum_x C_j(x)^2\ge c\lambda^2\) whenever
\(\mu L^{sj}\le c_0\), for a fixed \(c_0>0\).
\end{proposition}

\begin{proof}
We first assume \(K\) has finite range \(C_PR\), and later remove
that restriction by moment matching. Its zero row sum, symmetry and
second moment give \(|\widehat K(k)|\le Cg\min(R^2|k|^2,1)\).
The box multiplier comparison makes the latter minimum comparable
to \(\ell(k)\). Let
\[
 F=cB-K,\qquad \mathscr T=B^2-AF=B(B-cA)+AK.
\]
For large \(R\), \(F\) and \(B\) are bounded above and below by
fixed positive constants. Moreover
\begin{equation}\label{complete:den-comparison}
 \widehat{\mathscr T}(k)\asymp \mu+\ell(k)^\alpha.
\end{equation}
Indeed, with \(t=\widehat D(k)\),
\[
 \widehat{B-cA}(k)
 =\mu+c\zeta\ell(k)^\alpha
       +(1-c)\zeta\bigl(Q_N(1)-Q_N(t)\bigr).
\]
Here \(c=1+O(g)\), and the last bracket is nonnegative and at most
\(C_N\ell(k)\); the Hessian term is another \(O(g\ell(k))\)
perturbation. Since \(0<\alpha<1\), these perturbations are bounded
by a small multiple of \(\ell(k)^\alpha\).

The crucial point is to represent \(A\) on an internal Hilbert space
without introducing spatially nonlocal matrix entries. Put
\[
 w(u)=\frac{\sin(\pi\alpha)}\pi
                 u^{N-\alpha-1}(1-u)^\alpha,
 \qquad \mathcal H_0=L^2((0,1),du)\oplus\mathbb R.
\]
On \(\ell^2(\mathbb Z^d;\mathcal H_0)\), let
\(\mathcal D=\operatorname{diag}(I-uD,I)\). Define the finite-range
map \(\mathcal L:\ell^2(\mathbb Z^d)\to
\ell^2(\mathbb Z^d;\mathcal H_0)\) by
\[
 (\mathcal L f)(u)=\sqrt{\zeta w(u)}D^{N/2}f,
 \qquad (\mathcal L f)_0=\sqrt{\zeta\varepsilon_R}\,f.
\]
Then \(A=\mathcal L^*\mathcal D^{-1}\mathcal L\) as quadratic
forms. This is exactly \eqref{frac:tail-resolvent}. The inverse may
be unbounded on the internal space, but its displayed compression
is bounded, since \(w(u)/(1-u)\) is integrable.

Consider the bounded, finite-range, self-adjoint operator
\begin{equation}\label{complete:block-operator}
 \mathcal Q=
 \begin{pmatrix}
   \mathcal D&-\mathcal L F\\
   -F\mathcal L^*&FB^2
 \end{pmatrix}
\end{equation}
on \(\ell^2(\mathbb Z^d;\mathcal H_0\oplus\mathbb R)\).
Its range is a fixed multiple of \(R\), its norm is uniformly
bounded, and it is positive. To see positivity, complete the square
in the first component. The Schur complement is
\(FB^2-F A F=F\mathscr T\ge0\). This form calculation is valid
by approximation with \(\mathcal D+\delta I\) and passage to
\(\delta\downarrow0\), or directly because
\(\mathcal D^{-1/2}\mathcal L\) is bounded. The compressed inverse is
\begin{equation}\label{complete:block-compression}
 V_0=\mathcal L^*\bigl[\mathcal Q^{-1}\bigr]_{00}\mathcal L
       =\frac{AB^2}{B^2-AF},\qquad
 \mathscr C=B^{-2}V_0.
\end{equation}
All inverses in this assertion can equivalently be taken as positive
quadratic-form limits of resolvents. In particular no finite internal
dimension or uniform gap at zero frequency is being assumed.

We record quantitative estimates for those resolvents. Set
\(A_r=\mathcal L^*(\mathcal D+r)^{-1}\mathcal L\), \(r>0\).
At each frequency let \(h_r=B^2+r/F\). The Schur formula gives
\begin{equation}\label{complete:resolvent-compression}
 V_r=\mathcal L^*\bigl[(\mathcal Q+r)^{-1}\bigr]_{00}\mathcal L
       =\frac{A_rh_r}{h_r-F A_r},\qquad
 -V_r'=\frac{h_r^2(-A_r')+A_r^2}{(h_r-F A_r)^2}.
\end{equation}
The second identity follows by differentiation and \(Fh_r'=1\).
It is important to use this derivative: \(V_r\) alone would give
only one postmass power.

Write \(\rho=r+\ell(k)\) and restrict first to small \(r,\ell(k)\).
The integral defining \(A_r\) gives, uniformly,
\begin{align}
 A_r&\asymp1, &
 -A_r'&\le C\rho^{\alpha-1},\nonumber\\
 A-A_r&\asymp r\rho^{\alpha-1}, &
 h_r-F A_r&\asymp\mu+\rho^\alpha.
 \label{complete:resolvent-estimates}
\end{align}
For example, near \(u=1\), \(w(u)\asymp(1-u)^\alpha\), while
\(1+r-u\widehat D(k)\asymp r+\ell(k)+1-u\).
Thus the derivative is bounded by a constant times
\(\int_0^1v^\alpha(v+\rho)^{-2}\,dv\le C\rho^{\alpha-1}\).
The difference \(A-A_r\) is the positive integral with integrand
\(r v^\alpha/[(v+\ell)(v+\ell+r)]\); it is comparable to
\(r\rho^{\alpha-1}\). Combining this with
\eqref{complete:den-comparison} proves the last estimate, since
\(\ell^\alpha+r\rho^{\alpha-1}\asymp\rho^\alpha\).
The bounded white component does not affect these estimates.
Consequently
\begin{equation}\label{complete:V-derivative}
 -V_r'\le C\frac{(r+\ell)^{\alpha-1}}
                         {[\mu+(r+\ell)^\alpha]^2},
 \qquad V_r\asymp\frac1{\mu+(r+\ell)^\alpha}
\end{equation}
in this low-frequency region.

Away from \(\ell=0\), the denominators are uniformly bounded below.
Every derivative of \(A_r\) in \(r\), for \(0<r\le1\), is bounded
by \(C_m(|\widehat D(k)|^N+\varepsilon_R)\). The same is true
of \(-V_r'\). Indeed \(A_r^2\le C(-A_r')\), as follows directly
from the integral and Cauchy--Schwarz, or from the uniform upper
bound on \(\mathcal D+r\) and on \(A_r\).

At each frequency the operator in \eqref{complete:block-operator}
also satisfies
\begin{equation}\label{complete:matrix-gap}
 \mathcal Q(k)\ge c\ell(k)I.
\end{equation}
For small positive \(\ell\), completion of the square gives a lower
bound
\(c\ell\|z-Jv\|^2+c(\mu+\ell^\alpha)|v|^2\), where
\(J=\mathcal D^{-1}\mathcal L F\).
The integral estimate above gives \(\|J\|^2\le C\ell^{\alpha-1}\).
Hence \(\ell(\|z\|^2+|v|^2)\) is bounded by a constant times that
completed square. Away from zero the same argument has uniformly
bounded coefficients. This proves \eqref{complete:matrix-gap}.

Apply Lemma~\ref{complete:scalar-polynomials} to the bounded positive
operator \(\mathcal Q\).
After a fixed spectral rescaling its polynomials \(q_j\) obey
\[
 \deg q_j\le CL^j,\quad q_j(t)\ge0,\quad
 \sum_{j\ge0}q_j(t)=t^{-1},\quad
 q_j(t)\le C_p r^{-1}(1+t/r)^{-p},\qquad r=L^{-2j}.
\]
They also satisfy \(q_j(t)\ge c r^{-1}\) for \(0\le t\le c_*r\).
These are scalar statements; the spectral theorem proves them on
any Hilbert space, and polynomial degree proves the finite-range
property, independently of the internal dimension. The upper bound
and \eqref{complete:matrix-gap} imply
\[
 \mathcal L^*[q_j(\mathcal Q)]_{00}\mathcal L
 \le C_p r(1+\ell/r)^{-(p-2)}(-V_r').
\]
All terms are positive multipliers. Using
\eqref{complete:V-derivative}, their low-frequency bound is
\begin{equation}\label{complete:slice-fourier}
 C_p r^{-\alpha}(1+\mu r^{-\alpha})^{-2}
              (1+\ell/r)^{-(p-1-\alpha)}.
\end{equation}
Their high-frequency bound is
\(C_pr^{p-1}(|\widehat D|^N+\varepsilon_R)\).

The box multiplier comparison and a change of variables yield
\[
 \int_{|k|\le C/R}|k|^m(1+\ell(k)/r)^{-q}\,dk
   \le C R^{-d-m}r^{(d+m)/2}
\]
when \(q>(d+m)/2\). Moreover, for \(N>m+1\),
\[
 \int |k|^m(|\widehat D(k)|^N+e^{-R})\,dk
       \le C_{m,N}R^{-d-m}.
\]
For the latter use the product Dirichlet bound
\(|\widehat D(k)|\le C\prod_i\min(1,(R|k_i|)^{-1})+CR^{-d}\).
Choosing \(p\) so that \(p-1>3\alpha+k_*/2\), and enlarging it further for the low-frequency integral, proves
\eqref{complete:cov-bounds} and its weighted Fourier version for
these compressed pieces of \(V_0\).

Multiplication by \(B^{-2}\) preserves a positive finite-range
decomposition with these bounds. Indeed the same scalar construction
for the finite-range, uniformly gapped operator \(B\) gives positive
pieces for \(B^{-1}\) with arbitrarily rapid decay in their range
index. Convolve two such decompositions to obtain one for \(B^{-2}\),
and group products with the pieces of \(V_0\) by the maximum range
index. In the terms whose maximum comes from \(V_0\), sum the bounded
positive multiplier of the other factor. In the remaining terms use
the arbitrarily rapid decay and the finite weighted Fourier integral
of the full compressed covariance. These are precisely the two
estimates used in the product construction of
Proposition~\ref{prop:fc_newcovbound}.

For completeness the low-frequency lower bound follows from the
compressed spectral measure, not from an operator gap of fractional
order. Write \(V_r(k)=\int_0^S(t+r)^{-1}\,d\nu_k(t)\).
For \(\mu\le c_0r^\alpha\) and \(\ell(k)\le c_1r\),
\eqref{complete:V-derivative} implies
\(V_{\delta r}\ge a\delta^{-\alpha}r^{-\alpha}\) and
\(V_{c_*r}\le A c_*^{-\alpha}r^{-\alpha}\), if \(c_0,c_1\)
are taken small relative to \(\delta\). On \(t>c_*r\),
\((t+\delta r)^{-1}\le2(t+c_*r)^{-1}\). Choose fixed
\(\delta>0\) so small that the first lower bound exceeds four times
the second upper bound. It follows that
\(\nu_k([0,c_*r])\ge c r^{1-\alpha}\).
The scalar lower bound on \(q_j\) therefore gives a compressed
slice at least \(c r^{-\alpha}\) on that ball. The initial scale
uses the constant positive polynomial in the initial time interval.
The zero-scale positive piece of \(B^{-2}\) is bounded below, so
the same bound holds after the product construction.
After multiplying by \(\lambda R^s\), Plancherel on a ball of
volume \(cR^{-d}L^{-dj}\) proves the stated bubble lower bound,
because \(d=2s\).

Finally remove the finite-range assumption on \(K\). Match its
moments through an arbitrarily high fixed order \(P\) by a symmetric
finite-range kernel \(K_f\) of range \(C_PR\), zero row sum, and
weighted norm \(O_P(g)\), exactly as in
Proposition~\ref{prop:fc_newcovbound}. Let
\(T_f=B(B-cA)+AK_f\) and \(T=B(B-cA)+AK\). The positive ratio
\[
 M(k)=\widehat T_f(k)/\widehat T(k)
       =1-\widehat A(k)(\widehat K(k)-\widehat K_f(k))/\widehat T(k)
\]
is \(1+O(g)\). In rescaled coordinates \(\xi=Rk\), its derivatives
through any prescribed order \(m\) are integrable, with
\[
 \int|\partial_k^m(M-1)|\,dk\le C_{P,m}gR^{m-d},
\]
provided \(P\) is sufficiently large. Near zero, moment matching
makes the derivatives of \(K-K_f\) bounded by
\(Cg|\xi|^{P+1-m}\); differentiation of the fractional denominator
costs at most \(|\xi|^{-s-m}\). Derivatives of the fractional
numerator obey the same symbol bounds, so the resulting powers are
integrable for large \(P\), uniformly in \(\mu\ge0\). Away from
zero the denominator is bounded below and the filter smoothing,
with \(N\) chosen large after \(P\), gives the displayed integral
bound.
Truncate the Fourier series at range \(RL^n\), subtract decreasing
error constants, and take consecutive differences as in the proof of
Proposition~\ref{prop:fc_newcovbound}. This decomposes \(M\) into
positive trigonometric polynomials \(M_n\), with \(M_0\ge1/2\)
and arbitrarily rapid fixed decay in \(n\). Multiply them by the
already constructed covariance pieces and group by maximum index.
The weighted Fourier estimates give \eqref{complete:cov-bounds};
choose the decay exponent larger than \(d+s+k_*\) to preserve
its postmass factor. The factor \(M_0\ge1/2\) preserves the lower
bound. All operations are continuous in the indicated parameter
norms. This proves the proposition.
\end{proof}

\begin{corollary}[Critical Gaussian asymptotic and bubble envelopes]
\label{complete:covariance-consequences}
At \(\mu=0\), with \(0<\alpha<1\), the covariance in
\eqref{complete:target-cov} obeys
\[
 (\lambda R^s\mathscr C)(x)\sim
 \frac{\lambda R^s A(0)}{B(0)c\zeta a_{R,d}^\alpha}
 \frac1{(2\pi)^{d/2}}|x|^{-s},
\]
where \(A(0),B(0)\) in this display denote the zero-frequency
multipliers (row sums). Its normalized positive pieces satisfy
\[
 0\le b_j=a_0\sum_x(w_j(x)^2-w_{j-1}(x)^2)
       \le C\lambda^2(1+\mu L^{sj})^{-2},
 \qquad w_j=\sum_{i\le j}\lambda R^s\mathscr C_i,
\]
and \(b_j\ge c\lambda^2\) when \(\mu L^{sj}\le c_0\).
\end{corollary}
\begin{proof}
The analytic kernels \(B,K\) have expansions with remainder
\(O_R(|k|^2)\), while the only leading nonsmooth contribution to
\(B-cA\) is \(c\zeta a_{R,d}^\alpha|k|^s\).
Consequently \(\widehat{\mathscr T}(k)
=B(0)c\zeta a_{R,d}^\alpha|k|^s+O_R(|k|^2)\).
Division shows that the covariance multiplier is its displayed
positive multiple of \(|k|^{-s}\), plus a symbol of strictly higher
order: its derivatives of order \(m\) are bounded by
\(C_{R,m}|k|^{-s+\delta-m}\), where
\(\delta=\min(s,2-s)>0\). A smooth dyadic partition and integration
by parts give a Fourier remainder \(O_R(|x|^{-s-\delta})\).
The Riesz kernel coefficient is \((2\pi)^{-d/2}\), because
\(d=2s\). This proves the asymptotic.
For the bubble bound, positivity gives nonnegative mixed sums by
Plancherel. The range and pointwise estimates give
\(\sum_{i\le j}\|C_i\|_1\le C\lambda R^sL^{sj}\), and multiplying
this by the pointwise bound on \(C_j\) proves the upper envelope.
The square of the last piece and the lower bound in
Proposition~\ref{complete:covariance-prop} give the lower bound.
\end{proof}

\section{Convolution back to spin variables}
\label{frac:inverse-kernels}

For \(d<4\), consider the even-tail split with a strictly
positive diagonal addition. Its multiplier \(\widehat A\) is
smooth away from zero and has a positive constant term. Near
zero its difference from that constant, and the corresponding
difference for \(1/\widehat A\), satisfy symbol derivative
bounds \(C_{R,m}|k|^{s-m}\) of every fixed order \(m\).
This follows from \eqref{frac:symbol}, the analytic polynomial
part of \(P_N\), and the positive constant term in the inverse.
A smooth dyadic partition near zero, followed by integration
by parts on each annulus, gives
\[
 |A^{-1}(x)|+|(BA^{-1})(x)|
     \le C_R(1+|x|)^{-d-s}
\]
when \(B\) is finite-range. For instance, on an annulus of radius
\(2^{-j}\), the nonsmooth part of a derivative of order \(m\)
is \(O_R(2^{-j(s-m)})\). Its Fourier coefficient is bounded by
\(C_{R,M}2^{-j(d+s)}(1+2^{-j}|x|)^{-M}\); summing in \(j\)
proves the assertion for \(M>d+s\). The smooth part has
arbitrarily rapid decay. This replaces the exponential kernel
bound used in the original spin reconstruction.



\section{Power counting and the enlarged activity ball}
\label{frac:scaling}
Write \(s=d/2\), so that the fractional precision has order \(s\)
and the Gaussian field has scaling dimension
\[
 [\phi]=\frac{d-s}{2}=\frac d4=\frac s2.
\]
For a block of side \(b_j=bL^j\), its fluctuation radius in
canonical field units is therefore proportional to \(b_j^{-d/4}\).
The scaling degree of a local monomial with \(p\) fields and
\(r\) spatial derivatives is \(p d/4+r\). In particular,
\[
 \begin{array}{c|c|c}
 \text{monomial}&\text{dimension minus }d&\text{linear RG factor}\\ \hline
 \phi^2&-s&L^s\\
 \phi^4&0&1\\
 \phi^6&s&L^{-s}\\
 |\nabla\phi|^2&2-s&L^{s-2}
 \end{array}
\]
The quartic is marginal in every dimension considered here. The local
sextic remains irrelevant, including in dimension one. For \(d<4\),
the ordinary local kinetic term is irrelevant; the backward tuning
of a marginal local stiffness in the four-dimensional argument must
not be copied into these dimensions. The nonlocal fractional kinetic
term is part of the Gaussian precision and is not a local monomial
produced by polynomial localization.

For the proposed four-sector norm, the fluctuation source radius
at coalescence scale \(k\) is
\[
 \ell_{\sigma,j}=G_j b_{j\wedge k}^{s/2}2^{(j-k)_+},
 \qquad \ell_{\sigma\sigma,j}=\ell_{\sigma,j}^2.
\]
The large-field radii are
\[
 h_j=k_0G_j^{-1/4}b_j^{-s/2},\qquad
 h_{\sigma,j}=G_j^{1/4}b_{j\wedge k}^{s/2}2^{(j-k)_+},\qquad
 h_{\sigma\sigma,j}=h_{\sigma,j}^2.
\]
Thus \(h_{\sigma,j}h_j=k_0(2L^{-s/2})^{(j-k)_+}\le k_0\), and
\(\ell_{\sigma,j}/h_{\sigma,j}=G_j^{3/4}\) at every scale.
The full regulated estimates for these weights are proved in
Section~\ref{completion:regulated-localization}.

\subsection{Initialization scales}
Put \(g=R^{-d}\), choose \(b\asymp R\), and let
\(A_{\rm sc}=g^{-a_*}\), \(\lambda=A_{\rm sc}^{-2}\).
In the Hessian-subtracted variable \(u\), make the field change
\[
 \psi=\frac{R^{d/4}}{A_{\rm sc}}u.
\]
The coefficient of \(\psi^4/4\) in the on-site quartic is exactly
\[
 G_0=\frac{A_{\rm sc}^4g}{3}.
\]
The physical fluctuation and large-field radii on an initial block are
\[
 \ell'_u\asymp A_{\rm sc}R^{-d/2},\qquad
 h_u\asymp R^{-d/4}.
\]
Thus \(\ell'_u/h_u\to0\) whenever \(a_*<1/4\). A local sextic
with a bounded microscopic coefficient has block sizes
\[
 b^d(\ell'_u)^6=O(A_{\rm sc}^6g^2),\qquad
 b^d h_u^6=O(g^{1/2}).
\]
A quartic coefficient defect of order \(g\) has fluctuation size
\(O(A_{\rm sc}^4g^2)\), which is smaller than the displayed
sextic bound. These are power-counting identities. The all-field connected estimate is established in Section~\ref{completion:flow-closure}.

The initialization in Section~\ref{completion:flow-closure} gives the weighted norm bound
\[
 \|K_0\|\le U=C\bigl(A_{\rm sc}^6g^2+
                    G_0^{9/4}g^{1/2}\bigr),
\]
then the amplitude rescaling used in dimension four remains admissible
without any dimension-dependent change to its exponents. Indeed, for
\(1/12<a_*<1/10\),
\begin{align}
 \frac{U}{G_0^{9/4}}
   &\le C\bigl(g^{3a_*-1/4}+g^{1/2}\bigr)\longrightarrow0,
      \label{frac:eq:seed-small}\\
 \frac{G_0}{\lambda^3}
   &=\frac13g^{1-10a_*}\longrightarrow0,\qquad
 \frac{U}{\lambda G_0}
   \le C\bigl(g^{1-4a_*}+g^{7/4-7a_*}\bigr)
        \longrightarrow0.
      \label{frac:eq:seed-ratios}
\end{align}
These statements establish the consistency of the smallness regime;
the regulated map theorem is proved in Section~\ref{completion:regulated-localization}.

\subsection{The analytic input and its Cauchy consequence}
The following formulation specifies the input needed to transfer a
contraction to the larger initial activity. Norms use a fixed positive
reference coupling \(G\), including at complex potential inputs.
The activity space must include all four source sectors, with the
specified support, parity, symmetry and pre-coalescence conditions.
Assume the exact maps \(\Psi=K_+\) and
\(R=V_+-\Phi^{\rm pt}(V)\) are holomorphic on the joint balls
\[
 \|V-V_0\|<cG,\qquad \|K\|<cG^{9/4},
\]
and satisfy, throughout these balls,
\begin{equation}\label{frac:eq:analytic-input}
 \|\Psi(V,K)\|\le CG^{9/4},\quad
 \|\Psi(V,0)\|+\|R(V,0)\|\le C\chi^{3/2}G^3,
 \quad \|D_KR\|\le C.
\end{equation}
Assume separately that at real potential inputs
\(\|D_K\Psi(V,0)\|\le\kappa_0<1\), with all constants uniform
in the spread radius, scale, volume and mass cutoff.

\begin{lemma}[Cauchy extension of a uniform contraction]
\label{frac:lem:cauchy}
Under these analytic inputs, on fixed smaller balls and
\(\|K\|\le\epsilon G^{9/4}\),
\begin{align*}
 \|D_K\Psi(V,K)\|&\le\kappa_0+C\epsilon
                         &&\text{for real }V,\\
 \|\Psi(V,K)\|+\|R(V,K)\|
    &\le C\bigl(\chi^{3/2}G^3+\|K\|\bigr),\\
 \|D_V\Psi(V,K)\|+\|D_VR(V,K)\|
    &\le CG^{-1}\bigl(\chi^{3/2}G^3+\|K\|\bigr).
\end{align*}
\end{lemma}
\begin{proof}
Cauchy estimates in activity directions give
\(\|D_K\Psi\|\le C\) and
\(\|D_K^2\Psi\|\le CG^{-9/4}\). Integrating the second derivative
along the segment from zero to \(K\) proves the first assertion.
Integrating the first derivatives of \(\Psi,R\) gives the second.
Cauchy's estimate on a potential disc of radius proportional to \(G\)
gives the third. The constants depend only on the fixed radius losses.
Choose \(\epsilon\) before making \(G\) small, so that
\(\kappa_0+C\epsilon<1\). An estimate only at \(K=0\) would not
supply \eqref{frac:eq:analytic-input} or this conclusion.
\end{proof}

For a fractional covariance the one-step mass multiplier is \(L^s\).
The mass Green function is consequently
\(L^{-s(i-j+1)}\), replacing \(L^{-2(i-j+1)}\). A local kinetic
coordinate, if retained when \(d<4\), obeys a stable forward equation
with linear part \(z_{j+1}=L^{s-2}z_j\). With \(L\) large, this
contraction can be made smaller than any fixed mass-damping ratio.
Only the relevant mass needs a physical tuning in this case. In
\(d=4\) the local kinetic coordinate remains marginal and the original
backward tuning is required. Theorem~\ref{completion:map-theorem} supplies the analytic estimates; the covariance moments and physical-coordinate Wick formulas supply these normal-form equations and their remainders.

\section{The observable mechanism}
\label{frac:sec:observables}
Here is the precise last step once a fractional two-source construction
is available. Suppose the massless Gaussian covariance has positive
finite-range slices \(\Gamma_{j+1}\), with range at most \(Cb_j\),
\[
 \|\Gamma_{j+1}\|_\infty\le C\lambda b_j^{-s},
 \qquad C(x)=\sum_{j\ge0}\Gamma_{j+1}(x)
                   \sim c_0|x|^{-s},\quad c_0>0.
\]
Suppose its interacting two-source sum satisfies
\begin{equation}\label{frac:eq:source-input}
 q(a,c)=\sum_{j\ge0}v_{a,j}v_{c,j}\Gamma_{j+1}(a-c)+E(a,c),
 \qquad |E(a,c)|\le b_k^{-s}\eta_k,
 \quad \eta_k\longrightarrow0,
\end{equation}
where \(b_k\asymp|a-c|\). Assume
\(v_{a,j}\to v_{a,\infty}\) uniformly in the marked site, with
\(0<c\le v_{a,j}\le C\). For a small fixed pair of scales near
coalescence the same convergence is understood as their indices tend
to infinity. These observable estimates are established by the graded source argument in Section~\ref{completion:flow-closure}.

\begin{proposition}[Observable conclusion under the source estimates]
\label{frac:prop:observable}
If \(q\) is translation invariant, then there is \(v_\infty>0\)
independent of the marked site such that
\[
 q(a,c)\sim v_\infty^2 c_0|a-c|^{-s}.
\]
\end{proposition}
\begin{proof}
The range condition makes the covariance sum start at \(k-O(1)\).
The sum of the absolute slice bounds from that scale onward is at most
\(C\lambda b_k^{-s}\). Uniform convergence of \(v_{a,j}\) therefore
allows their replacement by their limiting values at cost
\(o(b_k^{-s})\). The source error has the same order. Thus
\[
 q(a,c)=v_{a,\infty}v_{c,\infty}C(a-c)+o(|a-c|^{-s}),
\]
uniformly in the two sites.
Translation invariance, comparing \((a,a+x)\) with \((0,x)\), gives
\(v_{a+x,\infty}/v_{x,\infty}\to
r(a)=v_{0,\infty}/v_{a,\infty}\). For every fixed positive integer
\(p\), multiplying this ratio at the shifts \(x,x+a,\ldots,x+(p-1)a\)
gives \(r(pa)=r(a)^p\). All these ratios are bounded above and
away from zero independently of \(p\); hence \(r(a)=1\).
This proves that the limiting coefficient is independent of the site
and proves the assertion.
\end{proof}



\section{Dimension-dependent estimates behind a fractional map theorem}
\label{completion:map-support}
This section records consequences of the fractional covariance estimates.
It separates these consequences from the nonlinear one-step map theorem.
Put \(s=d/2\), \(d\in\{1,2,3\}\), and \(b_j=bL^j\).
Assume translation-invariant, hypercubic-symmetric positive finite-range
pieces with
\[
 \operatorname{range}(\Gamma_{j+1})\le Cb_j,\qquad
 \|\nabla^r\Gamma_{j+1}\|_\infty
 \le C_r\lambda b_j^{-s-r}(1+\mathfrak m b_j^s)^{-2}
\]
for the finite collection of derivative orders required below. Constants
may depend on the fixed scale factor, but not on the mass cutoff.
Here \(\mathfrak m=\mu/R^s\) if \(\mu\) is the dimensionless
auxiliary mass used with \(b_j\asymp RL^j\); thus the damping is comparable to
\((1+\mu L^{sj})^{-2}\). Write
\(w_j=\sum_{i<j}\Gamma_{i+1}\), and choose \(J\) so that
\(\mathfrak m b_J^s\) is comparable to one.

\begin{lemma}[Moment estimates and cancellation of the marginal logarithm]
\label{completion:moments}
At every scale,
\begin{align*}
 |w_j(x)|&\le C\lambda(b+|x|)^{-s},&
 \|w_j\|_1&\le C\lambda b_j^s,\\
 \sum_x|x|^2|w_j(x)|&\le C\lambda b_j^{s+2},&
 \sum_x|x|^2w_j(x)^2&\le C\lambda^2b_j^2.
\end{align*}
Put \(\Gamma=\Gamma_{j+1}\) and
\(a_j=(1+\mathfrak m b_j^s)^{-2}\). Then
\begin{align*}
 \left|\sum_x((w_j+\Gamma)^2-w_j^2)(x)\right|
    &\le C\lambda^2a_j,\\
 \left|\sum_x\big((w_j+\Gamma)^3-w_j^3
                   -3\Gamma(0)w_j^2\big)(x)\right|
    &\le C\lambda^3a_j b_j^{-s},\\
 \sum_x|x|^2\left|(w_j+\Gamma)^3(x)-w_j^3(x)\right|
    &\le C\lambda^3a_j b_j^{2-s}.
\end{align*}
After the mass scale the normalized cumulative moments satisfy
\[
 b_j^{-s}\|w_j\|_1\le C\lambda L^{-s(j-J)},\qquad
 b_j^{-s-2}\sum_x|x|^2|w_j(x)|
       \le C\lambda L^{-2s(j-J)}\qquad(j\ge J).
\]
\end{lemma}
\begin{proof}
The range bound makes only pieces with \(b_i\ge c|x|\) contribute
at \(x\); their pointwise bounds form a geometric sum. Summing its
result over the support of \(w_j\) gives the two absolute moments.
Since \(d=2s\), summing \(|x|^2(b+|x|)^{-2s}\) gives
\(O(b_j^2)\), proving the fourth estimate.
The bubble bound follows by expanding the difference of squares and
using \(\|\Gamma\|_\infty\|w_j\|_1\) and
\(\|\Gamma\|_\infty\|\Gamma\|_1\).
Evenness and the second-difference bound give
\[
 |\Gamma(x)-\Gamma(0)|
       \le C\lambda a_j b_j^{-s-2}|x|^2.
\]
Thus the potentially logarithmic term in the cubic increment is
\(3\sum w_j(x)^2(\Gamma(x)-\Gamma(0))\), bounded by the fourth
moment estimate. The remaining terms are
\(3\sum w_j\Gamma^2+\sum\Gamma^3\), bounded by
\(C\lambda^3a_j^2b_j^{-s}\) and
\(C\lambda^3a_j^3b_j^{-s}\), respectively.
Inserting \(|x|^2\) in the uncorrected cubic expansion instead
uses \(\sum|x|^2w_j^2\), \(\sum|x|^2|w_j|\), and the
range of \(\Gamma\), proving the kinetic-moment estimate.
Finally, the postmass row-sum slices are bounded by
\(C\lambda\mathfrak m^{-2}b_i^{-s}\), whose sum is
\(O(\lambda b_J^s)\). Their second absolute moments are bounded by
\(C\lambda\mathfrak m^{-2}b_i^{2-s}\). Since \(s<2\), the latter sum
is dominated by its final term. Division by the displayed normalizing
powers gives exactly the two postmass estimates. In particular the
second moment does not decay as \(L^{-(s+2)(j-J)}\) under only the
natural mass power two.
\end{proof}

The normalized Wick mass increment is therefore \(O(\lambda a_j)\),
the bubble is \(O(\lambda^2a_j)\), the corrected cubic mass coefficient
is \(O(\lambda^3a_j)\), and the local kinetic coefficient has the same
size after multiplication by \(b_j^{s-2}\). For large fixed \(L\),
the postmass normalized moments are smaller than any prescribed fixed
power of \(\chi_j=2^{-(j-J)_+}\) up to the exponent permitted by
\(L^s\); in particular they are bounded by \(C\lambda\chi_j^{3/2}\).
No arbitrary covariance mass-decay power is needed for this conclusion.

\begin{lemma}[The decaying Gaussian regulator]
\label{completion:H}
For a block \(B\) of \(n=b_j^d\) sites, put
\[
 \ell_j=\ell_0b_j^{-s/2},\qquad
 h_j=k_0G^{-1/4}b_j^{-s/2},\qquad
 H_B(\phi)=\exp[-\kappa\|\phi_B\|_2^2/(nh_j^2)].
\]
If \(\zeta\) has covariance \(\Gamma_{j+1}\), then
\(\mathbb EH_B(\phi+\zeta)\le H_B(\phi)^{1/2}\) for sufficiently
small \(\lambda G^{1/2}\), uniformly in the mass scale.
\end{lemma}
\begin{proof}
The covariance restricted to \(B\) has operator norm at most
\(C\lambda a_j b_j^{d-s}=C\lambda a_jb_j^s\).
Writing \(q=\kappa/(nh_j^2)\), one obtains
\(2q\|\Gamma_B\|_{\rm op}\le C\kappa k_0^{-2}\lambda G^{1/2}\).
For this quantity at most one, Gaussian completion of the square gives
\[
 \mathbb EH_B(\phi+\zeta)
 =\det(I+2q\Gamma_B)^{-1/2}
 \exp[-q\langle\phi,(I+2q\Gamma_B)^{-1}\phi\rangle]
 \le H_B(\phi)^{1/2}.
\]
The crucial equality is \(d=2s\), not \(d=4\).
\end{proof}

\paragraph{The precise algebraic irrelevance gap.}
At zero field consider a scalar even local functional, with spatial
support of diameter \(O(b_j)\), and its coefficient pairing with
coarse-scale test tensors. Let \(\Delta=s/2\). Coarse test tensors of
field degree \(p\) have size \(\ell_{j+1}^p\), and their spatial
derivatives of order \(r\) have size
\(\ell_{j+1}^p b_{j+1}^{-r}\). Localize the constant, mass and quartic
terms. Use a canonical even quadratic Fourier symbol, or centered
lattice derivatives modulo total differences, to eliminate the
one-derivative quadratic density. Then the remaining pairings have
factors
\[
 L^{-6\Delta}=L^{-3s}\quad(p\ge6),\qquad
 L^{-2\Delta-2}=L^{-(s+2)}\quad(p=2),\qquad
 L^{-4\Delta-1}=L^{-(d+1)}\quad(p=4).
\]
For \(p=2\) this is the second-order spatial Taylor remainder; for
\(p=4\) a first-order remainder suffices. The higher-field terms follow
by the ratio \(\ell_{j+1}/\ell_j=L^{-\Delta}\). Reblocking adds at most
\(CL^d\) fine-block positions. Hence the small-field linear
irrelevance factor is
\begin{equation}\label{completion:gap}
 CL^{-\delta},\qquad
 \delta=\min(s,2-s,1)=
 \begin{cases}1/2,&d=1,3,\\1,&d=2.\end{cases}
\end{equation}
The same calculation for a single source, anchored at its marked site,
has no volume factor. After removing the linear source-field term,
the first remaining field term is cubic (gain \(L^{-s}\)) or contains
one spatial derivative (gain \(L^{-1}\)). In the mixed sector after
coalescence the constant is extracted and the next even field term
gains \(L^{-s}\); the growth of the mixed source radius costs only the
fixed factor four. A sufficiently large \(L\) absorbs this cost.

This is the small-field localization calculation. A proof of a full
activity-norm contraction also needs the regulator estimates and the
nonlinear reblocking/circle-product theorem; it is not implied by
\eqref{completion:gap} alone.

\paragraph{Dimension-neutral sizes in the nonlinear step.}
The identities
\[
 |B|\ell_j^4=\ell_0^4,\qquad
 |B|h_j^4=k_0^4G^{-1},\qquad
 \ell_j/h_j=(\ell_0/k_0)G^{1/4}
\]
explain the same coupling powers as in four dimensions. Centered
quartic and linear-source fluctuations have size
\(O(\chi_j^{1/2}G)\) at \(\ell\) and
\(O(\chi_j^{1/2}G^{1/4})\) at \(h\), provided the indicated Gaussian
Sobolev moments are bounded. Their third-order integral Taylor
remainders have sizes
\(O(\chi_j^{3/2}G^3)\) and
\(O(\chi_j^{3/2}G^{3/4})\), respectively. Multiplication of the latter
by the large-field weight \(G^{9/4}\) again gives \(G^3\).
A localized activity direction of radius \(\tau=cG^{9/4}\) has
sizes \(O(\tau)\) at \(\ell\) and \(O(c)\) at \(h\); taking \(c\)
small relative to \(k_0^4\) preserves quartic stability on the complex
potential disc. The next section proves the regulated localization and circle-product bounds needed in addition to these scaling identities.

\paragraph{Stop both source radii at coalescence.}
For a uniform four-sector mixed-norm proof, use the explicit radii
\[
 h_{\sigma,j}=G_j^{1/4}b_{j\wedge k}^{s/2}2^{(j-k)_+},\qquad
 h_{\sigma\sigma,j}=h_{\sigma,j}^2
       =G_j^{1/2}b_{j\wedge k}^{s}4^{(j-k)_+}.
\]
Thus
\[
 h_{\sigma,j}h_j=k_0(2L^{-s/2})^{(j-k)_+}\le k_0
\]
when \(L^{s/2}\ge2\); equality is not required in the source estimates.
Both small- and large-field source weights are stopped at the same
coalescence scale. Their ratios are uniform in the scale:
\[
 \ell_{\sigma,j}/h_{\sigma,j}=G_j^{3/4},\qquad
 \ell_{\sigma\sigma,j}/h_{\sigma\sigma,j}=G_j^{3/2}.
\]
The mixed radius grows by the factor four after coalescence, so a
mixed two-field remainder has the contraction factor \(4L^{-s}\),
up to adjacent coupling ratios. The mixed covariance term obeys
\[
 h_{\sigma\sigma,j}|\Gamma_{j+1}(a,c)|
       \le C\lambda\chi_jG_j^{1/2},
\]
since its extra postcoalescence factor is
\((4L^{-s})^{j-k}\le1\). If the single large-field source radius were
instead kept proportional to \(b_j^{s/2}\) after coalescence, converting
its low jets from the stopped small-field radius would introduce an
unbounded scale ratio. Stopping both radii avoids this loss.
The extended source domain is specified by these weighted coefficient
norms; nilpotence makes a fixed bound on unweighted source coefficients
unnecessary for the complex map. Bounds on actual physical source
coefficients follow separately from the one-source projection with an
unused distant second mark.

\begin{lemma}[The zero-field/large-field mixed norm]
\label{completion:Y-W}
Let \(E\) be the finite-volume test-field space and take field jets
through a fixed order \(p\ge9\). Assume the test-tensor norms permit
insertion of a field direction \(v\) at cost at most \(\|v\|_E\),
as for the product derivative test spaces. Write
\[
 |F|_{T_a(\phi)}=\sum_{q=0}^p\frac{a^q}{q!}
                    \|D^qF(\phi)\|_{E^{*q}},
 \qquad \ell/h=cG^{1/4},
\]
where the fixed factor \(c\) is harmless. Suppose
\[
 \mathcal G(\phi)=e^{c_0\|\phi\|_E^2/\ell^2},\qquad
 H(\phi)\le1,\qquad
 \widetilde G(t\phi)\le e^{c_1\|\phi\|_E^2/\ell^2}
 \quad(0\le t\le1),\qquad c_1<c_0.
\]
Define
\begin{align*}
 \|F\|_W&=\sup_\phi\frac{|F|_{T_\ell(\phi)}}{\mathcal G(\phi)}
       +G^{9/4}\sup_\phi
              \frac{|F|_{T_h(\phi)}}{H(\phi)\widetilde G(\phi)},\\
 \|F\|_Y&=|F|_{T_\ell(0)}
       +G^{9/4}\sup_\phi
              \frac{|F|_{T_h(\phi)}}{H(\phi)\widetilde G(\phi)}.
\end{align*}
Then \(\|F\|_Y\le\|F\|_W\le C\|F\|_Y\), with a constant independent
of the volume and \(G\). The statement applies sectorwise in the
nilpotent source algebra if the small-field source weight is at most
the corresponding large-field source weight.
\end{lemma}
\begin{proof}
The first inequality follows by evaluating the first supremum at zero.
For the reverse inequality put \(u=\|\phi\|_E/\ell\). For each
\(q<p\), Taylor-expand \(D^qF(\phi)\) from zero through order
\(p-q-1\) and use the integral remainder with \(D^pF(t\phi)\).
For \(q=p\), use its value directly. The tensor insertion assumption
and the factorial coefficients yield
\[
 |F|_{T_\ell(\phi)}
 \le C_p|F|_{T_\ell(0)}\sum_{r=0}^{p-1}u^r
   +C_p(\ell/h)^p\left(\sum_{r=0}^pu^r\right)
                       \sup_{0\le t\le1}|F|_{T_h(t\phi)}.
\]
The first term divided by \(\mathcal G(\phi)\) is bounded by a
constant times the zero-field norm. In the second term use
\(H(t\phi)\le1\) and the regulator bound, leaving a polynomial
in \(u\) times \(e^{-(c_0-c_1)u^2}\). Its supremum is finite.
Since \((\ell/h)^p\le C_pG^{9/4}\) for \(p\ge9\), this gives
exactly the second term in \(\|F\|_Y\).
For a source sector the same argument introduces only the ratio of
its small-field weight to its large-field weight, which is at most
one by hypothesis. Applying the estimate to each of the four sectors
and adding proves the last statement.
\end{proof}

For the explicit stopped source weights above the sector ratios are
\(G_j^{3/4}\) and \(G_j^{3/2}\), so they are at most one for small
\(G_j\). The finite-volume polymer version uses the sum of squared
local field seminorms in \(E\); the affine-quotient seminorm is bounded
by the full field seminorm. Choosing the growing regulator with a fixed
spare exponent gives the displayed domination uniformly on field segments.

\section{A fractional regulated localization estimate}
\label{completion:regulated-localization}
Put \(\Delta=s/2\), and fix \(p_\Phi=12\). On a block
neighborhood use the concrete seminorms
\[
 \|v\|_{\Phi_j}=\max_{|\gamma|\le p_\Phi}
                      b_j^{|\gamma|}\|\nabla^\gamma v\|_\infty,
 \qquad
 \|v\|_{\widetilde\Phi_j}=\inf_{p\ {
 affine}}\|v-p\|_{\Phi_j}.
\]
Fixed enlargements of block neighborhoods change only constants.
The field radii are \(\ell_j=\ell_0b_j^{-\Delta}\) and
\(h_j=k_0G_j^{-1/4}b_j^{-\Delta}\). Define growing and quotient
regulators by exponentiating small fixed positive multiples of the
squares of these seminorms divided by \(\ell_j^2\), and define
\(H_j\) by the block mean square divided by \(h_j^2\), with a
negative coefficient. Use strengthened regulator powers on intermediate
maps, leaving fixed positive margins between input and output powers.
All estimates below hold for these choices; replacing them by the
standard equivalent test-function regulators changes fixed constants.

\begin{lemma}[The affine complement estimate]
\label{completion:reg-comparison}
Let \(B\) be a coarse block of side \(Lb_j\), and let \(X\subset B\)
contain at most \(2^d\) fine blocks. Put \(\phi=p+r\), where \(p\)
is an affine approximation on the coarse neighborhood, and set
\[
 z=\left(|B|^{-1}\sum_Bp^2\right)^{1/2}/h_{j+1},\qquad
 q=\|r\|_{\Phi_{j+1}}/\ell_{j+1}.
\]
The approximation may be chosen within any fixed factor of the
quotient infimum. For fixed sufficiently large \(L\), sufficiently
small \(G_j\), and stable potentials with quartic coefficient
comparable to \(G_j\), the reference factor on \(B\setminus X\)
absorbs the input quotient regulator on \(X\), the inverse output
\(H_{j+1}\), and any fixed polynomial in \(z,q\). More precisely,
after Gaussian expectation their product divided by a fixed spare
power of the coarse quotient regulator is bounded by
\[
 C_L e^{-c z^2-cq^2}
\]
before the polynomial is absorbed. The same estimate holds for all
field derivatives through a fixed finite order, with a different fixed
polynomial before absorption. Nilpotent linear sources give only such
polynomials.
\end{lemma}
\begin{proof}
For an affine function, normalized suprema and normalized \(L^2\)
norms on a cube are equivalent. Removing at most \(2^d\) fine blocks
removes a fraction at most \(2^dL^{-d}\). The normalized supremum of
an affine function is bounded by its normalized \(L^2\) norm, so
\[
 |B|^{-1}\sum_{B\setminus X}p^2\ge c|B|^{-1}\sum_Bp^2
\]
for fixed large \(L\). Restriction of coarse test fields gives
\[
 \|\phi\|_{\Phi_j(X)}/h_j
 \le C L^{-\Delta}(z+G_j^{1/4}q),\qquad
 \|\phi\|_{\widetilde\Phi_j(X)}/\ell_j
 \le C L^{-\Delta}q.
\]
Indeed coarse scaled differences dominate the restricted fine ones;
affine functions disappear in the second inequality. Consequently the
input quotient regulator is at most \(e^{CL^{-s}q^2}\).

The identity \(|B|G_jh_{j+1}^4\asymp1\) permits the quartic Young
inequality with an arbitrary fixed quadratic coefficient \(A>0\):
\[
 \frac{G_j}{4}\sum_{B\setminus X}\psi_x^4
 \ge A|B|^{-1}\sum_{B\setminus X}\psi_x^2/h_{j+1}^2-C_A.
\]
The affine complement inequality and \((p+r)^2\ge p^2/2-r^2\)
then give, for \(\psi=\phi\),
\[
 V_{B\setminus X}(\phi)
       \ge cA z^2-C_A G_j^{1/2}q^2-C_A.
\]
The local mass and kinetic coefficients in their normalized domain
contribute at most \(C_LG_j^{1/2}(z^2+G_j^{1/2}q^2)\), and are
absorbed for small \(G_j\). This quadratic argument avoids any
uncontrolled positive \(G_jq^4\) term.
With \(\psi=\phi+\zeta\), one loses only
\(C_A|B|^{-1}\sum_B\zeta^2/h_{j+1}^2\). Its covariance quadratic
form has operator norm at most \(C_L\lambda G_j^{1/2}\).
The Gaussian determinant formula therefore bounds its exponential
moments uniformly for small \(\lambda G_j^{1/2}\).
The additional fluctuation quotient regulators are controlled by the
same formula applied to the discrete Sobolev quadratic forms through
order \(p_\Phi+d\). The covariance derivative assumption through twice
that order bounds their operator norm by \(C_L\lambda/\ell_0^2\)
and their trace by that constant times the number of blocks. The fixed
number of blocks here gives a fixed bound; on polymers it gives a fixed
exponential per block.

Dividing by \(H_{j+1}\) costs at most
\(e^{C\kappa(z^2+G_j^{1/2}q^2)}\). Choose \(A\) large enough to
leave a negative coefficient of \(z^2\). Next choose \(L\) so that
\(CL^{-s}q^2\) uses less than half the spare quotient exponent, and
then choose \(G_j\) small so that all \(G_j^{1/2}q^2\) terms use less
than half the remaining margin. This leaves
\(e^{-cz^2-cq^2}\). Field derivatives of the reference and nilpotent
source factors introduce fixed polynomials, absorbed in the retained
Gaussian decay. All choices precede taking the spread radius large.
\end{proof}

\begin{lemma}[Small-set contraction in the full mixed norm]
\label{completion:smallset}
Use field jets through order at least nine and the explicit mixed source
radius in Section~\ref{completion:map-support}. For the linearized
small-set reblocking map, with the local bulk and source counterterms
extracted and scalar quadratic symbols put in their even canonical
form, the mixed norm satisfies
\[
 \|\mathcal L_jF\|_{W_{j+1}}
 \le \left(CL^{-\delta}+C_L\lambda^{1/2}
                         +C_LG_j^{1/4}\right)\|F\|_{W_j},
 \qquad \delta=\min(s,2-s,1)>0.
\]
The background reference factors on the complement of a small set are
included in \(\mathcal L_j\). The constant multiplying \(L^{-\delta}\)
is independent of \(L\): the coarse quartic sum is normalized by
\(|B|G_jh_{j+1}^4\asymp1\), and the affine norm comparisons have uniform
constants. Gaussian moments may be bounded by two after the fixed
\(L\) is chosen and then \(C_L\lambda\) is made small. The remaining
nonuniform constants occur only in the displayed small-parameter terms.
\end{lemma}
\begin{proof}
It is enough to use the equivalent zero-field/large-field norm of
Lemma~\ref{completion:Y-W}. The small-field local Taylor pairings have
exactly the factors in \eqref{completion:gap}. The number of fine-block
positions in a coarse block is at most \(CL^d\); source sectors are
anchored and do not have this volume factor. Gaussian averaging does
not enlarge the support of the external field variables. For the lower field jets, its change of zero-field jets is at most
\(C_L\lambda^{1/2}\|F\|_{W_j}\), by Taylor expansion in the fluctuation
and the Gaussian moments. Top-order jets are bounded directly by their
coarse test-tensor restriction gain; no additional field derivative is
assumed for a Lipschitz estimate on those top jets. If the
counterterms are extracted after expectation, this error may instead
be incorporated in the localization itself. Critical terms regenerated
by Gaussian contractions must not be counted as geometrically irrelevant
without one of these two provisions.

For the large-field bound, treat each input source degree
\(r=0,1,2\) separately and Taylor-expand its field dependence through
order \(n_r-1\), where
\[
 n_r=9-3r.
\]
Low-order coefficient tensors are controlled by the zero-field norm.
After conversion to output large-field radii and multiplication by the
weight \(G_{j+1}^{9/4}\), their coupling factor is at most a fixed
constant times
\[
 G_j^{9/4-3r/4-p/4}\le C G_j^{1/4}
        \qquad(0\le p\le n_r-1).
\]
Here the inequality means that the factor is bounded above by
\(CG_j^{1/4}\) for small \(G_j\). All powers of the normalized affine
field and quotient field are absorbed by Lemma~\ref{completion:reg-comparison}.

For the integral remainder, each field direction costs
\(L^{-\Delta}(z+G_j^{1/4}q)\), and each restricted field test tensor
has the same \(L^{-\Delta}\) factor. The affine lemma again absorbs
the resulting polynomial. The large-field source radius changes by at
most \(CL^\Delta\) for a single source, and by at most four after
coalescence for the mixed sector. Thus the remainders have factors
\[
 L^{d-9\Delta}=L^{-5s/2}\quad(r=0),\qquad
 L^{-6\Delta+\Delta}=L^{-5s/2}\quad(r=1),\qquad
 4L^{-3\Delta}=4L^{-3s/2}\quad(r=2).
\]
Before coalescence the small-set mixed sector vanishes. Extra source
factors from the complement reference are bounded multipliers because
the mixed source weight is at most the product of the single weights;
they add only fixed field polynomials. Derivatives of field order above
\(n_r\) have at least the same restriction gain and are bounded directly
by the input large-field seminorm. Adjacent reference-coupling ratios
are \(1+O(G_j)\), which is absorbed by the last error term.
Combining the low jets and remainders proves the large-field estimate;
the zero-field estimate and norm equivalence prove the result.
\end{proof}

\begin{theorem}[The fractional one-step map]
\label{completion:map-theorem}
Let \(d=2s\in\{1,2,3\}\), and suppose the positive finite-range
covariance pieces have the bounds of Section~\ref{completion:map-support}
through total derivative order \(32\), with natural mass power two.
Use the concrete field norms above, field jets through order at least
nine, the four nilpotent sectors, and the explicit mixed large-field
source weight. Define the exact extraction, Gaussian integration,
second-order covariance redistribution, reblocking, and final
componentwise circle-product maps by the same algebraic formulas as
in Part II, with canonical mass dimension \(s\) and field dimension
\(s/2\). The local source domain is defined by bounded weighted source norms;
a fixed bound on all unweighted coefficients is not required for arbitrary
complex activity directions. The bulk potential domain has \(g\asymp G>0\) and normalized
mass and local kinetic coefficients bounded by fixed multiples of
\(G\). The reference couplings in consecutive norms obey
\(1-CG\le G_+/G\le1\); their values are fixed also at complex inputs.

Fix universal admissible \(k_0,\kappa\) and regulator powers first.
There are then choices of fixed \(L\), large-set parameter
\(\varrho\), and a joint complex ball radius \(c>0\),
for which the maps are holomorphic on
\[
 \|V-V_0\|<cG,\qquad \|K\|_{W_j}<cG^{9/4}.
\]
On fixed smaller balls, the activity map \(\Psi\) and bulk normal-form
remainder \(R\) satisfy
\begin{align*}
 \|\Psi(V,K)\|+|R(V,K)|
      &\le C(\chi_j^{3/2}G^3+\|K\|),\\
 \|D_KR(V,K)\|&\le C,\\
 \|D_K\Psi(V,0)\|
      &\le CL^{-\delta}+C_L\varrho^{\epsilon_0}
                     +C_L(\lambda^{1/2}+G^{1/4}),
       \qquad \delta=\min(s,2-s,1)>0,
\end{align*}
where \(\epsilon_0>0\) is fixed. In particular one can make the last
bound strictly smaller than \(1/32\), uniformly in the mass cutoff,
volume and sufficiently large spread radius. The Cauchy derivative
bounds of the enlarged-ball lemma hold as well. The constants depend
on the fixed norm choices, but not on the mass cutoff or volume.
\end{theorem}

\begin{proof}
The submultiplicative test-tensor estimates, localization at zero field,
and the finite nilpotent algebra are the same normed-algebra operations
in every dimension considered here. The only dimensional inputs to Gaussian integration are the finite
derivative bounds and the regulator comparison already proved.
Bulk localization retains the vacuum, mass, on-site quartic, and
isotropic local kinetic monomials. The kinetic coordinate is retained
even when irrelevant and is evolved with factor \(L^{s-2}\). Odd
quadratic derivative densities are put in canonical form modulo total
differences. Single-source localization retains the linear source-field
monomial at every scale. Only the mixed localization is constrained to vanish
on small sets before coalescence. Consequently projection to a single
nilpotent mark commutes exactly with the maps, independently of the
position of the unused second mark; this is the all-scales convention
needed for the distant-mark source bound. The small-set derivative is
Lemma~\ref{completion:smallset}. We give the remaining counting and
remainder bounds explicitly.

\paragraph{Regulator comparison for large sets.}
The quotient regulator also reblocks with a strict margin. On each fine
block choose the affine Taylor polynomial at its center. The discrete
second-order Taylor bound and the scaled derivatives through order
\(p_\Phi\) imply
\[
 \|\phi\|_{\widetilde\Phi_j(B_{\rm fine})}
    \le C L^{-2}\|\phi\|_{\widetilde\Phi_{j+1}(B)}.
\]
For derivative orders zero and one this is the Taylor remainder;
for orders at least two it is the factor \(L^{-|\gamma|}\le L^{-2}\).
After summing over at most \(L^d\) fine blocks and dividing by the
squared radii, the coefficient of the coarse quotient exponent is
\[
 CL^{d-4}(\ell_{j+1}/\ell_j)^2=CL^{s-4},
\]
which is strictly smaller than any fixed regulator margin for large
\(L\). Subsets of fine blocks only decrease the left side.
On occupied sites the product of the fine \(H\) factors has quadratic
coefficient larger than the coarse coefficient by \(L^s\). On
unoccupied sites the quartic reference, by quadratic Young's inequality,
supplies the missing coarse quadratic coefficient at a fixed exponential
cost per coarse block. The Gaussian shift uses the same determinant
bounds as the affine lemma, with a fixed exponential loss per input
block. Hence dense occupied sets cause no growing-regulator loss: one
uses the \(Y\) norm, the decaying \(H\), and this quotient comparison,
rather than attempting to compare products of growing regulators.

\paragraph{Circle-product summability.}
The nonlinear estimate uses a purely combinatorial bound. For a connected
polymer let \(A(X)=\varrho^{(|X|-2^d)_+}\). Connected shapes with
\(n\) fine blocks and a prescribed anchor number at most \(C_d^n\).
A depth-first walk through a connected set has length at most \(2n\);
split it into segments of length at most \(L/2\). Each segment meets at
most \(2^d\) coarse blocks. Therefore
\[
 |\overline X|\le2^d(1+4n/L).
\]
If \(n\) is bounded by a fixed multiple of \(2^d\), a sufficiently
large \(L\) makes \(|\overline X|\le2^d\). For larger \(n\), the
displayed inequality gives \(|\overline X|\le n/2\) after increasing
\(L\). Thus large input sets gain a positive power of \(\varrho\)
in comparison with their output size weight, and this gain absorbs
\(C_d^n\) when \(\varrho\) is small.
For a product of \(m\) nonempty component activities the elementary
count
\[
 \sum_{i=1}^m(|X_i|-2^d)_++2^d(m-1)
       \ge\left(\sum_i|X_i|-2^d\right)_+
\]
shows that component amplitudes at most \(\varrho^{2^d}\) supply the
missing large-set weight. Fixed exponential losses per block are
absorbed by the strengthened input weight margins. At zero field,
products of activities of norm at most \(\tau\) are bounded by
\(\tau^m\). At large field, with \(\tau=cG^{9/4}\), their weighted
sizes are at most
\[
 G^{9/4}(\tau G^{-9/4})^m=\tau c^{m-1}.
\]
Hence every circle product and connected expansion used in the exact
step is bounded by \(C\tau\), uniformly in the volume, for small fixed
\(c\). The same counting bounds derivatives, with the extra polynomial
choice factors absorbed by a smaller weight margin.

\paragraph{Second-order \(W\) terms and exact redistribution.}
The polynomial \(W_j\) is counted as degree two in the potential.
Its norm estimates always mean the polynomial multiplied by its stable
reference factor, or its zero-field polynomial majorant followed by
quartic/regulator absorption. The bare polynomial does not have a
finite decaying \(H\) norm. In this precise sense,
\begin{equation}\label{completion:W-bounds}
 \|e^{-V}W_j\|_{\ell,G}\le C\chi_jG^2,\qquad
 \|e^{-V}W_j\|_{h,H\widetilde G}\le C\chi_jG^{1/2}.
\end{equation}
Here and below stable interpolated references may replace \(e^{-V}\).

For clarity, the scalar quartic--quartic Wick contractions give the
following finite list. Write \(w=w_j\), suppress fixed combinatorial
constants, and use the canonical localization of the quadratic symbol:
\[
\begin{array}{c|c|c|c}
\text{contractions}&\text{external fields}&
 \text{normalized covariance moment}&\text{large-field size}\\ \hline
1&6&b_j^{-s}\|w\|_1& O(\chi_jG^{1/2})\\
2&4&b_j^{-1}\sum_x|x|w(x)^2&O(\chi_jG)\\
3&2&b_j^{s-3}\sum_x|x|^3|w(x)|^3&O(\chi_jG^{3/2})\\
4&0&\text{removed by localization}&0
\end{array}
\]
All three small-field sizes are \(O(\chi_jG^2)\). The second row
uses the first spatial Taylor remainder after quartic localization;
using \(\sum w^2\) instead would introduce an inadmissible logarithm.
The third row uses the third spatial remainder after mass and kinetic
localization. A second remainder is enough for a noncanonical bound,
but the third remainder makes the postmass estimate immediate.
The pointwise finite-range slice sum gives
\[
 |\nabla^rw_j(x)|\le C_r\lambda(b+|x|)^{-s-r}
                  (1+\mathfrak m|x|^s)^{-2}.
\]
If \(\xi=b_J\), the second-row moment is at most \(C\lambda^2\xi\)
and the third-row moment at most \(C\lambda^3\xi^{3-s}\) after the
mass scale. Their normalized decays are respectively
\(L^{-(j-J)}\) and \(L^{-(3-s)(j-J)}\), while the first row decays
as \(L^{-s(j-J)}\). These dominate \(\chi_j\) for a sufficiently
large fixed \(L\).

The remaining Wick cases do not require a new dimension-dependent
estimate. A mass vertex supplies \(\nu b_j^s=O(G)\). A kinetic vertex
supplies \(z_{\rm phys}b_j^{s-2}=O(G)\); move its differences by
summation by parts to external fields whenever a single covariance is
contracted. For two contractions, use the corresponding weighted
moments of \(w\nabla^2w\) and \((\nabla w)^2\) after quadratic
localization. For example,
\(\sum|x|^3(|w\nabla^2w|+|\nabla w|^2)\le C\lambda^2b_j\),
with a bound \(C\lambda^2\xi\) after the mass scale. This is a direct
shell sum of the displayed pointwise derivative bound. A
kinetic--kinetic one-contraction term can be written with four external
differences and \(w\); its normalized row moment has the first-row
bound. All zero-row-sum derivative terms have no hidden mass constant.
Vacuum terms are localized away. Single-source terms have only one
possible contraction; source--quartic terms use the first row, and
source--quadratic terms use the first spatial remainder or the
external kinetic differences. The mixed source contraction is a constant
and is extracted. Stopped source radii only decrease these bounds.
This proves \eqref{completion:W-bounds} for every sector.

Introduce a formal parameter \(t\) multiplying the potential. The exact
second-order identity is
\[
 \mathbb E[e^{-tV}(1+t^2W_j)]
 =1-t\mathbb EV+t^2\left(\tfrac12\mathbb EV^2+\mathbb EW_j\right)
       +O(t^3).
\]
Put \(R_2=\tfrac12\operatorname{Cov}(V,V)+\mathbb EW_j\).
The quadratic correction to the output potential is
\(-\operatorname{Loc}R_2\), and the unlocalized second-order factor is
\((1-\operatorname{Loc})R_2\). Expanding
\[
 e^{-t\mathbb EV+t^2\operatorname{Loc}R_2}
       [1+t^2(1-\operatorname{Loc})R_2]
\]
shows exactly the same coefficients through degree two. This is the
covariance redistribution used in Maps~1--5; it is an algebraic identity
and does not depend on the field dimension.

\paragraph{All higher-order terms.}
Centered potential increments have the respective sizes
\[
 \ell:\quad C_L\chi_j^{1/2}G_j,\qquad
 h:\quad C_L\chi_j^{1/2}G_j^{1/4}.
\]
There are four kinds of remainders after the exact second-order
redistribution:
\begin{itemize}
\item Cubic integral remainders in centered potential increments have
respective sizes
\[
 C_L\chi_j^{3/2}G_j^3,\qquad C_L\chi_j^{3/2}G_j^{3/4}.
\]
\item One \(W_j\) times a centered increment has the same two sizes,
by \eqref{completion:W-bounds}. Terms with two \(W\) factors have
sizes at most \(C_L\chi_j^2G_j^4\) and
\(C_L\chi_j^2G_j\), which are smaller after the large-field weight.
\item Bilinear differences from extraction obey
\[
 \|W(V-Q)-W(V)\|_{\ell,\mathrm{stable}}
       \le C\chi_jG\tau,\qquad
 G^{9/4}\|W(V-Q)-W(V)\|_{h,\mathrm{stable}}
       \le C\chi_jG^{1/2}\tau.
\]
This follows from bilinearity, \(\|Q\|_\ell\le C\tau\),
\(\|Q\|_h\le C\tau G^{-9/4}\), and the covariance contraction
factor \(\ell^2/h^2=O(G^{1/2})\). Field-independent source constants
do not enter \(W\).
\item In the final canonical chart let
\(a=\chi_jG^2+\tau+\|K\|\) bound the redundant stiffness difference.
The singleton difference has stable bounds
\(C\chi_jGa\) at \(\ell\) and
\(C\chi_jG^{-1/2}a\) at \(h\). Its exponential multiplier is bounded
by the spare regulator powers. After multiplication by \(G^{9/4}\),
these bounds are admissible in the same output norm.
\end{itemize}
Gaussian expectation of products of the segment remainders is bounded
by the determinant/Sobolev argument even when segment parameters differ
between blocks. Thus all pure-potential outputs have mixed size at most
\(C_L\chi_j^{3/2}G_j^3\). This conclusion counts \(W\) as degree two;
it does not assert that every higher-order term contains three fresh
centered increments.

For large output polymers the required size weight is also obtained
from the perturbative factors. Each single-block increment has size at
most \(C_LG^{1/4}\) at \(h\), and each quadratic \(W\) factor has size
at most \(C_LG^{1/2}\). Finite covariance range below a fixed fraction
of the coarse block implies that each factor covers only a fixed
number of coarse blocks. A connected contribution occupying \(n\)
coarse blocks therefore contains at least \(cn-C_0\) such perturbative
orders, for fixed positive \(c,C_0\). Retain three orders for the
\(G^3\) bound above; the remaining factor
\((C_LG^{1/4})^{cn-C_0}\) supplies any prescribed strengthened
large-set weight when \(R\) is large after \(\varrho\) is fixed.
There are at most a fixed exponential number of connected choices.
Boundedly many exceptional sizes are absorbed into \(C_\varrho\).
Terms containing input activities use the preceding circle-product
count and are bounded by \(C_L\|K\|\).

The exact Maps~1--5 formulas, their second-order covariance
redistribution, and the componentwise final chart change in Part II
are algebraic identities on these same normed algebras. Localization
uses bounded finite zero-field inverses; the quartic remains stable on
the complex potential disc. Consequently the preceding bounds give
holomorphy and the uniform rough estimate on the whole joint ball,
not only at zero activity. Combining the small-set lemma with the
large-set circle count gives
\[
 \|D_K\Psi(V,0)\|
       \le CL^{-\delta}+C_L\lambda^{1/2}+C_LG^{1/4}
          +o_\varrho(1).
\]
The normal-form bulk coefficient bounds follow from
Lemma~\ref{completion:moments}. Fix \(k_0,\kappa\) and the regulator powers universally first; then
choose \(L\), then \(\varrho\), then the complex ball radius \(c\),
and finally the
spread radius so that \(\lambda,G\) are small. This makes the zero-activity derivative at most \(1/32\). Choose
\(\epsilon\) afterwards so that the enlarged-ball Cauchy correction is
less than \(1/32\). This gives \(\|D_K\Psi(V,K)\|<1/16\), uniformly
on \(\|K\|\le\epsilon G^{9/4}\), as required by the recursion.

\end{proof}

\section{A uniform recursion with an irrelevant local kinetic coordinate}
\label{frac:sec:recursion}
The following is an algebraic result. Its one-step hypotheses are verified by Theorem~\ref{completion:map-theorem} for the covariance constructed above. Put \(s=d/2\), with \(d\in\{1,2,3\}\), and fix \(L\) so
large that \(a_z=L^{s-2}<1/8\). Fix \(d_0=1/4\).

\begin{proposition}[Uniform fractional normal-form tuning]
\label{frac:prop:recursion}
Let \(\chi_j=2^{-(j-J)_+}\), where
\(J\in\mathbb N\cup\{\infty\}\), and suppose
\[
 x_{j+1}=x_j-b_jx_j^2,\qquad
 0\le b_j\le C\lambda^2\chi_j,\qquad
 b_j\ge c\lambda^2\quad(j\le J).
\]
Set \(f_j=\chi_j^{3/2}x_j^3+Ud_0^j\).
Consider the recursion
\begin{align*}
 g_{j+1}&=g_j-b_jg_j^2+e_j^g(v_j,K_j),\\
 z_{j+1}&=a_z z_j-\theta_jg_j^2+e_j^z(v_j,K_j),\\
 \mu_{j+1}&=L^s\mu_j+\eta_jg_j-\xi_jg_j^2
       -\omega_jg_j\mu_j-\pi_jg_jz_j+e_j^\mu(v_j,K_j),\\
 K_{j+1}&=P_j(v_j,K_j),\qquad v_j=(g_j,z_j,\mu_j),
\end{align*}
where
\[
 |\eta_j|\le C\lambda\chi_j,\qquad
 |\theta_j|+|\xi_j|\le C\lambda^3\chi_j,\qquad
 |\omega_j|+|\pi_j|\le C\lambda^2\chi_j.
\]
Assume that on the boxes
\[
 |g_j-x_j|\le\varepsilon x_j,\quad
 |z_j|+|\mu_j|\le M\lambda x_j,\quad
 \|K_j\|\le M_K f_j
\]
the maps obey
\begin{align*}
 |e_j(v,K)|&\le C(\chi_j^{3/2}x_j^3+\|K\|),\\
 \|P_j(v,0)\|&\le C\chi_j^{3/2}x_j^3,\qquad
 \|P_j(v,K)-P_j(v,K')\|\le\kappa\|K-K'\|,
                    \quad \kappa<1/8,\\
 |e_j(v,K)-e_j(v',K')|
   &\le C_{M,M_K}\bigl(f_j|v-v'|/x_j+\|K-K'\|\bigr),\\
 \|P_j(v,K)-P_j(v',K)\|
   &\le C_{M,M_K}f_j|v-v'|/x_j.
\end{align*}
The first two constants are uniform over fixed coordinate boxes;
constants in the last two lines may depend on their fixed sizes.
Suppose
\[
 \lambda\to0,\qquad x_0/\lambda^3\to0,\qquad
 U/(\lambda x_0)\to0.
\]
For sufficiently small parameters and every prescribed
\(\|K_0\|\le U\), there is a controlled solution with
\(g_0=x_0\), \(z_0=0\), and \(\mu_j\to0\), satisfying
\[
 |g_j-x_j|\le\varepsilon x_j,\qquad
 |z_j|+|\mu_j|\le C\lambda\chi_j x_j,\qquad
 \|K_j\|\le C f_j.
\]
It is the unique solution in a fixed box of these weighted sizes.
The contraction factor and all constants are independent of \(J\).
For a finite horizon the same conclusion holds, with \(\mu_T=0\)
in place of the condition at infinity; \(z\) retains its initial
condition \(z_0=0\) and is not assigned an additional terminal value.
If the maps, reference data and prescribed initial activity are
continuous in external parameters in a common Banach-space chart,
the tuned initial mass is continuous as well.
\end{proposition}

\begin{proof}
Telescoping the reference recursion gives, uniformly in \(J\),
\begin{equation}\label{frac:eq:recursion-sums}
 \frac{x_0}{1+C\lambda^2x_0j}\le x_j\le x_0,
 \qquad \sum_{i\ge j}\chi_i x_i^p
       \le C_p\lambda^{-2}x_j^{p-1}\quad(p=2,3).
\end{equation}
Before the cutoff use \(b_i\ge c\lambda^2\); afterwards the
\(\chi_i\)-sum is geometric and \(x_i\asymp x_{J+1}\).
For \(A_{j,i}=\prod_{i<l<j}(1-2b_lx_l)\),
\(1-2t\le(1-t)^2\) gives
\begin{align}
 A_{j,i}&\le(x_j/x_{i+1})^2,\qquad
 \sum_{i<j}A_{j,i}b_ix_i^2\le Cx_j,\nonumber\\
 \sum_{i<j}A_{j,i}f_i
   &\le Cx_j\left(\frac{x_0}{\lambda^2}+\frac U{x_0}\right).
      \label{frac:eq:forward-sum}
\end{align}
For the last estimate, before \(J\) divide the summands by
\(x_j^2\) and sum \(x_i\), using
\(\log(x_i/x_{i+1})\ge b_ix_i\) and
\(t\log(x_0/t)\le Cx_0\). After \(J\), use the geometric
\(\chi_i\)-sum. The transient uses the lower bound in
\eqref{frac:eq:recursion-sums} and
\(\sum_i d_0^i(1+i)^2<\infty\).

Put
\[
 a=\frac{x_0}{\lambda^2}+\frac U{x_0},\qquad
 c_* =\frac{x_0^2}{\lambda}+\frac U{\lambda x_0}.
\]
Both tend to zero under the stated hypotheses. Since
\(f_{j+1}/f_j\ge1/4\) for small parameters and \(a_z<1/8\),
\[
 \sum_{i<j}a_z^{j-1-i}f_i\le C f_j.
\]
Also, writing \(w_j=\chi_jx_j\), the adjacent ratios of \(w_j\)
are at least \(\tfrac12(1-C\lambda^2x_0)>1/4\), so
\[
 \sum_{i<j}a_z^{j-1-i}\chi_i x_i\le C\chi_jx_j,
 \qquad
 \sum_{i\ge j}L^{-s(i-j+1)}f_i\le C f_j.
\]
The latter estimate follows from monotonicity of \(f_i\) and the
geometric mass Green function. In particular,
\begin{equation}\label{frac:eq:f-over-mass-weight}
 \sup_j\frac{f_j}{\lambda\chi_jx_j}\le Cc_*.
\end{equation}
For its transient term use \(\chi_j\ge2^{-j}\) and the reference
lower bound, which give
\(d_0^j x_0/(\chi_jx_j)\le
(2d_0)^j(1+C\lambda^2x_0j)\le C\).

Write \(t_j=g_j-x_j\). Define a fixed-point map by solving the
linearized \(t\) equation forward from \(t_0=0\), with forcing
\(-b_jt_j^2+e_j^g\). Solve the stable kinetic equation forward:
\[
 z_j=\sum_{i<j}a_z^{j-1-i}(-\theta_i g_i^2+e_i^z).
\]
Writing \(F_i^\mu=\eta_i g_i-\xi_i g_i^2-
\omega_i g_i\mu_i-\pi_i g_i z_i+e_i^\mu\), solve the mass
backward:
\[
 \mu_j=-\sum_{i\ge j}L^{-s(i-j+1)}F_i^\mu.
\]
The output activity is the original forward update
\(K_{j+1}=P_j(v_j,K_j)\). All forcings on the right use the input
sequence. For a finite horizon truncate the backward mass sum at
\(T-1\); the other formulas are unchanged.

Use the maximum of the four sup norms divided respectively by
\(\varepsilon x_j,M\lambda x_j,M\lambda x_j,M_Kf_j\).
Choose \(M_K\) so large that \(4C/M_K+4\kappa<1\), and then
choose \(M\) large enough for the leading mass forcing
\(\eta_jg_j\). Separate bounds on \(z,\mu\) fit the hypotheses
with \(M\) replaced by \(2M\). The kinetic forcing has normalized
size \(O(\lambda^2x_0+c_*)\); the nonleading mass forcings have
normalized size \(o(1)\), by
\eqref{frac:eq:f-over-mass-weight}. The normalized forward quartic
output is at most \(C\varepsilon+Ca/\varepsilon\), by
\eqref{frac:eq:forward-sum}. Thus a small fixed \(\varepsilon\),
followed by small parameters, makes the product box invariant.

For explicit contraction estimates let
\(D_t,D_z,D_\mu,D_K\) denote the four normalized input distances.
The forward Green estimate yields
\[
 D_t^{\rm out}\le(C\varepsilon+o(1))D_t
                       +o(1)(D_z+D_\mu+D_K).
\]
Here the error coefficients are bounded by fixed multiples of
\(a,M\lambda a/\varepsilon,M_Ka/\varepsilon\).
For the forward kinetic coordinate,
\[
 D_z^{\rm out}\le o(1)(D_t+D_z+D_\mu+D_K).
\]
Indeed its leading \(t\)-coefficient is
\(O((\varepsilon/M)\lambda^2x_0)\); its error coefficients are
bounded by fixed multiples of
\((\varepsilon/M)c_*,\lambda c_*,(M_K/M)c_*\).
For the mass coordinate the only nonvanishing leading row entry is
from \(\eta_jg_j\):
\[
 D_\mu^{\rm out}\le(C\varepsilon/M+o(1))D_t
                         +o(1)(D_z+D_\mu+D_K).
\]
The other perturbative mass coefficients contain
\(\lambda^2x_0\) or \(\lambda^2x_0^2\), and its error entries
are controlled by \(c_*\). Finally,
\[
 D_K^{\rm out}\le4\kappa D_K+
 \frac{4C_{M,M_K}}{M_K}
       \bigl(\varepsilon D_t+M\lambda(D_z+D_\mu)\bigr).
\]
Choose \(\varepsilon\) after the fixed box sizes, then make the
parameters small. The maximum row sum is bounded by some \(q<1\)
independent of \(J\) and of the finite horizon. Banach's theorem
therefore gives the fixed point and uniqueness in the box.

At that fixed point, all mass forcings are bounded by
\(C\lambda\chi_jx_j\), using
\eqref{frac:eq:f-over-mass-weight} and the preliminary box bounds.
The backward geometric sum preserves this weight. The kinetic
forcing has the same bound and its forward geometric sum preserves
it as well. This proves the stronger displayed estimates with
\(\chi_j\). They tend to zero at a finite mass cutoff; when
\(J=\infty\), \(x_j\to0\) by the lower bubble bound.
For continuity, iterate the same uniformly contracting sequence map.
At each fixed scale its forward coordinates use finite sums. The only
backward sum is the mass sum, which has a locally uniform geometric
majorant \(CL^{-s(i-j+1)}\lambda x_0\). Dominated convergence
therefore makes every fixed coordinate of each finite iterate
continuous. The common contraction factor makes convergence of the
initial mass uniform on a small parameter neighborhood, and proves
its asserted continuity.
\end{proof}


\section{Identifying a tuned endpoint for a summable interaction}
\label{frac:sec-endpoint}

The endpoint argument does not require a finite interaction range.
It does require a summable majorant for the periodic correlations.
For a fractional interaction this majorant is a renewal kernel;
exponential decay would be false even at a positive auxiliary mass.
The following statement separates this issue from the RG estimates
needed to obtain the susceptibility input.

\begin{proposition}[A summable-interaction endpoint criterion]
\label{frac:endpoint}
Let \(J:\mathbb Z^d\to[0,\infty)\) be symmetric and translation
invariant, with \(J(0)=0\), \(\sum_xJ(x)<\infty\), and with its
support generating \(\mathbb Z^d\). On the torus
\(\Lambda_D=(\mathbb Z/D\mathbb Z)^d\), use the periodized
interaction
\[
 J_D(\bar x)=\sum_{n\in\mathbb Z^d}J(x+Dn)
 \quad(\bar x\ne0);
\]
the diagonal term is omitted, since it changes only the partition
function's normalization. Write
\(\chi_D(\beta)=\sum_{\bar x\in\Lambda_D}
\langle\sigma_0\sigma_{\bar x}\rangle_{D,\beta}\).
Let \(I_0\subset(0,\infty)\) be a compact interval. Suppose that,
for every sufficiently small auxiliary mass \(\mu>0\), there are
tuned parameters \(\beta_{\mu,D}\in I_0\) along tori with
\(D\to\infty\), and that
\begin{equation}\label{frac:eq-susceptibility-input}
 \chi_D(\beta_{\mu,D})
 =\frac{a(\beta_{\mu,D})}{\mu}-b(\beta_{\mu,D})
       +O_\mu(D^{-s}),\qquad s>0,
\end{equation}
where \(a,b\) are continuous on \(I_0\) and
\(\inf_{I_0}a>0\). The error constant is uniform in the chosen
tuned parameters at fixed \(\mu\).

For any subsequential limit
\(\beta_{\mu,D}\to\beta_\mu\), one has
\[
 \beta_\mu<\beta_c(J),\qquad
 \chi^{\mathrm f}(\beta_\mu)
       =\frac{a(\beta_\mu)}{\mu}-b(\beta_\mu),
\]
and the periodic two-point functions converge locally and in the
summed sense to the free infinite-volume two-point function at
\(\beta_\mu\). Every choice of these subsequential limits satisfies
\[
 \beta_\mu\longrightarrow\beta_c(J)
       \qquad(\mu\downarrow0).
\]
For each fixed pair of sites their free correlations therefore
converge to the critical free correlations. No monotonicity of
\(\beta_\mu\), or uniqueness of the tuning, is assumed.
\end{proposition}

\begin{proof}
We use the Griffiths inequalities, the Simon--Lieb inequality, and
the standard sharpness identification of the susceptibility
threshold with the spontaneous-magnetization critical point for
summable ferromagnetic Ising interactions; see \cite{DCT}.
In particular susceptibility is finite and the Gibbs state is
unique at every fixed strictly subcritical inverse temperature.

Fix \(\mu>0\) and a subsequence with
\(\beta_{\mu,D}\to\beta_\mu\). Equation
\eqref{frac:eq-susceptibility-input} bounds the periodic
susceptibilities uniformly along this subsequence. If
\(\beta<\beta_\mu\), every fixed finite free box can eventually
be embedded in \(\Lambda_D\) with
\(\beta\le\beta_{\mu,D}\). Its pair interactions are no larger
than the corresponding periodized interactions. The Griffiths
inequalities therefore bound its partial susceptibility by the
periodic susceptibility. Increasing the free boxes gives
\(\chi^{\mathrm f}(\beta)\le C_\mu\). Moreover,
\begin{equation}\label{frac:eq-free-left-continuity}
 \chi^{\mathrm f}(\beta_\mu)
   =\sup_{\beta<\beta_\mu}\chi^{\mathrm f}(\beta).
\end{equation}
Indeed free susceptibility is the supremum of the finite free-box
partial susceptibilities, each of which is continuous and
increasing in \(\beta\); the two suprema commute. This proves
\(\chi^{\mathrm f}(\beta_\mu)<\infty\) before any uniqueness
or summability of a limiting periodic state is used.

For a finite box \(S\) containing the origin put
\[
 \varphi_\beta(S)=\sum_{u\in S,\,v\notin S}
 \langle\sigma_0\sigma_u\rangle^{\mathrm f}_{S,\beta}
                  \tanh(\beta J(v-u)).
\]
There is a box \(S\) with \(\varphi_{\beta_\mu}(S)<1\).
To see this, bound the free-box correlation by its infinite-volume
value at \(\beta_\mu\), and use
\(\tanh(\beta J)\le\beta J\). For each fixed \(u\), the
sum over \(v\notin S\) tends to zero as the boxes increase.
It is bounded by \(\beta_\mu\|J\|_1\), so dominated convergence
with the summable function
\(u\mapsto\langle\sigma_0\sigma_u\rangle^{\mathrm f}_{\beta_\mu}\)
proves \(\varphi_{\beta_\mu}(S)\to0\).
For fixed \(S\), \(\varphi_\beta(S)\) is continuous in
\(\beta\): the correlations are finite-volume continuous, and
the remaining series is dominated on compact intervals by a
constant times \(J\). Choose
\(\beta_+>\beta_\mu\) and \(a_1<1\) with
\(\varphi_{\beta_+}(S)<a_1\). The finite-set criterion for
sharpness gives \(\beta_+\le\beta_c(J)\), and hence
\(\beta_\mu<\beta_c(J)\).

We next prove the uniform summability needed for the moving-temperature
limit. Choose \(\varepsilon>0\) so small that the nonnegative
exit kernel
\begin{equation}\label{frac:eq-exit-kernel}
 T_*(v)=\boldsymbol1_{\{v\notin S\}}
   \sum_{u\in S}
   \bigl(\langle\sigma_0\sigma_u\rangle^{\mathrm f}_{S,\beta_+}
                    +\varepsilon\bigr)
       \tanh(\beta_+J(v-u))
 \quad\text{satisfies}\quad \|T_*\|_1=a_*<1.
\end{equation}
Such a choice is possible because the added mass is at most
\(\varepsilon |S|\beta_+\|J\|_1\).
Embed \(S\) in all sufficiently large tori. The additional image
bonds between any two of its sites tend to zero, by summability
of \(J\). Finite-volume continuity and monotonicity imply,
uniformly for \(0\le\beta\le\beta_+\),
\[
 \langle\sigma_0\sigma_u\rangle^{\mathrm f}_{S,J_D,\beta}
 \le \langle\sigma_0\sigma_u\rangle^{\mathrm f}_{S,J,\beta_+}
                      +\varepsilon\qquad(u\in S)
\]
when \(D\) is sufficiently large. Here the subscript \(S,J_D\)
means that only the periodized bonds internal to \(S\) are kept.
Define the torus exit kernel
\[
 T_{D,\beta}(\bar v)=\boldsymbol1_{\{\bar v\notin S\}}
 \sum_{u\in S}
  \langle\sigma_0\sigma_u\rangle^{\mathrm f}_{S,J_D,\beta}
                         \tanh(\beta J_D(\bar v-u)).
\]
For nonnegative numbers, \(\tanh(\sum_i t_i)\le\sum_i\tanh t_i\).
Using this inequality and \eqref{frac:eq-exit-kernel} gives
\[
 T_{D,\beta}\le\mathcal P_DT_* ,\qquad
 (\mathcal P_Df)(\bar x)=\sum_{n\in\mathbb Z^d}f(x+Dn).
\]
If \(\bar v\notin S\), every representative \(v+Dn\) is
outside the embedded box \(S\), so the indicator in
\eqref{frac:eq-exit-kernel} causes no loss in this comparison.
For \(\bar v\in S\) its left side is zero.
In particular \(\|T_{D,\beta}\|_1\le a_*<1\).

Let \(G_{D,\beta}(\bar x)=
\langle\sigma_0\sigma_{\bar x}\rangle_{D,\beta}\).
The Simon--Lieb inequality outside \(S\), and the bound
\(G_{D,\beta}\le1\) inside \(S\), give on the whole torus
\[
 G_{D,\beta}\le\boldsymbol1_S+T_{D,\beta}*G_{D,\beta}.
\]
Iterating this inequality, the remainder after \(n\) iterations
has supremum at most \(a_*^n\). Nonnegative periodization
commutes with convolution. Consequently
\begin{equation}\label{frac:eq-renewal-envelope}
 G_{D,\beta}\le\mathcal P_DH,\qquad
 H=\left(\sum_{n\ge0}T_*^{*n}\right)*\boldsymbol1_S,
 \qquad \|H\|_1\le\frac{|S|}{1-a_*}<\infty,
\end{equation}
for all sufficiently large \(D\) and \(\beta\le\beta_+\).
This is a summable bound; no exponential bound is asserted.
Take one shortest representative of each torus point. For every
\(K>0\), each image of a representative with \(|x|>K\) also
has norm greater than \(K\). Hence
\[
 \sum_{\substack{x\text{ shortest representative}\\|x|>K}}
       G_{D,\beta}(\bar x)
 \le\sum_{|z|>K}H(z)\longrightarrow0
       \qquad(K\to\infty),
\]
uniformly in these volumes and temperatures.

Since \(\beta_\mu<\beta_c(J)\), the infinite-volume Gibbs state
there is unique. Every local subsequential limit of the periodic
laws at \(\beta_{\mu,D}\) satisfies the DLR equations at
\(\beta_\mu\). For a summable interaction this follows by
first truncating the interactions entering a fixed conditional
law to a finite neighborhood; the discarded contribution is
bounded uniformly by a constant times the tail of \(J\).
The finite-neighborhood equations pass to the local limit, and
then the truncation is removed. Uniqueness therefore gives
local convergence to the free state. The uniform tail bound in
\eqref{frac:eq-renewal-envelope} permits summation, proving
\[
 \chi^{\mathrm f}(\beta_\mu)
 =\lim_{D\to\infty}\chi_D(\beta_{\mu,D})
 =\frac{a(\beta_\mu)}{\mu}-b(\beta_\mu).
\]

The compact parameter interval keeps \(a\) bounded away from
zero and \(b\) bounded, so these susceptibilities diverge as
\(\mu\downarrow0\). Each \(\beta_\mu\) is strictly below
\(\beta_c(J)\). If an accumulation point were smaller than
\(\beta_c(J)\), choose a fixed strictly subcritical temperature
between them. Monotonicity and subcritical susceptibility
finiteness would give a common finite upper bound along that
subsequence, a contradiction. Thus every accumulation point
is \(\beta_c(J)\).
Finally, the critical free pair correlation is the supremum
of its finite-free-volume correlations. As for
\eqref{frac:eq-free-left-continuity}, continuity of those finite
correlations and commutation of the two increasing suprema give
left continuity at \(\beta_c(J)\). Monotonicity then gives
convergence along any of the tuned temperatures just considered,
including choices that are not monotone.
\end{proof}

\section{Initialization, physical tuning, and observable closure}
\label{completion:flow-closure}
We describe the consequences of the one-step theorem, keeping its use
separate from the covariance construction. In this section the
one-step estimates of Theorem~\ref{completion:map-theorem} are used
for the exact four-sector map, with constants independent of the
spread radius, mass cutoff and torus volume.

\subsection{Finite-range reference and the initial activity}
The even-tail split has
\(\|\beta Q\|_1\le q_*<1\), range at most \(C_NR\), and
\(\|Q\|_\infty=O(R^{-d})\). The diagonal of \(Q\) is also
\(O(R^{-d})\). In the cavity comparison it is an omitted self-edge,
whose contribution is bounded by the same \(O(R^{-d})\) error as
an omitted ancestor edge. The proof of Proposition~\ref{prop:cavity}
therefore gives, with \(g=R^{-d}\),
\[
 \sup_i\sum_x |(S-S_0)(i,x)|e^{\eta|i-x|/R}\le Cg,
 \qquad S=D^2F(0),\quad S_0=(I-\beta Q)^{-1}.
\]
Indeed every branch of the cavity recursion has an error at depth
\(n\) at most \(C(n+1)g\), and its weighted row contraction is
\(q'<1\). Summing \(Cg\sum(n+1)(q')^n\) proves the bound.
The joint field/link estimates and real coercivity use the same
recursion and contain no spatial-dimensional identity. With
\(B=I-\beta Q\), the exact change \(h=Bu\) makes the leading
mean-field quartic exactly \(\sum u_i^4/12\), independently of
all reference interpolation links. The residual Hessian satisfies
\eqref{complete:K-assumption}, so Proposition~\ref{complete:covariance-prop}
applies to its nonconstant part.

Choose \(b\asymp R\), with its fixed aspect sufficiently large
for the finite-range reference and covariance collars. The all-field
connected initialization of Proposition~\ref{prop:transformed-init}
uses only these cavity bounds and
\[
 |B|h_u^4\asymp1,\qquad h_u\asymp R^{-d/4},\qquad
 \ell'_u\asymp A_{\rm sc}R^{-d/2},\qquad \ell'_u/h_u\to0.
\]
Its Taylor and regulator estimates in these dimensions are explicitly
\[
 |B|g(\ell'_u)^4+|B|(\ell'_u)^6
       \le C A_{\rm sc}^6g^2,
 \qquad |B|g h_u^4+|B|h_u^6\le Cg^{1/2}.
\]
The joint link coefficient extraction is unchanged: its graph has
bounded degree depending only on \(d,N\), the adaptive radii are held
fixed during differentiation, and the forest factorization uses the
finite-range \(B_t\), not the long-range Gaussian numerator. Thus
in the fractional norm, including its sources,
\[
 \|K_0\|\le U=C(A_{\rm sc}^6g^2+G_0^{9/4}g^{1/2}),\qquad
 G_0=A_{\rm sc}^4g/3,\qquad \lambda=A_{\rm sc}^{-2}.
\]
Take \(A_{\rm sc}=g^{-a_*}\), \(1/12<a_*<1/10\).
The ratios \eqref{frac:eq:seed-small}--\eqref{frac:eq:seed-ratios}
put the initial activity inside the enlarged ball and ensure the
smallness conditions of Proposition~\ref{frac:prop:recursion}.

\subsection{Mass matching at finite volume}
Use the Hessian lift and periodization described in
Section~\ref{sec:finitevolume}. Only the reference Hessian is lifted;
its exponential bound follows from the finite-range reference. The
fractional numerator itself is periodized directly. The auxiliary-space
construction commutes with periodization, since each polynomial
piece is a translation-invariant finite-range kernel.

The potential Wick formulas are algebraic Gaussian identities. The
bounds in Lemma~\ref{completion:moments} give the bubble and cubic
coefficients required in Proposition~\ref{frac:prop:recursion}.
In physical field units the local kinetic coefficient is unchanged
by reblocking; after multiplication by \(b_j^{s-2}\) its linear
factor is \(L^{s-2}\). Thus that proposition applies with \(z_0=0\),
forward evolution of \(z\), and only a backward terminal condition
for the mass. The initial normal-form change is the identity because
the cumulative covariance is zero. For each external covariance
parameter it returns a continuous initial mass
\(\mu_0^*\) with \(|\mu_0^*|\le C\lambda G_0\), uniformly in
the finite horizon and the auxiliary mass.

Let \(r\) be the row sum of the residual Hessian before extracting
its zero-row-sum part. As in the exact algebra of
Section~\ref{sec:matching-physical},
\[
 r(\beta,\mu)
 =\frac{(1-\beta q)(1-\mu-\beta)}{\beta p}
      -(1-\beta q)^2\widehat E(0),\qquad
 E=S-S_0.
\]
The complex cavity comparison gives \(\partial_\beta\widehat E(0)=O(g)\),
so \(\partial_\beta r=-1+O(g)\) on the physical box. The bare
normalized mass is
\(\mu_0^{\rm bare}=b^s r/(\lambda R^s)\). Therefore the continuous map
\[
 r\longmapsto (R/b)^s\lambda\mu_0^*(r)
\]
has image \(O(\lambda^2G_0)=O(g)\). The constant can be chosen
before the interval size: use a fixed comparison domain
\(|\beta-1|+\mu+\|K\|_{\exp}\le\varepsilon_0\) for the covariance
bounds. The initialization absorbs the mass factor uniformly once
\(C_1g^{1/2}\) is small. Thus none of the image constants needs to
grow with \(C_1\). First choose a sufficiently large fixed interval
\([-C_1g,C_1g]\), then take \(R\) large.
It is a self-map and has a fixed point. Continuity here can be checked
without assuming a common variable-weight Banach chart: at fixed finite
horizon, every finite fixed-point iterate is continuous in the local
field jets by Gaussian domination and the finite algebraic maps. The
common contraction makes the scalar initial-mass iterates converge
uniformly on the physical interval. Their limit is continuous. At the fixed point the
bare physical mass agrees with the tuned mass, whereas \(g_0=G_0\)
and \(z_0=0\) already agree. No second physical tuning of a marginal
local stiffness is used in \(d<4\). The argument works uniformly
on finite tori with terminal normalized mass zero.

\subsection{The final Gaussian integral}
At fixed positive auxiliary mass, let the torus side tend to infinity
and stop the local flow with \(L^d\) blocks remaining. For a covariance
piece of physical range \(r_j\asymp RL^j\) beyond the mass scale,
Proposition~\ref{complete:covariance-prop} implies
\[
 \|\nabla^k C_j\|_\infty
       \le C\lambda\mu^{-2}R^{2s}r_j^{-(d+s+k)}.
\]
Periodization on a torus of side \(D_N\) costs at most
\(C(1+r_j/D_N)^d\) images. Summing the pieces with
\(r_j\gtrsim D_N\) gives
\[
 \|\nabla^k C_{\rm tail}^{\Lambda_N}\|_\infty
       \le C\lambda\mu^{-2}R^{2s}D_N^{-(d+s+k)}.
\]
After division by the squared fluctuation radius and its derivative
units this is \(O_{R,\mu}(D_N^{-d})\), hence tends to zero.
The Sobolev/determinant argument for the fixed number of remaining
blocks gives the last-scale regulator moments and
\[
 Z_{\rm tail}(0)\ge\tfrac12,\qquad
 \|Z_{\rm tail}\|_{T_{\ell_{N-1}}(0)}\le C_L.
\]
It follows that
\[
 |\Lambda_N|^{-1}
 |D^2_{\boldsymbol1,\boldsymbol1}\log Z_{\rm tail}(0)|
       \le C_LD_N^{-s},
\]
since \(d=2s\) and \(\ell_{N-1}^{-2}\asymp D_N^s\).
The exact row-sum susceptibility calculation in
\eqref{frac:suscept-identity} and Proposition~\ref{frac:endpoint}
then identify every massless tuned temperature with the physical
critical temperature. This argument uses the summable renewal
envelope, not exponential decay of massive fractional correlations.

At fixed \(\mu>0\), choose a subsequence on which the compact physical
parameters converge. The lifted reference Hessians converge in a weaker
exponential kernel norm by the cavity locality bound. Proposition~\ref{complete:covariance-prop}
gives continuity of every fixed covariance piece. The initial activity
on a fixed polymer is continuous by its exact finite induced-volume
forest formula. Every fixed number of RG steps uses finitely many local
polymer inputs and Gaussian integrals with the uniform regulator
majorants. Its field jets therefore converge locally. Repeating this
argument as \(\mu\downarrow0\) gives the massless fixed-scale
trajectory. No compactness of an infinite-dimensional activity ball is
being inferred merely from a norm bound.

\subsection{Linear and mixed sources}
The localization and Gaussian identities for sources are unchanged
in physical field coordinates. Before coalescence, the fluctuation
source radius obeys \(\ell_{\sigma,j}\ell_j\asymp x_j\).
For a single-source projection at an arbitrary fixed scale, place an
unused second mark outside that scale's dependence neighborhood.
This keeps its fluctuation radius equal to \(x_jb_j^{s/2}\) and
bounds the extraction error by \(Cf_j/x_j\), independently of any
later coalescence of the actual pair.

Let \(\nu_j\) be the mass in the original, physical polynomial
coordinates, and put
\[
 w_j^{(1)}=\sum_yw_j(y),\qquad \eta'_j=3\Gamma_{j+1}(0),
 \qquad
 d_j=(\nu_j+\eta'_jg_j)w_{j+1}^{(1)}-\nu_jw_j^{(1)}.
\]
The one-source Wick formula and the extraction estimate give
\[
 \lambda_{a,j+1}=(1-d_j)\lambda_{a,j}+e_{a,j},
 \qquad |e_{a,j}|\le Cf_j/x_j.
\]
The inverse normal-form bounds and Lemma~\ref{completion:moments}
give
\[
 b_j^s|\nu_j|\le C\chi_jx_j,
 \qquad
 b_j^s|\nu_{j+1}-\nu_j-\eta'_jg_j|
          \le C(\chi_jx_j^2+f_j),\qquad
 |w_j^{(1)}|\le C\lambda b_j^s.
\]
The leading kinetic term has zero row sum and contributes no linear
source coefficient under constant-field localization. Its higher
terms are covered by the displayed errors.
Telescoping yields
\[
 \sum_{j=k}^{n-1}d_j
 =\nu_nw_n^{(1)}-\nu_kw_k^{(1)}
 +\sum_{j=k}^{n-1}(\nu_j+\eta'_jg_j-\nu_{j+1})w_{j+1}^{(1)}.
\]
The boundary term at infinity vanishes. The remaining sum, the sum
of \(d_j^2\), and \(\sum|e_{a,j}|\) are bounded by
\[
 C\left(x_0/\lambda^2+U/x_0\right)=o(1).
\]
Their tails tend to zero uniformly in the marked site and as the
mass cutoff tends to infinity. Expanding \(\log(1-d_j)\) and using
variation of constants proves that the source coefficients remain
in a fixed positive interval and converge uniformly to positive
limits. This also closes the assumed bounded source domain in the
one-step theorem.

For a pair with coalescence index \(k\), the mixed extraction is zero
on small polymers before \(k\). Afterwards its absolute value is at
most
\[
 \frac{Cf_j}{x_j^2 b_k^s4^{j-k}}.
\]
The Gaussian part is exactly the contraction of the two linear
sources. Thus
\[
 q_\infty(a,c)=\sum_{j\ge0}\widehat\lambda_{a,j}
                 \widehat\lambda_{c,j}\Gamma_{j+1}(a-c)+E(a,c),
 \quad
 |E(a,c)|\le\frac C{b_k^s}
       \left(\chi_k^{3/2}x_k+\frac{Ud_0^k}{x_k^2}\right).
\]
At zero mass this is \(o(b_k^{-s})\). At fixed positive mass the
last torus source integral is
\(O(\ell_{\sigma\sigma,N-1}^{-1})\to0\), since
\(x_{N-1}\) remains positive and \(4^{N-1-k}\to\infty\).
For fixed marks, the remaining tails are dominated uniformly in
the mass cutoff by a constant times
\(4^{-(j-k)}(1+\lambda^2x_0j)^2\), a summable sequence. Therefore
the source sums pass through both limiting operations. The
observable argument of Proposition~\ref{frac:prop:observable} makes
the limiting coefficient independent of the site and gives a
positive multiple of \(|a-c|^{-s}\). The exact convolution identity
\eqref{frac:spin-identity} completes the transfer to spin variables.

\section{Summary of the trajectory and observable estimates}\label{frac:inputs}
The construction in Sections~\ref{complete:covariance}--\ref{completion:flow-closure}
establishes the following estimates for \(d<4\). For \(d=4\), use
Part~\ref{part:fourdim}. All statements are
for a fixed sufficiently large \(R\). They concern the exact
Hubbard--Stratonovich representation from the positive split in
Section~\ref{frac:split} for \(d<4\); for \(d=4\), use the original split. Write its normalized interaction as
\(D^+=P+Q\), put \(p=\sum_xP(x)\), \(q=\sum_xQ(x)\), and
\(B=I-\beta Q\). Here \(\beta\) is the inverse temperature of
\(D^+\), while \(\tau=\beta/(1+\delta_R)\) is the parameter in
\eqref{frac:weight}. Here \(\delta_R=\varepsilon_R\) for \(d<4\),
and \(\delta_R=P_0(0)\) for the original four-dimensional split.
The field normalization is
\(\psi=\sqrt\lambda R^{s/2}u\), with \(h=Bu\) and \(\lambda>0\).
Use \(\mu>0\) for the auxiliary mass; it is a mass coefficient,
not necessarily the square of a parameter.

\paragraph{(F1) Tuned massive trajectories and finite-volume bounds.}
For each sufficiently small \(\mu>0\), on tori of side \(D_N\to\infty\)
there exist exact tuned bulk trajectories with parameters in the physical box \( |\beta-1|\le C_0g\), with \(\beta\), \(\beta p\), and \(1-\beta q\) bounded
above and below by positive constants. Their covariance and local activities have convergent
subsequences at every fixed scale, first as \(N\to\infty\), then as
\(\mu\downarrow0\). The covariance row sum is
\[
 C_0^{\Lambda_N}=\frac{\lambda R^s\beta p}{(1-\beta q)\mu}.
\]
The last-scale bulk integral \(Z_{\rm tail}\) has a positive value at
zero, bounded away from zero, and satisfies
\begin{equation}\label{frac:tail-derivative-input}
 \frac1{|\Lambda_N|}
  \left|D^2_{\boldsymbol1,\boldsymbol1}\log Z_{\rm tail}(0)\right|
     \le C_{R,\mu}D_N^{-s}.
\end{equation}
These are finite-volume estimates on the actual tuned trajectory.
They are not supplied merely by a formal choice of a critical mass.

\paragraph{(F2) Critical covariance and control of observable sectors.}
By Proposition~\ref{complete:covariance-prop} and
Corollary~\ref{complete:covariance-consequences} for \(d<4\), and the covariance estimates of Part~\ref{part:fourdim} for \(d=4\), every massless
subsequential limit has a positive covariance
\(C=\sum_{j\ge0}\Gamma_{j+1}\), with block lengths \(b_j=bL^j\),
finite-range pieces of range at most \(cb_j\), and
\[
 \|\Gamma_{j+1}\|_\infty\le C_R\lambda b_j^{-s},
 \qquad C(x)\sim c_{d,R}|x|^{-s},\qquad c_{d,R}>0.
\]
The graded RG identities give, for all sufficiently distant marks,
\begin{equation}\label{frac:source-input}
 G_\psi(a,c)=\sum_{j\ge0}\widehat\lambda_{a,j}
                 \widehat\lambda_{c,j}\Gamma_{j+1}(a-c)+E(a,c),
 \qquad |E(a,c)|\le b_k^{-s}\eta_k,
\end{equation}
where \(b_k\asymp|a-c|\), \(\eta_k\to0\), and
\(\widehat\lambda_{a,j}\to\lambda_{a,\infty}\) uniformly in the mark.
The source coefficients and their limits are uniformly positive and bounded. The massive
source sums and their errors are uniformly summable in the scale
index for each fixed distant pair as \(\mu\downarrow0\).
The exact final torus source integral tends to zero at fixed positive
mass. Thus the source identities pass first to infinite volume and
then to the massless limit. The source argument in Section~\ref{completion:flow-closure} proves these estimates in addition to the bulk bounds.

\paragraph{How the estimates are supplied.}
The auxiliary-space construction gives the covariance slices and bubble
bounds. The regulated localization and one-step map theorems give
uniform contraction on the whole joint complex neighborhood. The
weighted recursion gives the tuned mass. Section~\ref{completion:flow-closure}
then supplies finite-volume tail estimates and both observable sectors.
In particular the postmass estimates used throughout have the natural
power two, and the local kinetic coordinate has its forward irrelevant
multiplier \(L^{s-2}\).

\section{The critical spin asymptotic}
\label{frac:reduction}
\begin{proof}[Proof of Theorem~\ref{frac:main}]
First, the established estimate (F1) gives the exact susceptibility calculation needed in
Proposition~\ref{frac:endpoint}. Completing the square in a uniform
source before all RG integrations gives
\[
 \chi_\psi^{\Lambda_N}
 =C_0^{\Lambda_N}+(C_0^{\Lambda_N})^2|\Lambda_N|^{-1}
 D^2_{\boldsymbol1,\boldsymbol1}\log Z_{\rm tail}(0).
\]
Conditionally on the spins, \(h\) is Gaussian with mean \(\beta P\sigma\)
and covariance \(\beta P\). Since \(\psi=\sqrt\lambda R^{s/2}B^{-1}h\),
its row-sum covariance obeys the exact identity
\[
 \chi_\psi^{\Lambda_N}
 =\lambda R^s\frac{\beta p+(\beta p)^2\chi_\sigma^{\Lambda_N}}
                         {(1-\beta q)^2}.
\]
Together with \eqref{frac:tail-derivative-input}, these equations imply
\begin{equation}\label{frac:suscept-identity}
 \chi_\sigma^{\Lambda_N}(\beta_{\mu,N})
   =\frac{1-\beta_{\mu,N}q}{\beta_{\mu,N}p\,\mu}
       -\frac1{\beta_{\mu,N}p}+O_{R,\mu}(D_N^{-s}).
\end{equation}
Proposition~\ref{frac:endpoint} applies to the summable ferromagnetic
interaction \(D^+\), whose diagonal is irrelevant. Hence every massive
limiting temperature is strictly subcritical, and every massless
accumulation point is its critical temperature \(\beta_c^+\). In the
original parameter this is \(\tau_c=\beta_c^+/(1+\delta_R)\).
Critical free correlations are the limits of these massive
subcritical correlations, without requiring monotonic tuning.

The conditional Gaussian representation and summability of its
kernels identify the massless auxiliary covariance with the one
constructed from the critical free spin state. It is therefore
translation invariant and bounded at all pairs, including the near
pairs not covered by the source flow. Indeed spin correlations are
bounded by one, while the conditional Gaussian covariance and its mean kernels
are summable at every fixed \(R\). Next, (F2) and the source argument
of Section~\ref{frac:sec:observables} give
\begin{equation}\label{frac:aux-asymptotic}
 G_\psi(a,c)=\lambda_\infty^2c_{d,R}|a-c|^{-s}
                    +o(|a-c|^{-s}),\qquad \lambda_\infty>0.
\end{equation}
The exact auxiliary-to-spin identity is
\begin{equation}\label{frac:spin-identity}
 G_\sigma=(\lambda R^s)^{-1}P_* *G_\psi*P_*-(\beta P)^{-1},
 \qquad P_*=B(\beta P)^{-1}.
\end{equation}
The kernels \(P_*\) and \((\beta P)^{-1}\) satisfy
\(O_R((1+|x|)^{-d-s})\) bounds for \(d<4\), as proved in
Section~\ref{frac:inverse-kernels}. Their tails are summable; exponential
summability is not claimed. Moreover
\[
 \sum_xP_*(x)=\frac{1-\beta_c^+q}{\beta_c^+p}>0.
\]
For completeness, if \(|H(y)|\le C(1+|y|)^{-d-s}\) and a bounded
function satisfies \(F(x)\sim A|x|^{-s}\), \(0<s<d\), then
\begin{equation}\label{frac:convolution}
 (H*F)(x)\sim A\Big(\sum_yH(y)\Big)|x|^{-s}
\end{equation}
whenever the displayed coefficient is nonzero, and in all cases
the difference from that leading expression is \(o(|x|^{-s})\).
To prove this, first sum over \(|y|\le M\), then over
\(M<|y|\le|x|/2\); the latter is bounded by
\(C|x|^{-s}\sum_{|y|>M}|H(y)|\). On
\(|y|>|x|/2\), \(|x-y|\le |x|/2\), use
\(|H(y)|\le C|x|^{-d-s}\) and
\(\sum_{|z|\le |x|/2}|F(z)|\le C|x|^{d-s}\), giving
\(O(|x|^{-2s})\). The remaining region is bounded by the same
order using the tail of \(H\) and \(|F(x-y)|\le C|x|^{-s}\).
Take \(|x|\to\infty\), then \(M\to\infty\). The argument also
preserves the global \(O((1+|x|)^{-s})\) bound, so it can be used twice.

Applying \eqref{frac:convolution} in \eqref{frac:spin-identity}, and
observing that the final inverse kernel is \(o(|x|^{-s})\), yields
\[
 C_{d,R}=\frac{\lambda_\infty^2 c_{d,R}}{\lambda R^s}
       \left(\frac{1-\beta_c^+q}{\beta_c^+p}\right)^2>0.
\]
For \(d=4\) the same convolution step uses the exponential inverse
kernels of Part~\ref{part:fourdim}. This proves Theorem~\ref{frac:main}, using the estimates established above.

\end{proof}

\part{The four-dimensional specialization}\label{part:fourdim}
\section*{Four-dimensional model and theorem}
The spread parameter is denoted by \(R=\rho\) below,
with \(R\) an integer; for a real parameter take \(R=\lfloor\rho\rfloor\).
Specifically,
\[
 D_R(x)=\frac{\boldsymbol1_{\{0<|x|_\infty\le R\}}}{(2R+1)^4-1},
 \qquad (\Delta_R f)(x)=\sum_yD_R(y)[f(x+y)-f(x)].
\]
The spin weight is \(\exp(\beta^{\rm box}\sum_{x,y}D_R(x-y)\sigma_x\sigma_y/2)\).
Let \(\beta_c^{\rm box}=\beta_c(D_R)\) be the usual
spontaneous-magnetization critical inverse temperature. We write
\(\langle\cdot\rangle^f_\beta\) for the infinite-volume free state,
obtained as the increasing-volume limit with free boundary conditions,
and use the Euclidean norm in the asymptotic below.

\begin{theorem}\label{thm:main}
For the four-dimensional Ising model with the normalized
box interaction \(D_R\), and all sufficiently large fixed
\(R\), its infinite-volume critical free two-point function obeys
\[
 \langle\sigma_0\sigma_x\rangle^f_{\beta_c^{\rm box}}
       \sim C_R |x|^{-2}
       \qquad (|x|\longrightarrow\infty)
\]
with \(C_R>0\).
\end{theorem}

\paragraph{Proof conventions.}
The conclusion concerns each sufficiently large \emph{fixed} \(R\).
The constant in the asymptotic may depend on \(R\). A different small
parameter \(\varrho\in(0,1)\) is used for polymer-size weights and
must not be confused with the spread radius.

The proof uses the finite-range RG framework of Bauerschmidt, Brydges
and Slade, with the decay-storing regulator of \cite{ParkMap,ParkTorus}.
The local algebra and norm estimates of that framework are used as
inputs; the covariance, initial activity, enlarged-domain argument,
and physical parameter matching for the box model are constructed
below. A componentwise circle-product formula specifies the last
change of RG coordinates exactly. The first two sections of this part also record
the sextic-prefactor estimates; the eventual initialization keeps the
whole sextic remainder in the activity and controls its transient.
The physical small parameter is \(g=R^{-4}\); a running quartic
coupling is denoted by \(G\) or \(x_j\) in the flow sections.

\section{Local initialization with a decaying regulator}

Put
\[
U(t)=\tfrac12t^2-\log\cosh t,\qquad
U_r(t)=\tfrac r2t^2+U(t),\qquad
p_r(t)=\tfrac r2t^2+\tfrac1{12}t^4.
\]
Then \(U''(t)=\tanh^2t\), and
\begin{equation}\label{eq:utaylor}
 U(t)\ge c\min(t^4,t^2),\qquad
 U(t)=\tfrac1{12}t^4-\tfrac1{45}t^6+O(t^8).
\end{equation}
Consider a box \(B\) of side \(R\) and cardinality \(n=R^4\). In microscopic field units let
\[
 h=R^{-1},\qquad \ell=R^{-2},\qquad g=R^{-4},
 \qquad |r|\le C_0g.
\]
This is a local estimate under a mesoscopic small-field normalization. The application to the box interaction requires obtaining such a normalization for the exact weight.

For a function of the fields on a finite set, use the absolute Taylor majorant
\[
 |f|_{a,\phi}=\sum_{\alpha}
 \frac{a^{|\alpha|}}{\alpha!}\,|\partial^\alpha f(\phi)|.
\]
It is submultiplicative and bounds the field part of Park's \(T_a(\phi)\) norm, since the test tensors are bounded by \(a^k\) at degree \(k\). We use this stronger majorant in the local estimates. Let
\[
 H_B(\phi)=\exp\{-\kappa\norm{\phi_B}_2^2/(nh^2)\}.
\]
The other regulators used below have the following properties. The growing regulator \(G_B\) dominates \(\exp(c\norm{\phi_B}_\infty^2/\ell^2)\), and \(\widetilde G_B\) dominates
\[
 \exp\left(c\ell^{-2}\inf_{p\text{ affine}}
                 \norm{\phi_B-p}_\infty^2\right)
\]
for fixed \(c>0\). These follow from the block test-function norms in \cite[\S 2.4]{ParkMap}. Regulators on a bounded union of blocks are products over the blocks.

\begin{proposition}[Local sextic subtraction]\label{prop:local}
Set
\[
 P_B=\sum_{y\in B}p_r(\phi_y),\qquad
 S_B=\tfrac1{45}\sum_{y\in B}\phi_y^6,\qquad
 E_B=e^{-\sum_{y\in B}U_r(\phi_y)}-e^{-P_B}(1+S_B).
\]
For a fixed sufficiently small \(\kappa>0\), uniformly in \(R\) sufficiently large,
\begin{equation}\label{eq:norms}
 \sup_\phi\frac{|E_B|_{\ell,\phi}}{G_B(\phi)}
     \le Cg^3,
 \qquad
 \sup_\phi\frac{|E_B|_{h,\phi}}{H_B(\phi)\widetilde G_B(\phi)}
     \le Cg.
\end{equation}
The first bound fails at order \(g^3\) for the direct subtraction with \(S_B=0\). The latter subtraction also leaves an error of order \(g^{1/2}\) at \(\phi=h\mathbf 1\).
\end{proposition}

\begin{proof}
Write \(M(t)=\min(t^4,t^2)\). For \(|t|\le1\), the first three derivatives of \(U\) are bounded by \(C|t|^3,C|t|^2,C|t|\). The sum of the terms of degree at least four in its absolute Taylor majorant is \(O(h^4)\), by uniform analyticity of \(U''\) in a fixed strip. For \(|t|\ge1\) the first derivative is \(O(|t|)\), and the remaining Taylor terms sum to \(O(h^2)\). Young's inequality and \(|r|\le C_0h^4\) give, for \(0<a\le h\),
\[
 |e^{-U_r}|_{a,t}\le e^{Ch^4-cM(t)},\qquad
 |e^{-p_r}|_{a,t}\le e^{Ch^4-ct^4}.
\]
The same bound with \(M\) holds for the exponential of the negative interpolated potential \((1-s)p_r+sU_r\), \(0\le s\le1\). These bounds follow by taking the exponential of the absolute nonconstant Taylor terms and retaining the exact constant term in the exponent. They give \(C\exp(-c\sum_B M(\phi_y))\) on the block.

Let \(\Delta=\sum_B U_r-P_B\). The identity
\[
 E_B=-e^{-P_B}(\Delta+S_B)
       +\Delta^2\int_0^1(1-s)e^{-P_B-s\Delta}\,ds
\]
is exact. On the single-site level, the absolute Taylor majorants of \(U_r-p_r\) and \(U_r-p_r+t^6/45\) are bounded by \(C(|t|+a)^6\) and \(C(|t|+a)^8\), respectively. Near zero this follows from the convergent expansion, and away from zero it follows from the uniform strip estimates and polynomial growth. With \(M_\phi=\norm{\phi_B}_\infty\), we obtain
\begin{equation}\label{eq:localmajorant}
 |E_B|_{a,\phi}\le C\big[n(M_\phi+a)^8+n^2(M_\phi+a)^{12}\big]
                              e^{-c\sum_B M(\phi_y)}.
\end{equation}
For \(a=\ell\), division by \(G_B\) absorbs the polynomial and gives \(O(n\ell^8+n^2\ell^{12})=O(g^3)\). For \(a=h\), let \(p\) be an affine approximation on the block. On affine fields the supremum is bounded uniformly by the normalized block \(L^2\) norm. Consequently
\[
 \frac{M_\phi}{h}
 \le C\left(\frac{\norm{\phi_B}_2}{\sqrt n h}
      +\frac{\norm{\phi_B-p}_\infty}{h}\right).
\]
Use \(M(t)\ge c'h^2t^2-C'h^4\), \(nh^4=1\), and \(\ell\le h\). A fraction of the exponential in \eqref{eq:localmajorant} controls the first term in parentheses after division by \(H_B\), and \(\widetilde G_B\) controls the second. The result is \(O(nh^8+n^2h^{12})=O(g)\). For the lower bound on direct subtraction, the sixth Taylor coefficient tested along the constant direction has size \(n\ell^6/45=g^2/45\). Evaluating the direct subtraction at \(h\mathbf 1\) gives \((e^{-1/12}/45+o(1))nh^6\).
\end{proof}

The Gaussian part of the regulator propagates exactly. If \(\zeta_B\) has covariance \(\Gamma_B\) and \(b=\kappa/(nh^2)\), completing the square gives
\begin{equation}\label{eq:gaussH}
 \E e^{-b\norm{\phi_B+\zeta_B}_2^2}
 =
 \det(I+2b\Gamma_B)^{-1/2}
 \exp\{-b\langle\phi_B,(I+2b\Gamma_B)^{-1}\phi_B\rangle\}.
\end{equation}
In particular it is at most \(H_B(\phi)^{1/2}\) if \(2b\norm{\Gamma_B}_{\mathrm{op}}\le1\). This is the mechanism of \cite[Appendix C]{ParkMap}. In the initial comparison, Park's large-field weight \(g^{9/4}\) makes the second bound in \eqref{eq:norms} at most order \(g^3\). The direct comparison instead has a weighted contribution of order \(g^{11/4}\).

\section{Gaussian transport with quartic and source factors}

Here we use the canonical field at scale \(j\). Let \(X\) contain at most \(M\) blocks of side \(L^j\), with \(M\) fixed. Set
\[
 n=L^{4j},\qquad \ell=L^{-j},\qquad h=g^{-1/4}\ell.
\]
For the local statements below assume
\begin{equation}\label{eq:covass}
 \max_{x,y\in X}|\Gamma_{xy}|\le C_\Gamma\ell^2,\qquad
 \E(1+\norm{\zeta_X}_\infty/\ell)^k\le C_{k,M}
 \quad(k\ge1).
\end{equation}
The second condition holds for the usual smooth finite-range slices, by a block Sobolev estimate and Gaussian moments. It is included explicitly among the hypotheses here.

Use commuting nilpotent variables \(s_a,s_b\) with \(s_a^2=s_b^2=0\). The absolute Taylor majorant includes the sum of the four source sectors, weighted by \(1,\tau_a,\tau_b,\tau_a\tau_b\). We assume \(\tau_\#\ell\lesssim g\) for the fluctuation norm and \(\tau_\# h\lesssim1\) for the large-field norm. For a bulk sextic \(S_X=u\sum_{x\in X}\phi_x^6\) assume
\begin{equation}\label{eq:sizesix}
 n|u|\ell^6\le Cg^2 d_j,\qquad 0\le d_j\le1.
\end{equation}

\begin{lemma}[Linear-source transport]\label{lem:wick}
Assume \(a,b\in X\) and let \(J_s(\phi)=\lambda_a s_a\phi_a+\lambda_b s_b\phi_b\), with bounded \(\lambda_a,\lambda_b\). Define
\[
 \mathcal H_{6,v}(t)=t^6+15vt^4+45v^2t^2+15v^3 .
\]
With \(v_x=\Gamma_{xx}\) and \(b_x=\lambda_a s_a\Gamma_{xa}+\lambda_b s_b\Gamma_{xb}\), one has
\begin{equation}\label{eq:wickshift}
 \E\left[e^{J_s(\phi+\zeta)} S_X(\phi+\zeta)\right]
 = e^{J_s(\phi)+\lambda_a\lambda_b s_as_b\Gamma_{ab}}\,
       u\sum_{x\in X}\mathcal H_{6,v_x}(\phi_x+b_x).
\end{equation}
After dividing by the displayed exponential, the two single-source coefficients are given by
\[
 6u\lambda_\#\sum_{x\in X}\Gamma_{x\#}
             (\phi_x^5+10v_x\phi_x^3+15v_x^2\phi_x),
\]
and the double-source coefficient is
\[
 30u\lambda_a\lambda_b\sum_{x\in X}\Gamma_{xa}\Gamma_{xb}
             (\phi_x^4+6v_x\phi_x^2+3v_x^2).
\]
Their weighted Taylor norms at zero are, respectively,
\[
 \begin{array}{c|cc}
  & \ell\text{ scale} & h\text{ scale}\\ \hline
  \text{one source}&O_M(g^3d_j)&O_M(gd_j)\\
  \text{two sources}&O_M(g^4d_j)&O_M(g^{3/2}d_j).
 \end{array}
\]
For a finite-range slice of range \(r_\Gamma\), the last polynomial vanishes if \(|a-b|>2r_\Gamma\).
\end{lemma}

\begin{proof}
Complete the square in the Gaussian expectation with a linear source and then take the coefficient in each sector. Equivalently, apply the identity with real source parameters and take the Taylor polynomial with each source degree at most one. Each source replaces one field by \(\lambda_\#\Gamma_{x\#}\) and carries an additional factor \(\tau_\#\). Thus one source costs at most \(C\tau_\#\ell^2/\mathfrak h\) relative to the bulk size at field radius \(\mathfrak h\). At \(\ell\) this is \(O(g)\) and at \(h\) it is \(O(g^{1/2})\). The bulk sizes at the two radii are \(O_M(g^2d_j)\) and \(O_M(g^{1/2}d_j)\), by \eqref{eq:sizesix}. The support statement follows from the two covariance factors.
\end{proof}

\begin{proposition}[Local mixed quartic and source bound]\label{prop:mixed}
Assume \eqref{eq:covass} and \eqref{eq:sizesix}. Let
\[
 V_X=\sum_{x\in X}(g\phi_x^4/4+\nu\phi_x^2/2),\qquad
 n|\nu|\ell^2\le Cg,\qquad I_X=e^{-V_X+J_s}.
\]
Put \(\mathcal E_X=\Cov_\Gamma(S_X(\phi+\zeta), I_X(\phi+\zeta))\).
In the full source-weighted Taylor majorant,
\[
 \sup_\phi \frac{|\mathcal E_X|_{\ell,\phi}}{G_X(\phi)}
 \le C_M g^3d_j,\qquad
 \sup_\phi \frac{|\mathcal E_X|_{h,\phi}}{H_X(\phi)\widetilde G_X(\phi)}
 \le C_M gd_j.
\]
These estimates bound the double-source sector as well as the two single-source sectors.
\end{proposition}

\begin{proof}
Take two centered Gaussian vectors \((\zeta_t,\zeta'_t)\), each with marginal covariance \(\Gamma\) and with cross covariance \(t\Gamma\). Gaussian differentiation of the cross covariance, followed by integration in \(t\), gives
\begin{equation}\label{eq:interp}
 \mathcal E_X
 =\int_0^1\sum_{x,y\in X}\Gamma_{xy}
      \E\big[\partial_x S_X(\phi+\zeta_t)
                   \partial_y I_X(\phi+\zeta'_t)\big]\,dt .
\end{equation}
The identity holds coefficient by coefficient in the source algebra. For \(\mathfrak h\in\{\ell,h\}\) write \(p(\psi)=1+\norm{\psi_X}_\infty/\mathfrak h\). Submultiplicativity, the quartic Taylor estimate in the proof of Proposition~\ref{prop:local}, and the source bounds give
\begin{align*}
 \sum_{x\in X}|\partial_xS_X|_{\mathfrak h,\psi}
 &\le C_M n|u|\mathfrak h^5 p(\psi)^5,\\
 \sum_{y\in X}|\partial_y I_X|_{\mathfrak h,\psi}
 &\le C_M\big(n g\mathfrak h^3+n|\nu|\mathfrak h+\tau_a+\tau_b\big)
            p(\psi)^5 e^{-c g\sum_X\psi_y^4}.
\end{align*}
Sources outside \(X\) are omitted. In the second line the factors from \(e^{J_s}\) and its field derivatives have been bounded as polynomials using nilpotence. At \(\ell\) the prefactors in the two lines are bounded by \(C_Mg^2d_j/\ell\) and \(C_Mg/\ell\). Using \eqref{eq:covass} in \eqref{eq:interp}, the expectation leaves only a polynomial in \(1+\norm{\phi_X}_\infty/\ell\). This is absorbed by \(G_X\).

At \(h\) the corresponding prefactors are \(C_Mg^{1/2}d_j/h\) and \(C_M/h\). The covariance factor \(C_\Gamma\ell^2\) then gives \(C_Mgd_j\). To keep the decay in the expectation, for any fixed \(A>0\) bound the quartic factor by
\[
 C_{A,M}\exp\{-A\norm{\phi_X+\zeta'_t}_2^2/(nh^2)\}.
\]
This follows by completing a square at each site, since \(ngh^4=1\). Cauchy--Schwarz separates the polynomial fluctuation factor from this Gaussian factor. Apply \eqref{eq:gaussH} to the latter, using
\(\norm{\Gamma_X}_{\mathrm{op}}\lesssim_M n\ell^2\) and
\(\ell^2/h^2=g^{1/2}\). By taking \(A\) a sufficiently large fixed multiple of \(\kappa\), there is enough decay left to divide by \(H_X\) and absorb powers of the block \(L^2\) averages. The remaining powers of the deviations from affine fields are absorbed by \(\widetilde G_X\), as in Proposition~\ref{prop:local}.
\end{proof}

For the single Gaussian integration there is therefore an exact comparison
\begin{equation}\label{eq:comparison}
 \E_\Gamma[(1+S_X)I_X]=(1+\overline S_X)\E_\Gamma I_X+\mathcal E_X,
 \qquad \overline S_X=\E_\Gamma S_X(\phi+\zeta).
\end{equation}
The terms of degrees zero, two and four in \(\overline S_X-S_X\) are explicit from \(\mathcal H_{6,v}\). They are the required bulk counterterms at first order in the sextic. If \(u\) is held fixed in microscopic units, then \(n\ell^6\) decreases by \(L^{-2}\) per step. In a stable quartic reference factor, replacing \(1+S_X+\overline S_X-S_X\) by \((1+S_X)e^{\overline S_X-S_X}\) produces terms quadratic in the sextic, since \(e^A-1-A=A^2\int_0^1(1-t)e^{tA}\,dt\). The quartic part of \(A=\overline S_X-S_X\) is smaller than the quartic in \(V_X\) by \(O_M(gd_j)\). The stable interpolation bound thus applies after multiplication by the quartic reference. For nondegenerate \(\Gamma\), restricting \(\phi+\zeta\) to a fixed bounded set gives a Gaussian lower bound on the bulk convolution. Thus a positive quartic counterterm cannot be exponentiated directly against the fully averaged weight.

\section{Exact preprocessing for the standard box}

Let \(D_R\) be uniform on \(\{x\in\Z^4:0<|x|_\infty\le R\}\). On a large periodic torus the zero-field weight is \(\exp(\beta\langle\sigma,D_R\sigma\rangle/2)\). Let \(f_R\) be a nonnegative, even, hypercubic-symmetric smooth bump sampled on \(\Z^4\), supported on \(|x|_\infty\le R/3\), and normalized to sum to one. Take
\[
 P_0=\theta(f_R*f_R)+e^{-R}\delta_0
\]
with a sufficiently small fixed \(\theta>0\). Then \(P_0(x)\le D_R(x)\) for \(x\ne0\). This follows from \(\norm{f_R*f_R}_\infty\le\norm{f_R}_\infty\lesssim R^{-4}\). Replace \(D_R\) by
\[
 D=\frac{D_R+P_0(0)\delta_0}{1+P_0(0)},\qquad
 P=\frac{P_0}{1+P_0(0)},\qquad Q=D-P .
\]
This adds a spin-independent diagonal energy and reparametrizes inverse temperature. Write \(\delta_R=P_0(0)\). The two temperature parameters satisfy
\(\beta=(1+\delta_R)\beta^{\rm box}\), and consequently
\[
 \beta_c^D=(1+\delta_R)\beta_c^{\rm box}.
\]
Below \(\beta\) always denotes the inverse temperature for \(D\);
the endpoint identified by the tuning argument is \(\beta_c^D\). Write \(p=\sum P\), \(q=\sum Q=1-p\). We have \(p\asymp1\), \(Q(0)=0\), \(Q\ge0\), \(\norm{Q}_\infty\lesssim R^{-4}\), and \(\widehat P(k)>0\). For \(\beta\) in a fixed sufficiently small neighborhood of one, \(J=\beta Q\) obeys \(\norm{J}_1\le q_*<1\).

Let \(A=\beta P\), and let
\[
 F(h)=\log \E_{J,0}e^{\langle h,\sigma\rangle}
\]
be the pressure increment of the reference Ising system with interaction \(J\). The Hubbard--Stratonovich transformation applied to \(A\) alone gives the exact density for \(h\)
\begin{equation}\label{eq:HS}
 \frac1Z\exp\{-\tfrac12\langle h,A^{-1}h\rangle+F(h)\}.
\end{equation}
Thus the reference spins have been integrated out. For a parameter \(m^2>0\) define
\[
 S_0=(I-J)^{-1},\qquad c=\frac{1-m^2-\beta q}{\beta p},\qquad
 M_m=I-J-cA .
\]
For \(m^2\) small and \(\beta\) close to one, \(c>0\) and is bounded above and below uniformly. The Gaussian covariance and residual interaction in \eqref{eq:HS} can be taken as
\begin{equation}\label{eq:preCov}
 C_m=(A^{-1}-cS_0)^{-1}=A(I-J)M_m^{-1},\qquad
 W(h)=\tfrac c2\langle h,S_0h\rangle-F(h).
\end{equation}
All identities are exact. In particular the mass parameter here specifies a Gaussian decomposition, rather than an approximation to the spin measure.

\begin{proposition}[Finite-range slices]\label{prop:frd}
On \(\Z^4\), the covariance \(C_m\) admits positive finite-range pieces at lengths \(O(RL^j)\), \(j\ge0\). For any fixed orders \(k,N\ge0\) their kernels obey
\begin{equation}\label{eq:slicebound}
 \norm{\nabla^k C_{m,j}}_\infty
 \le C_{k,N,L}R^{-4-k}L^{-(2+k)j}(1+m^2 L^{2j})^{-N}.
\end{equation}
The bound includes the first piece, with \(j=0\). In particular the massless auxiliary fluctuation has site variance \(O(R^{-4})\).
\end{proposition}

\begin{proof}
The multiplier of \(M_m\) is comparable to
\(m^2+\min(R^2|k|^2,1)\), uniformly in the indicated parameter range. Indeed, \(J+cA\) has mass \(1-m^2\) and its off-diagonal kernel is bounded above and below by constant multiples of the box kernel. For the box kernel the lower bound follows from the product Dirichlet sum for the cube. For \(|k|\le c/R\) it also follows by Taylor expansion and the box second moment. In the annulus \(c/R\le |k|\le C/R\) the rescaled Dirichlet sums converge uniformly to the product of sinc functions, which stays strictly below one there. For \(|k|\ge C/R\), the coordinate of largest absolute value bounds the corresponding normalized Dirichlet sum by \(O(1/C)\). The numerator \(A(I-J)\) is positive semidefinite and has range \(O(R)\). Repeated summation by parts for the sampled smooth bump yields
\[
 |\widehat{A(I-J)}(k)|\le C_s(1+R|k|)^{-s}
\]
for every fixed \(s\), including the diagonal term \(e^{-R}\).

Apply the positive polynomial decomposition of the inverse to \(M_m\), after a fixed scalar rescaling of its spectrum if needed. The polynomials \(P_t\) have degree \(O(t)\), and the integral uses the kernels \(tP_t(M_m)\,dt\). For \(0<t\le1\), \(tP_t\) is a fixed positive constant. For \(t>1\),
\[
 0\le P_t(\lambda)\le C_s(1+t\sqrt\lambda)^{-s}.
\]
See \cite[Lemma 2.7]{Bauerschmidt} and \cite[Lemma 3.8]{BPR}. Multiply the pieces by \(A(I-J)\), which commutes with them and preserves positivity. The range becomes \(O(R(1+t))\). In Fourier inversion, a derivative of order \(k\) costs at most \(C|p|^k\) at frequency \(p\). In the region \(|p|\le R^{-1}\) substitute \(R(1+t)p\), and in the complementary region use the decay of the numerator and of \(P_t\). For \(t>1\) this bounds the differentiated integrand by
\[
 C_{k,N}R^{-4-k}t(1+t)^{-4-k}(1+m^2t^2)^{-N}.
\]
For the first interval \([0,1]\) the numerator alone gives \(O(R^{-4-k})\). Integrating over the subsequent intervals \([L^{j-1},L^j]\) proves \eqref{eq:slicebound}.
\end{proof}


The same construction permits a small kinetic counterterm. Write
\(B=A(I-J)\). For \(\zeta\) in a sufficiently small fixed interval, replace
\(M_m\) by \(M_m+\zeta B(-\Delta_{\rm nn})\). The new covariance has
precision \(C_m^{-1}+\zeta(-\Delta_{\rm nn})\), and its denominator is
still finite range and obeys the same symbol comparison. Indeed,
\(\widehat B\) is uniformly bounded, and the multiplier of
\(-\Delta_{\rm nn}\) is at most a constant times \(\min(R^2|p|^2,1)\)
on the Brillouin zone. Multiplication by \(B\) still preserves positivity
of the polynomial slices.

\section{Locality and stability of the exact reference pressure}

Write \(\epsilon=R^{-4}\). For a real field \(h\), let \(m_i^{\rm ex}(h)\)
be the magnetization of the reference system. Let \(z\) be the unique real
solution to
\begin{equation}\label{eq:MF}
 z_i=h_i+\sum_k J_{ik}\tanh z_k,\qquad m_i^{\rm mf}=\tanh z_i .
\end{equation}
The row sum bound on \(J\) gives uniqueness by contraction.

\begin{proposition}[Uniform reference comparison]\label{prop:cavity}
There are constants \(r,\eta>0\), independent of \(R\), of the torus
volume, and of the real background \(h\), with the following property.
Expand \(m_i^{\rm ex}(h+x)-m_i^{\rm mf}(h+x)\) in \(x\) with coefficients
\(b_{i,\alpha}\). If \(\tau(i,S)\) denotes the shortest tree length for the periodic \(\ell^\infty\)
metric joining \(i\) and \(S\), allowing intermediate vertices,
then
\begin{equation}\label{eq:weighted}
 \sum_{\alpha}|b_{i,\alpha}|r^{|\alpha|}
       \exp\{\eta\tau(i,\operatorname{supp}\alpha)/R\}
 \le C\epsilon.
\end{equation}
The two magnetizations separately have uniformly bounded norms of this
type. In particular, the pressure defect has uniformly small derivative
rows, with an exponential tree weight, at every real background.
\end{proposition}

\begin{proof}
For series at anchor \(i\), use the norm on the left of
\eqref{eq:weighted}. The union of two trees rooted at \(i\) gives
submultiplicativity. Changing the anchor from \(k\) to \(i\) costs at
most \(e^{\eta|i-k|/R}\), since one may add an edge between the two
anchors. The norm of \(x_i\) is \(r\).

We give the recursion used for the exact magnetization. On a finite
graph with real fields, the half log ratio of the weights conditioned
on \(\sigma_i=1\) and \(\sigma_i=-1\) is \(h_i\) plus half the log of the
ratio of the two partition functions on the graph with \(i\) deleted.
In this ratio the fields at the neighbors are shifted by \(+J_{ik}\)
and \(-J_{ik}\). Order the neighbors and telescope by changing one shift
at a time. For a change at \(k\), the ratio on the smaller graph is
\(\cosh(H_k+J_{ik})/\cosh(H_k-J_{ik})\), where \(H_k\) is the cavity
half log ratio with the other shifts already chosen. Thus, at a node
of the recursion with \(d\) ancestors removed,
\begin{equation}\label{eq:cavity}
 H_i=h_i+x_i+\xi_i+
       \sum_{k\text{ available}} f_{J_{ik}}(H_k^{\rm child}),
 \qquad
 f_b(z)=\tfrac12\log\frac{\cosh(z+b)}{\cosh(z-b)}.
\end{equation}
Here \(|\xi_i|\le C d\epsilon\). The deleted vertices remove row mass at
most \(C(d+1)\epsilon\). The recursion ends after finitely many
deletions. At the root its magnetization is \(\tanh H_i\).

For real \(z\), \(|f_b'(z)|\le\tanh b\le b\). On a fixed sufficiently
thin strip, the higher Taylor coefficients around any real point
are bounded by \(C b\) times a geometric sequence. This follows from
the strip formula for \(f_b'\), and Cauchy's estimate on a smaller
strip. The same argument applied to
\(f_b(z)-b\tanh z\) gives a bound \(C b^3\) on its Taylor norm at a
small fixed radius. These statements also hold after substituting
a series whose nonconstant part is sufficiently small in an anchored
norm, by submultiplicativity.

Choose \(\eta>0\) small enough that \(e^\eta q_*<1\). Then choose
\(r>0\) small. For a sufficiently large fixed \(K\) and all series
with nonconstant norm at most \(Kr\), the Lipschitz bound for the
composition with \(f_b\) is \(b(1+CKr)\) in the anchored norm.
To see this, integrate \(f_b'\) along the segment between the two
series. The constant part along that segment is real, and the
nonconstant part has norm at most \(Kr\).
The same statement holds for the composition with \(b\tanh\).
Choose \(r\) so that \(q'=e^\eta q_*(1+CKr)<1\) and
\(r+q'Kr\le Kr\). Backward induction in \eqref{eq:cavity} proves
\(\norm{H_i-H_i(0)}_i\le Kr\). Iterating \eqref{eq:MF} about its real
solution in the same Banach algebra gives the corresponding analytic
series for the mean-field solution.

Compare a node of \eqref{eq:cavity} with the full mean-field solution
at its anchor. The shifts and omitted terms cost
\(C(d+1)\epsilon\). The cubic approximation to \(f_b\) costs at most
\(C\sum_k J_{ik}^3\le C\epsilon^2\). Differences at child nodes enter
with weights whose sum is at most \(q'\), including the anchor change.
Expanding this inequality down the finite recursion yields at the
root a bound
\[
 C\epsilon\sum_{d\ge0}(d+1)(q')^d\le C'\epsilon.
\]
The composition with \(\tanh\) has uniformly bounded Lipschitz constant
on these series. This proves \eqref{eq:weighted}.
\end{proof}

For clarity, let \(I_{\rm bin}(t)\) be the binary relative entropy with
respect to the fair spin, continuously defined for \(|t|\le1\).
The mean-field pressure increment is
\[
 F_{\rm mf}(h)=
 \sup_{m\in[-1,1]^\Lambda}
 \left\{\langle h,m\rangle+\tfrac12\langle m,Jm\rangle
                  -\sum_x I_{\rm bin}(m_x)\right\}.
\]
Its Hessian as a function of the maximizing variable in the open
cube is negative definite, since \(I_{\rm bin}''\ge1\).
The maximizer is the magnetization in \eqref{eq:MF}, and
\(\nabla F_{\rm mf}=m^{\rm mf}\). The exact pressure satisfies
\(\nabla F=m^{\rm ex}\). Both pressures vanish at zero and are even.
Applying Proposition~\ref{prop:cavity} at zero to their gradients
and integrating along the field proves that \(F-F_{\rm mf}\) has an
exact anchored Taylor expansion on \(\norm h_\infty<r\), with
degree-\(k\) kernels bounded by \(C\epsilon r^{-k}\) in the weighted
row norm (after reducing \(r\) by a fixed factor if necessary).
In particular, if \(S=D^2F(0)\),
\begin{equation}\label{eq:sus}
 \sup_i\sum_k |(S-S_0)_{ik}|e^{\eta|i-k|/R}
 \le C\epsilon .
\end{equation}

\begin{corollary}[Coercivity after Hessian subtraction]\label{cor:stable}
For all real \(h\), uniformly in the volume and in \(R\) large enough,
\begin{align}
 \tfrac12\langle h,Sh\rangle-F(h)
   &\ge c_1\sum_i\frac{h_i^4}{1+h_i^2},\\
 W(h)&\ge c_1\sum_i\frac{h_i^4}{1+h_i^2}
       - C(\epsilon+|\beta-1|+m^2)\norm h_2^2 .
       \label{eq:stabilityW}
\end{align}
\end{corollary}

\begin{proof}
Put \(u=S_0h\) and \(V_{\rm ent}(t)=I_{\rm bin}(t)-t^2/2\). Completing
the square in the variational formula gives
\[
 \tfrac12\langle h,S_0h\rangle-F_{\rm mf}(h)
 =\inf_{m\in[-1,1]^\Lambda}
  \left\{\tfrac12\langle m-u,(I-J)(m-u)\rangle
                     +\sum_i V_{\rm ent}(m_i)\right\}.
\]
Because \(I-J\ge(1-q_*)I\) and \(V_{\rm ent}(t)\ge t^4/12\),
the right-hand side is bounded below by
\(c\sum_i\min(u_i^4,u_i^2)\). This follows pointwise by considering
whether \(|m_i|\ge |u_i|/2\). For
\(\Psi(t)=t^4/(1+t^2)\), convexity and
\(\Psi(bt)\le b^4\Psi(t)\) for \(b\ge1\) imply
\(\sum_i\Psi(((I-J)u)_i)\le (1+q_*)^4\sum_i\Psi(u_i)\), by
Jensen's inequality and the absolute row and column bounds. Thus
the mean-field deficit is at least \(c'\sum_i\Psi(h_i)\).

Let \(E_4=F-F_{\rm mf}-\frac12\langle h,(S-S_0)h\rangle\).
It vanishes through degree three at zero. The uniform row bounds
from Proposition~\ref{prop:cavity} give
\[
 \sup_h\sup_i\sum_k|\partial_i\partial_kE_4(h)|\le C\epsilon,
 \qquad
 \sup_h\sup_i\sum_{k,l,m}|\partial_i\partial_k\partial_l\partial_m
 E_4(h)|\le C\epsilon.
\]
Split \(h=a+b\) with \(a_i=h_i\) on \(|h_i|\le1\) and zero otherwise.
Fourth-order Taylor expansion, symmetry of the derivative tensor, and
\(|a_i a_k a_l a_m|\le(|a_i|^4+|a_k|^4+|a_l|^4+|a_m|^4)/4\)
give \(|E_4(a)|\le C\epsilon\sum_i|a_i|^4\).
The Hessian bound yields
\(\norm{\nabla E_4(a)}_\infty\le C\epsilon\) and
\[
 |E_4(a+b)-E_4(a)|\le C\epsilon(\norm b_1+\norm b_2^2)
                    \le C'\epsilon\norm b_2^2.
\]
Consequently \(|E_4(h)|\le C\epsilon\sum_i\Psi(h_i)\), which
can be absorbed for large \(R\). Equation~\eqref{eq:sus} and
\(c-1=(1-m^2-\beta)/(\beta p)\) give the last claim.
\end{proof}


\begin{corollary}[Taylor stability of restricted reference weights]\label{cor:stablerestr}
Let \(X\subset\Lambda\), and let \(W_X\) be \(W\) evaluated at fields
supported in \(X\), with derivatives taken only in those fields.
Set \(a=R^{-1}\). If \(|\beta-1|+m^2\lesssim R^{-4}\), then
\[
 |e^{-W_X}|_{a,h}
 \le \exp\{C|X|a^4-c\textstyle\sum_{i\in X}\Psi(h_i)\}.
\]
In particular for an \(R\)-scale block this gives its Gaussian
regulator bound, uniformly in the external volume.
\end{corollary}

\begin{proof}
Write \(U_*(h)=\frac12\langle h,Sh\rangle-F(h)\).
The first three derivatives of \(U_*\) at zero vanish.
The uniform Taylor bounds on the magnetization in
Proposition~\ref{prop:cavity} give bounded absolute derivative rows
at every real field. Split the field into \(h^{<}\), supported on
\(|h_i|\le1\), and \(h^{>}\), supported on \(|h_i|>1\).
For \(k=1,2,3\), integrate the Taylor formula from zero to \(h^{<}\)
using the fourth derivative rows. By symmetry of the derivative
tensor and the arithmetic--geometric mean inequality, the sum
of the absolute entries of the \(k\)-th derivative at \(h^{<}\),
over indices in \(X\), is at most
\(C\sum_{i\in X}|h_i^{<}|^{4-k}\).
Moving from \(h^{<}\) to \(h^{<}+h^{>}\) adds at most
\(C\sum_{i\in X}|h_i^{>}|\), by the uniform \((k+1)\)-st derivative
rows and symmetry. The contribution of orders at least four
to the absolute Taylor majorant is \(O(|X|a^4)\), by the fixed
analytic radius of the magnetization. Consequently the
nonconstant Taylor terms of \(U_*\), at radius \(a\), are bounded
by \(\varepsilon\sum_X\Psi(h_i)+C_\varepsilon|X|a^4\).
Here the small-field terms use Young's inequality, while on
the large-field sites \(|h_i|\le |h_i|^2\).
The quadratic difference between \(W\) and \(U_*\) has absolute
row sum \(O(a^4)\) by \eqref{eq:sus} and the assumed parameter range.
Its value and nonconstant Taylor terms obey the same bound.
Combine the nonconstant estimates with the constant term of
Corollary~\ref{cor:stable} and exponentiate.
\end{proof}

The leading interactions can be computed explicitly. With \(u=S_0h\),
the expansion of the mean-field deficit through degree six is
\begin{equation}\label{eq:entsextic}
 \tfrac12\langle h,S_0h\rangle-F_{\rm mf}(h)
  =\tfrac1{12}\sum_i u_i^4+\tfrac1{30}\sum_i u_i^6
          -\tfrac1{18}\langle u^3,S_0u^3\rangle
          +O(\text{degree }8).
\end{equation}
Here \(u^3\) is pointwise. Indeed, the stationarity equation in the
preceding infimum gives \(m=u-S_0u^3/3+O(\text{degree }5)\), and
\(V_{\rm ent}(m)=m^4/12+m^6/30+O(m^8)\). The second-order correction
from the quadratic term is \(\langle u^3,S_0u^3\rangle/18\), and
the correction from the quartic term is
\(-\langle u^3,S_0u^3\rangle/9\). This proves the formula.
The actual degree-four Taylor kernel differs from its mean-field
value by \(O(R^{-4})\) in the weighted row norm. Its constant-field
coefficient remains positive for large \(R\), and has dimensionless
size of order \(R^{-4}\) on a block of side \(R\) with fluctuation
scale \(R^{-2}\). The degree-six term for this box preparation
already includes a bilocal polynomial. The local log-cosh prefactor
of Proposition~\ref{prop:local} alone would not represent it. The
pressure comparison gives the following block size estimates in
the fluctuation norm before any localization of the Taylor kernels.
For degree \(k\), the exact-minus-mean-field term has size at most
\(C_k R^4R^{-4}(R^{-2})^k\) per anchored block. For \(k=2,4,6\)
these are \(O(g),O(g^2),O(g^3)\), respectively, with
\(g=R^{-4}\). The quadratic kernel and the quartic correction thus
require treatment before the \(O(g^3)\) remainder estimate can be used.

% This section uses the notation in Exact preprocessing and Proposition~\ref{prop:cavity}.
\section{An exact connected expansion of the reference weight}

Partition the torus into blocks of side \(b\), where \(R<b\le M R\)
for a fixed \(M\). Join two distinct blocks when some microscopic
reference edge joins them. This graph has maximum degree \(\Delta\)
bounded independently of the volume and range. For a set \(X\) of
blocks let \(\Lambda_X\) be their union. In the following construction
we restrict the spin system to \(\Lambda_X\) as well as restricting
the field. The reference pressure in a proper subvolume is therefore
computed with free boundary, even when the ambient volume is a torus.

For a real symmetric array \(t=(t_{BC})\in[0,1]^{\binom X2}\), let
\(J_X(t)\) be the restriction of \(J\) to \(\Lambda_X\), with each
cross-block edge between \(B\) and \(C\) multiplied by \(t_{BC}\).
Keep intra-block edges unchanged. Set
\[
 W_X(t,h)=\frac c2\langle h,(I-J_X(t))^{-1}h\rangle
       -\log\frac{Z_{J_X(t)}(h)}{Z_{J_X(t)}(0)}.
\]
Use the same scalar \(c\) as in \eqref{eq:preCov} for every subvolume
and every interpolation point. No auxiliary Gaussian integration
is performed in this construction.

\begin{proposition}[All-field forest expansion]\label{prop:forest}
Assume \(\|J\|_1\le q_*<1\), \(\|J\|_\infty\lesssim R^{-4}\),
\(J\ge0\), and \(|c-1|\lesssim R^{-4}\). Put \(a=R^{-1}\).
There exist connected activities \(A_X\), depending only on the
fields on \(\Lambda_X\), for which
\begin{equation}\label{eq:forest_partition}
 e^{-W_X(\boldsymbol 1,h)}=
       \sum_{\pi\in\mathcal P(X)}\prod_{Y\in\pi} A_Y(h).
\end{equation}
Here \(\mathcal P(X)\) denotes set partitions, and \(A_Y=0\) if
\(Y\) is disconnected in the block graph. For constants independent
of the volume and range, the Taylor seminorm satisfies
\begin{equation}\label{eq:forest_bound}
 |A_X|_{a,h}\le C_0^{|X|}
 \exp\!\left\{C|\Lambda_X|a^4-c_0\sum_{i\in\Lambda_X}\Psi(h_i)\right\}.
\end{equation}
The expansion also holds after multiplying the left side by
\(\exp(\sum_{z\in\{u,v\}\cap\Lambda_X}\sigma_z h_z)\) in the
nilpotent source algebra. The activity on \(Y\) then acquires exactly
the source factor with sites in \(\Lambda_Y\). In a seminorm weighting
\(\sigma_z\) by \(\tau_z\), the right side of \eqref{eq:forest_bound}
is multiplied by at most
\(\prod_{z\in\{u,v\}\cap\Lambda_X}(1+\tau_z(|h_z|+a))\).
\end{proposition}

\begin{proof}
We first extend the stability bound to a complex neighborhood of
all real interpolation points. In this paragraph the notation
\(t+v\) means a perturbation in the variables indexing pairs of
blocks, not a change in the field. There is a fixed \(r>0\) such
that the proof of Proposition~\ref{prop:cavity} works in the joint
Taylor algebra of field increments and block-link increments at
radius \(r\) in each family. For the unweighted bounds needed here
use the coefficient norm \(\sum_{\alpha,\gamma}|a_{\alpha,\gamma}|
 r^{|\alpha|+|\gamma|}\). To check this extension, at a microscopic
edge replace the real parameter \(d=J_{ij}t_{B_iB_j}\) in
\(f_d\) by \(d+J_{ij}v_{B_iB_j}\). At an intra-block edge there is
no increment. The increments in the shifts from deleted neighbors
have total norm at most \(r\sum_j J_{ij}\). They therefore cost
\(O(r)\) uniformly even deep in the recursion. The joint Taylor
expansion of \(f_d\), using its strip representation, bounds
its derivative in the message coordinate by
\(J_{ij}(1+O(r))\) after substitution of any series with
nonconstant norm \(O(r)\). Indeed the constant coefficient has
absolute value at most \(\tanh d\le J_{ij}\), while the remaining
coefficients cost \(O(rJ_{ij})\) by Cauchy's inequality on fixed
strips for the message and the rescaled coupling variable. Also
\(f_d(z)-d\tanh z\) costs \(O(J_{ij}^3)\) in the joint norm, since
it vanishes to order three in \(d\) and is analytic uniformly in
real \(z\). These bounds give the same backward induction with
row contraction \(q_*(1+O(r))<1\). For the mean-field messages the
argument is the same, with \(f_d(z)\) replaced by \(d\tanh z\).
Comparing a cavity node of depth \(k\) to the full mean-field
solution costs \(O((k+1)R^{-4})\) from missing edges and shifts,
uniformly in the joint norm. Summing this error with the
geometric row contraction bounds the exact-minus-mean-field
magnetization by \(O(R^{-4})\) in the joint norm. This argument
applies to every finite induced subgraph. It defines the
complex pressures by integration of the holomorphic
magnetization in the field, with pressure normalized at zero.
This agrees with the real log partition function on the real
parameters and hence with its analytic continuation.

Here are consequences for the pressure stability. At a real
background \(h\), the second and fourth field derivative rows
of \(W_X(t+v,h)-W_X(t,h)\) have absolute sums at most
\(C\varrho\) when \(\|v\|_\infty\le\varrho\le r/2\). This follows
by retaining the terms with positive link degree in the joint
Taylor expansion. At \(h=0\) the second field derivative row
improves to \(C\varrho R^{-4}\), because the exact and mean-field
rows differ by \(O(R^{-4})\) in the joint norm, the latter
Hessian is \((I-J_X(t+v))^{-1}\), and \(|c-1|\lesssim R^{-4}\).
The same bounds hold for the corresponding field derivatives
of the real pressure on the boundary faces \(t_{BC}=0,1\), by
the analytic construction. The pressure difference is even and
zero at zero field. Subtract its Hessian and repeat the split
into coordinates of magnitude at most one and greater than one
from Corollary~\ref{cor:stable}. The fourth derivative controls
the former coordinates and the second derivative controls
increments on the latter. Thus
\[
 |W_X(t+v,h)-W_X(t,h)|
 \le C\varrho\sum_i\Psi(h_i)+C\varrho R^{-4}\|h\|_2^2.
\]
Take \(\varrho\) sufficiently small and use the real coercivity
bound, uniform in \(t\) because the row sum is still at most
\(q_*\). This proves complex coercivity with a fixed positive
fraction of \(\sum_i\Psi(h_i)\) and quadratic loss
\(O(R^{-4})\|h\|_2^2\). The proof of
Corollary~\ref{cor:stablerestr} now applies with the same
absolute derivative row bounds for complex links. Taking the
absolute Taylor majorant in the field of radius \(a\), it gives
\begin{equation}\label{eq:forest_complex}
 |e^{-W_X(t+v,\cdot)}|_{a,h}
 \le \exp\!\left\{C|\Lambda_X|a^4-c_0\sum_i\Psi(h_i)\right\}
 \quad (\|v\|_\infty\le\varrho).
\end{equation}
For example, the quadratic loss is absorbed by
\(R^{-4}s^2\le \varepsilon\Psi(s)+C_\varepsilon a^4\), and the
nonconstant field Taylor terms are estimated by the derivative
rows in the same way as in the real proof. All constants here
are uniform in \(X\) and the real interpolation point.

Apply the Brydges--Kennedy--Abdesselam--Rivasseau forest formula
\cite{AR} on the complete graph with vertex set \(X\) to
\(e^{-W_X(t,h)}\). Given a forest \(\mathcal F\) and interpolation
parameters \(s_e\in[0,1]\) for its edges, its evaluation array has
entry \(\min_{e\text{ on path}}s_e\) for vertices joined by a forest
path and zero for vertices in different forest components.
The formula sums the integral of \(\partial_{\mathcal F}e^{-W_X}\)
at those arrays, where \(\partial_{\mathcal F}=\prod_{e\in\mathcal F}\partial_{t_e}\).
If an edge of \(\mathcal F\) is inactive in the microscopic
interaction, its derivative vanishes. At each evaluation array,
\(J_X(t)\) is a direct sum over the forest components. Its
inverse \((I-J_X(t))^{-1}\) has the same direct-sum property,
and the normalized pressure adds across the components.
This factorization continues to hold upon varying the
coordinates indexed by \(\mathcal F\) with all intercomponent
coordinates fixed to zero. The mixed derivative and the integral
therefore factor over the components. Define \(A_Y\) to be the sum
of the terms corresponding to spanning trees on \(Y\), using
\(W_Y\), and on a singleton use the weight with no link variables.
This proves \eqref{eq:forest_partition} and spatial connectivity.

For a spanning tree on \(Y\), multivariable Cauchy bounds its
mixed derivative in the Taylor seminorm by \(\varrho^{-(|Y|-1)}\)
times \eqref{eq:forest_complex}, reducing \(\varrho\) by a fixed
factor if necessary. One can apply the coefficient formula
before taking the Taylor majorant: integrate each field Taylor
coefficient on the same link polytorus and then sum its absolute
value. There are at most \(\Delta^{|Y|-1}\) spanning trees,
by orienting each tree toward a fixed root and recording one
of at most \(\Delta\) parents at each other vertex. This proves
\eqref{eq:forest_bound}. The field-linear source factor is
independent of the link variables and already factors over
forest components, so the same construction proves the source
statement by the product rule for the field Taylor seminorm.
\end{proof}

To put the partition expansion in circle-product notation, set
\(I(B)=A_{\{B\}}\) and define
\[
 K(Y)=\sum_{\substack{\pi\in\mathcal P(Y)\\|C|\ge2\text{ for every }C\in\pi}}
                \prod_{C\in\pi}A_C,
 \qquad K(\varnothing)=1.
\]
Partitions containing a disconnected part have zero weight.
Thus \(K\) factors over the geometric components of \(Y\),
and \eqref{eq:forest_partition} is exactly \((I\circ K)(X)\).
The number of partitions into connected parts is at most
\(2^{\Delta |Y|/2}\): for every part choose a spanning tree by
a fixed deterministic rule, and encode the partition by the
union of the chosen edge sets. Consequently \eqref{eq:forest_bound}
continues to hold with an enlarged \(C_0\) for this \(K\).
This is an exact all-field connected representation with
Gaussian damping. The bound has an exponential constant in
polymer size and supplies neither the decaying large-set factor
nor the small \(g^3\) norm required by the one-step RG theorem.
The singleton factor here is the full restricted weight, so
matching it to the quartic reference factor is a separate
localization problem.

\begin{proposition}[Size of the raw adjacent-block activity]\label{prop:rawforest}
Use the standard box split and take the minorant small enough that
\(Q_{xy}\ge cR^{-4}\) for \(0<|x-y|_\infty\le R\). Let \(B_1,B_2\)
share a face, write \(X=\{B_1,B_2\}\), and let \(\ell\asymp R^{-2}\)
be the fluctuation-field scale with fixed block aspect \(b/R\).
For \(R\) sufficiently large, the source-free activity of
Proposition~\ref{prop:forest} obeys
\[
 |K(X)|_{\ell,0}\ge c'R^{-4}.
\]
Here this activity is compared with the singleton factor prescribed
by the forest representation.
\end{proposition}

\begin{proof}
Write \(J_t=J_X(t)\), \(S_t=(I-J_t)^{-1}\), and
\(r_t=S_t\boldsymbol 1\), where \(\boldsymbol 1\) is the constant
field on \(\Lambda_X\). For \(t=0,1\), the mean-field deficit
computed in \eqref{eq:entsextic} has degree-four coefficient
\(\sum_i r_{t,i}^4/12\) at constant field \(h=z\boldsymbol 1\).
The argument for \eqref{eq:entsextic} applies without translation
invariance, as it only uses \(S_t=(I-J_t)^{-1}\). The comparison
of exact and mean-field pressure jets shows that the degree-four
coefficient of \(W_X(t,z\boldsymbol 1)\) differs from this by
\(O(R^{-4}|\Lambda_X|)\). Its degree-two coefficient is
\(O(R^{-4}|\Lambda_X|)\) as well, from the Hessian comparison and
\(|c-1|\lesssim R^{-4}\).

Entrywise positivity of the walk expansion gives
\(S_1-S_0\ge J_1-J_0\ge0\) and \(r_{t,i}\ge1\). There are at least
\(cR^4\) sites in the first block in a slab next to the common face
and away from the other faces. For each such site \(i\), at least
\(cR^4\) sites of \(B_2\) lie within its interaction box. Thus
\(((J_1-J_0)\boldsymbol 1)_i\ge c'\) on the slab. It follows that
\[
 \sum_i(r_{1,i}^4-r_{0,i}^4)\ge4\sum_i(r_{1,i}-r_{0,i})\ge c''R^4.
\]
For two blocks the exact tree integral is
\(K(X)=e^{-W_X(1,h)}-e^{-W_X(0,h)}\). Because \(W_X(t,0)=0\)
and it is even, the degree-four coefficient of \(e^{-W_X(t,z\boldsymbol1)}\)
is minus the quartic coefficient of \(W_X\) plus half the square
of its quadratic coefficient. After scaling \(z\) by \(\ell\),
the exact-minus-mean-field quartic terms and the quadratic squares
are both \(O(R^{-8})\). The mean-field quartic difference has
magnitude at least \(c''R^4\ell^4\asymp R^{-4}\). Evaluating the
fourth field Taylor coefficient in the scaled constant direction
is bounded in absolute value by the Taylor seminorm. This proves
the claimed lower bound.
\end{proof}

\begin{corollary}[A mixed quartic term cannot be removed by block-local sextics]
Let \(X=\{B_1,B_2\}\) be as in Proposition~\ref{prop:rawforest}.
There are nonnegative smooth lattice functions \(f_i\), supported
in \(B_i\), such that, with \(h=\ell(t_1f_1+t_2f_2)\),
\[
 [t_1^2t_2^2]e^{-W_X(1,h)}\le -c g,
 \qquad g=R^{-4},
\]
for all sufficiently large \(R\). If \(I_i(h_{B_i})\) are even
block-local factors, normalized by \(I_i(0)=1\), with quadratic
Taylor coefficients of size \(O(g)\) in these directions, then
\[
 \left|[t_1^2t_2^2]\left(e^{-W_X(1,h)}-I_1I_2\right)\right|
 \ge c' g.
\]
In particular this difference cannot have an \(O(g^3)\) Taylor norm.
\end{corollary}

\begin{proof}
Choose patches of side proportional to \(R\) on opposite sides
of the common face, and smooth bumps \(f_i\) bounded below on
smaller patches. Choose their supports so that every pair of
points in the two smaller patches is at distance at most \(R\).
Then \(\|\nabla^k f_i\|_\infty\le C_kR^{-k}\).
For \(S=(I-J_X(1))^{-1}\), the coefficient in the leading quartic is
\[
 [t_1^2t_2^2]\frac1{12}\sum_x(Sh)_x^4
 =\frac{\ell^4}{2}\sum_x(Sf_1)_x^2(Sf_2)_x^2\ge cR^4\ell^4.
\]
Indeed \(Sf_1\ge f_1\), and \(Sf_2\ge J_X(1)f_2\ge c\)
on the first smaller patch. The reference comparison bounds the
exact quartic error by \(O(g^2)\) in these directions. All coefficients
of \(W_X^{(2)}(1,h)\) are \(O(g)\), so its square contributes only
\(O(g^2)\) when exponentiating. This proves the first assertion.
The mixed coefficient of \(I_1I_2\) is the product of their two
quadratic coefficients, hence is \(O(g^2)\). This proves the second.
By rescaling the bumps by a fixed constant depending on the block
aspect and the finite test-derivative order, the symmetrized tensor
\(\ell^4 f_1^{\otimes2}\otimes f_2^{\otimes2}\) is an admissible
test tensor in the RG Taylor seminorm as well. A degree-six prefactor
does not change the degree-four coefficient.
\end{proof}

Thus the raw connected expansion has a two-block term already
at the order of the marginal quartic coupling. Improving bounds on
large polymers or increasing a regulator coefficient cannot remove
this low-degree term. A different reference factor or a localization
of the two-block quartic is needed before applying the small-activity
contraction.

\begin{proposition}[Large-set decay at zero field]\label{prop:forest_zero}
For every fixed \(\omega\in(0,1)\), one can choose a fixed block
aspect \(b/R\) sufficiently large that the activities above satisfy,
for \(|X|\ge2\),
\[
 |K(X)|_{\ell,0}\le C_\omega R^{-4}\omega^{|X|},
 \qquad \ell\asymp R^{-2},
\]
for all sufficiently large \(R\). This assertion concerns the
zero-field Taylor seminorm, before the large-field part of the
RG norm is imposed.
\end{proposition}

\begin{proof}
For an active block link \(e\), let \(E_e\) be the set of microscopic
endpoints of reference edges carrying its parameter. For a set
\(S\) of active link labels let \(\tau(i,T,S)\) be the shortest length
of a tree joining a site \(i\), a set \(T\) of field sites, and
at least one point of each \(E_e\) for \(e\in S\). The union of
witness trees proves submultiplicativity of the weight
\(e^{\eta\tau(i,T,S)/R}\). Changing the anchor across a reference edge
costs at most \(e^\eta\). In the recursion, a link parameter that
multiplies the coupling on an edge incident on \(i\) has tree
weight one at \(i\), since its support contains \(i\). Thus the
proof of joint analyticity in Proposition~\ref{prop:forest}
works also with this tree weight, upon taking \(\eta>0\) small enough
that \(e^\eta q_*<1\), and then taking the Taylor radius \(r\) small.
In particular the pressure may be expanded at zero field by
anchoring each field monomial at the site where the magnetization
is integrated along the field. In that representation the degree-two
coefficient rows of \(W_X\) have joint weighted norm \(O(R^{-4})\),
and all coefficients of field degree at least four have a
uniform joint weighted norm. There are no terms of field degree zero
or odd degree. The assertions include nonconstant link coefficients,
by the joint comparison with the mean-field Hessian.

We only need square-free link monomials, since the forest formula
differentiates each of its marked links once. Work in the quotient
of the link power-series algebra by the ideal generated by the
squares of its variables. Fix an arbitrarily large coefficient
radius \(D>1\) for this quotient algebra. The following elementary
geometry allows the change from radius \(r\) to radius \(D\):
\[
 |S|\le C_1\bigl(1+b^{-1}\tau(i,T,S)\bigr).
\]
Indeed, choose a witnessing point in \(E_e\) for each marked link.
Every block containing a witnessing point accounts for at most
\(\Delta\) distinct labels. A fixed positive fraction of the distinct
blocks can be chosen pairwise at block distance at least three, by
greedily removing the bounded block neighborhoods. Their witnessing
points have mutual periodic distance at least \(b\). Disjoint balls
of radius \(b/3\) around these points require a tree visiting all
of them to have length at least a constant times \(b\) times
one fewer than their number. This proves the inequality,
including on a torus, with an absolute adjustment for one witness.
Choose \(b/R\) sufficiently large that
\[
 (D/r)^{|S|}e^{-\eta\tau(i,T,S)/R}\le C_D.
\]
For each anchor this bounds the unweighted square-free link
coefficient norm at radius \(D\) by a fixed multiple of the
old joint norm. Set \(\epsilon=R^{-4}\). In the same algebra,
with the field Taylor norm at zero of radius \(\ell\), summing over
anchors now gives
\[
 \|W_X(t+v,h)\|_{\ell,D}
    \le C_D|\Lambda_X|(\epsilon\ell^2+\ell^4)
    \le C'_D \epsilon |X| .
\]
The constant may depend on the chosen fixed aspect. We use a
quotient norm only in the link variables; the field Taylor
norm is the full absolute coefficient sum. The exponential
therefore obeys
\[
 \|e^{-W_X(t+v,h)}-1\|_{\ell,D}
      \le C'_D\epsilon |X| e^{C'_D\epsilon |X|}.
\]
Extract the square-free coefficient indexed by a spanning tree,
which costs \(D^{-(|X|-1)}\). Sum over at most
\(\Delta^{|X|-1}\) trees. Since \(D\) can be made arbitrarily
large, the resulting bound on \(A_X\) has any prescribed
geometric decay factor, times \(C_D\epsilon\), after absorbing
\(|X|\) and the factor \(e^{C'_D\epsilon |X|}\) into the rate.
For a partition into \(k\) connected parts of size at least two,
retain the factors \(D^{-(|C|-1)}\) from the estimate for each part.
Their product is at most \(D^{-|X|/2}\). Also
\((C'_D\epsilon)^k\le C'_D\epsilon\) once \(R\) is large enough
that \(C'_D\epsilon\le1\), and the products of the remaining factors
\(|C|e^{C'_D\epsilon|C|}\Delta^{|C|-1}\) cost only a fixed
exponential in \(|X|\). This exponential and the connected-partition
count are independent of \(D\) after increasing \(R\) as above.
Choose \(D\) large to prove the bound for \(K\).
\end{proof}

\begin{proposition}[Large-set decay with the decaying regulator]\label{prop:forest_large}
Fix \(\omega\in(0,1)\). For a sufficiently large fixed block aspect
\(M=b/R\), set \(a=b^{-1}\), take \(\ell=R^{-2}\), and put
\[
 H_X(h)=\exp\left(-\kappa a^2\sum_{i\in\Lambda_X}|h_i|^2\right).
\]
Let \(\widetilde G_X\) be the product of the affine-quotient
regulators at block side \(b\) and fluctuation scale \(\ell\), with
at least four lattice derivatives in the defining test norm. Then
for some fixed \(\kappa>0\), all real fields, and \(|X|\ge2\),
\[
 |K(X)|_{a,h}\le C_\omega\omega^{|X|}H_X(h)\widetilde G_X(h).
\]
The bound can include up to two field-linear nilpotent sources with
\(a\tau_z\) uniformly bounded. The same bounds hold for the individual
connected tree activities. This is a large-set estimate for the
raw forest chart, without the \(g^3\) smallness conclusion.
\end{proposition}

\begin{proof}
Use the anchored joint norm in the proof of
Proposition~\ref{prop:forest_zero}, now at an arbitrary real
field. Let \(D>1\) be the square-free link coefficient radius.
Take \(M\) so large that
\[
 (D/r)^{|S|}e^{-\eta\tau(i,T,S)/R}
 \le C_D e^{-\eta\tau(i,T,S)/(2R)}.
\]
Here \(C_D\) can be chosen independently of \(M\) once this
lower bound on \(M\) is met. If \(S\ne\varnothing\), at least one
microscopic link endpoint in its witness tree lies within
\(R\) of a coarse block boundary. Define
\(w(i)=\exp(-\eta'\operatorname{dist}(i,\partial B_i)/R)\),
where \(B_i\) is the block containing \(i\), and choose
\(0<\eta'\le\eta/2\). The leftover tree factor is bounded by
\(Cw(i)\). Thus at any real background the absolute sums of the
field derivative rows of positive link degree, measured at link
radius \(D\), have a bound \(C_Dw(i)\) at anchor \(i\) for each
fixed derivative order at least two. The first derivative need not
be bounded at arbitrary field; below it is used only at a field
of supremum norm at most one. When summed in field Taylor degree at
radius at most \(r/2\), their coefficient bounds are geometric.
For the Hessian at zero the same statement has the extra factor
\(\epsilon=R^{-4}\), by joint comparison with the mean-field Hessian.

Let \(\mathcal Q(v,h)=W_X(t+v,h)-W_X(t,h)\) in the square-free
link algebra. We claim
\begin{equation}\label{eq:forest_shell_Q}
 \|\mathcal Q(v,\cdot)\|_{a,h;D}
 \le C_D\sum_{i\in\Lambda_X}w(i)
          \big(\Psi(h_i)+\epsilon h_i^2+a^4\big).
\end{equation}
For the quadratic term at zero, use the weighted Hessian
bound and symmetry of its field indices, bounding
\((|h_i|+a)(|h_j|+a)\) by the average of their squares.
For the remaining term subtract this quadratic. Its field
Taylor coefficients through degree three vanish at zero.
For derivative order \(k\le3\) at the field restricted to
\(|h_i|\le1\), integrate the fourth derivative along the
segment from zero. Bound each field product times \(a^k\)
by a constant times the sum of the fourth powers of the
participating field coordinates and \(a^4\). Because the
fourth derivative is symmetric in all field indices, anchor
at the index carrying each fourth power before using the row
bound. The \(a^4\) terms use the sum of the rows. Next move
to the full field. For \(1\le k\le3\), the difference is
controlled by the \((k+1)\)-st derivative, anchoring at the
index of the large field increment. For \(k=0\), use the
first derivative at the restricted field and the integrated
Hessian for the increment. The first derivative is bounded
rowwise by \(C_Dw(i)\) since the restricted field has supremum
at most one. Thus the cost on the large coordinates is bounded
by \(C_D\sum_{|h_i|>1}w(i)(|h_i|+h_i^2)\), which is at most a
constant times the \(\Psi\) contribution. All Taylor orders
\(k\ge4\) together are bounded by \(C_Da^4\sum_i w(i)\) using the
geometric row bounds. This proves \eqref{eq:forest_shell_Q}.

We record the elementary boundary estimate used to control
\eqref{eq:forest_shell_Q}. For \(p=4\) and each block \(B\), write
\(D_{p,B}=\max_e b^p\|\nabla_e^p h\|_{\infty,B}\), where only
forward differences whose full stencil lies inside \(B\) are
included. Then
\begin{equation}\label{eq:forest_trace}
 \sum_{i\in B}w(i)\Psi(h_i)
 \le C M^{-1/2}\sum_{i\in B}\Psi(h_i)
       +C|B|M^{-1-p}D_{p,B}^2 .
\end{equation}
To prove this, bound \(w\) by the sum of the corresponding weights
for the eight faces. For one face and a site \(i\) within \(b/4\)
of it, let \(e\) point inward and set \(q=\lfloor b/\sqrt M\rfloor\).
For \(1\le s\le q\) the \(p\)-th finite-difference identity writes
\(h_i\) as a linear combination of \(h_{i+lse}\), \(1\le l\le p\),
with coefficients of bounded absolute sum and error at most
\(s^p\|\nabla_e^p h\|_{\infty,B}\). Take \(M\) large enough that
all the stencils are inside the block. Convexity and fixed-scale
bounds on \(\Psi\), together with \(\Psi(u)\le u^2\), give a bound
on \(\Psi(h_i)\) by the corresponding sum of \(\Psi(h_{i+lse})\)
and \(C(q/b)^{2p}D_{p,B}^2\). Average over \(s\).
For fixed \(l\) and a given image site, the weights of its
possible preimages are bounded by a geometric sum of order \(R\)
along the normal line. The average therefore costs \(CR/q\).
The weighted count of sites in the block is \(O(|B|/M)\).
Sites more than \(b/4\) from the face contribute an
exponentially small coefficient times \(\sum_B\Psi\).
Summing over faces proves \eqref{eq:forest_trace}.

The affine-quotient regulator dominates
\(\exp(c\ell^{-2}D_{p,B}^2)\). Indeed, subtracting any affine
function preserves the \(p\)-th differences inside \(B\),
and the field norm on the block neighborhood includes their
normalized supremum. Since \(|B|\ell^2\asymp M^4\), the last
term in \eqref{eq:forest_trace} costs at most
\(CM^{3-p}\log\widetilde G_B\). Also, for large \(R\) at fixed
\(M\), \(\epsilon s^2\le\Psi(s)+C\epsilon^2\) and
\(\epsilon^2\le a^4\). First fix \(D\), then take \(M\)
sufficiently large. Equations \eqref{eq:forest_shell_Q} and
\eqref{eq:forest_trace} imply
\[
 \|\mathcal Q(v,\cdot)\|_{a,h;D}
 \le \delta\sum_{i\in\Lambda_X}\Psi(h_i)
       +\tfrac12\log\widetilde G_X(h)+|X|
\]
with any prescribed fixed \(\delta>0\). In the last term we used
\(a^4|B|=1\) and \(\sum_B w\le C|B|/M\). The constants in the
real restricted-weight stability bound are independent of
\(D,M\), provided \(R\) is sufficiently large depending on \(M\).
In fact its proof applies at any \(a\le R^{-1}\), with the
quadratic loss contributing at most \(O(|\Lambda_X|\epsilon^2)\),
which is at most \(O(|\Lambda_X|a^4)\) here.

Use submultiplicativity in the square-free link algebra to
exponentiate \(\mathcal Q\), and multiply by the real weight.
Taking \(\delta\) smaller than half the real damping coefficient
bounds the resulting joint norm by
\[
 e^{C|X|}\exp\left(-c'\sum_i\Psi(h_i)\right)
                  \widetilde G_X(h)^{1/2},
\]
where \(C\) is independent of \(D\). By
\(\Psi(s)\ge c a^2 s^2-Ca^4\), the exponential supplies
\(H_X\) with slack. For a source on \(B\), its extra polynomial
is controlled with this slack and the unused fraction of
\(\widetilde G_B\). Explicitly, the affine-quotient inequality
in the proof of Proposition~\ref{prop:local} bounds \(|h_z|/a\)
by a constant times \(a\|h_B\|_2\) and the affine deviation divided
by \(a\), since \(|B|a^4=1\) and \(\ell\le a\) for large \(R\).
The field-linear source factors can therefore be absorbed uniformly
for \(a\tau_z\) bounded.

Finally extract the square-free coefficient of each spanning tree,
which costs \(D^{-(|X|-1)}\). This leaves an exponential constant
per block independent of \(D\), including the number of spanning
trees and the normalized damping loss. In the partition formula
for \(K\), the product of the factors \(D^{-(|C|-1)}\) is at
most \(D^{-|X|/2}\), since all parts have size at least two.
The partition count costs only another fixed exponential.
Choose \(D\) large to obtain the prescribed geometric factor,
then choose \(M\) as above.
\end{proof}


% Uses the cavity, covariance and regulator notation of spread_out_ising_rg.tex.
\section{Localizing the mixed quartic}

For a symmetric quartic tensor \(C(i_1,i_2,i_3,i_4)\), write
\[
 Q_4(\phi)=\sum_{\boldsymbol i}C(\boldsymbol i)\prod_{k=1}^4\phi_{i_k},
 \qquad q_i=\sum_{i_2,i_3,i_4}C(i,i_2,i_3,i_4).
\]
We use ordered indices. The tensor includes whatever factorial
is used in the convention for the Taylor polynomial. Assume that
its anchored exponential moment is bounded, in the sense that
\begin{equation}\label{eq:quartic-tensor-local}
 \sup_{i_1}\sum_{i_2,i_3,i_4}|C(\boldsymbol i)|
           e^{\eta\tau(\boldsymbol i)/R}\le C_4,
\end{equation}
where \(\tau(\boldsymbol i)\) is the length of a shortest lattice
tree joining the four sites, measured in the periodic lattice metric
when the volume is a torus.

\begin{proposition}[A two-difference quartic remainder]\label{prop:quartic-local}
The exact localization
\begin{equation}\label{eq:q4exactlocal}
 \sum_i q_i\phi_i^4-Q_4(\phi)=\sum_{\boldsymbol i}C(\boldsymbol i)
   \left(\frac14\sum_{k=1}^4\phi_{i_k}^4-\prod_{k=1}^4\phi_{i_k}\right)
\end{equation}
has a representation by quartic polynomials \(D_4(X)\) supported
on connected block sets \(X\) of side \(b\). Every monomial in the
representation contains at least two first differences after path
telescoping. If the block test norm has at least two discrete
derivatives, then for any fixed \(\varrho\in(0,1)\), for \(b/R\)
sufficiently large,
\begin{equation}\label{eq:q4testlocal}
 \|D_4(X)\|_{T_{\mathfrak h,\phi}}
 \le C_\varrho b^4\mathfrak h^4 (R/b)^2\varrho^{|X|}
                    (1+N_X(\phi;\mathfrak h))^4 .
\end{equation}
Here \(N_X\) is the maximum of \(|\phi_i|/\mathfrak h\) on
\(\Lambda_X\) and of \(b|\nabla_e\phi_i|/\mathfrak h\) over unit
steps with both endpoints in \(\Lambda_X\). The estimate is
uniform in the volume. With \(\ell_b=(Rb)^{-1}\) and \(g=R^{-4}\),
its zero-field prefactor is \(g(R/b)^2\).
\end{proposition}

\begin{proof}
Symmetry of \(C\) gives \eqref{eq:q4exactlocal}. Set \(u=\phi_{i_1}\)
and \(d_k=\phi_{i_k}-u\) for \(k=2,3,4\). The expression in
parentheses is exactly
\begin{align*}
 &u^2\left(\tfrac32\sum_{k=2}^4d_k^2-
                         \sum_{2\le k<l\le4}d_kd_l\right)\\
 &\hspace{1cm}+u\left(\sum_{k=2}^4d_k^3-d_2d_3d_4\right)
                         +\tfrac14\sum_{k=2}^4d_k^4 .
\end{align*}
For each tuple, choose coordinate shortest paths from \(i_1\)
to the other three sites. Average over coordinate orders and
periodic shortest-path ties when symmetry is needed. The block
set visited by their union is connected, and has at most
\(C(1+\tau(\boldsymbol i)/b)\) blocks. In fact each coordinate
segment of length \(s\) meets at most \(2+s/b\) blocks, and
the total path length is at most \(C\tau(\boldsymbol i)\).
Assign the tuple's contribution to its visited block set and
write each \(d_k\) as a sum of unit first differences along
its path. This defines \(D_4(X)\) and proves the exact identity.

In the Taylor pairing of a degree-four monomial, a difference
factor is either evaluated on \(\phi\) or acts as a difference
in one argument of a test tensor. The former costs at most
\(\mathfrak h N_X/b\) for a unit step. The latter costs
\(b^{-1}\) in the test-function norm. When there are two or
more such factors, they act on distinct slots of the tensor,
including when the underlying sites coincide. Use two of these
differences for the gain, and if there are more, bound each extra
difference by the sum of its two endpoint evaluations. There
are only finitely many choices of differentiated factors. The
Taylor seminorm contribution of a tuple is therefore bounded by
\[
 C\mathfrak h^4(1+N_X)^4(\tau(\boldsymbol i)/b)^2.
\]
Reserve half the exponential weight in \eqref{eq:quartic-tensor-local}
for the second moment, giving \(C(R/b)^2\). The other half gives
\(\exp(-c b|X|/R)\) up to a fixed correction in the exponent
for the bounded-size case. Equivalently, by taking \(b/R\) large,
one can bound the remaining factor by \(C_\varrho\varrho^{2|X|}\),
with \(C_\varrho\) independent of larger \(b/R\). There are at most
\(b^4|X|\) choices of the anchor \(i_1\) for a given \(X\).
Absorb \(|X|\) in the geometric factor. This gives
\eqref{eq:q4testlocal}.
\end{proof}

For the exact box pressure the leading tensor at every real
forest interpolation point is
\[
 C_{\rm mf,t}(\boldsymbol i)=\tfrac1{12}\sum_x
                         \prod_{k=1}^4 S_t(x,i_k),
 \qquad S_t=(I-J_X(t))^{-1}.
\]
Each row of \(S_t\) has uniformly bounded exponential moment
at length \(R\) by its geometric series, since the weighted row
sum of \(J_X(t)\) is below one for small enough \(\eta\). Symmetry
and a star tree at \(x\) prove \eqref{eq:quartic-tensor-local} for
\(C_{\rm mf,t}\). The cavity comparison proves the same bound with
right side \(O(g)\) for the exact-minus-mean-field quartic tensor.
For the leading tensor the localized coefficient is
\[
 q_{i,t}^{\rm mf}=\tfrac1{12}\sum_x S_t(x,i)(S_t\mathbf 1)_x^3
 \ge \tfrac1{12}.
\]
The exact coefficient differs by \(O(g)\), so the sitewise localized
quartic is uniformly positive even in an induced subvolume.
At full translation invariance it is constant, with leading value
\(\tfrac1{12}(1-\beta q)^{-4}\), where \(q=\sum Q\) is the
reference interaction mass. At every real interpolation point
the full leading quartic is also positive uniformly in \(\sum_i\phi_i^4\),
by bounded invertibility of \(S_t\) on \(\ell^4\). Thus interpolation
between the local quartic and the full quartic preserves quartic
coercivity for large \(R\).

The polynomial estimate can also be used inside a stable weight.
For example, multiply \(D_4(X)\) by the product of the positive
localized quartic factors on \(\Lambda_X\). At large-field scale
\(a=b^{-1}\), the extra factor \((1+N_X(\phi;a))^4\) is absorbed
by unused Gaussian damping and the affine-quotient regulator.
Indeed, for each block neighborhood subtract an affine function.
The affine function's supremum and scaled slope on that
neighborhood are bounded by its normalized block \(L^2\) norm,
while the deviation and its scaled derivative are bounded by
the quotient seminorm. The identity \(b^4a^4=1\) converts
the normalized \(L^2\) term into the Gaussian damping factor.
At the fluctuation scale the growing regulator absorbs the same
polynomial. The submultiplicative Taylor seminorm of the positive
quartic weight contributes at most a fixed exponential in
\(|X|\), which can be absorbed by starting with a stronger
geometric factor in \eqref{eq:q4testlocal}. This does not by
itself give the circle-product estimates for a renormalized
chart, where one has to integrate the complete weights.

The gain \((R/b)^2\) is the natural size of this localization.
On a translation-invariant torus take the real field
\(\phi_i=\ell_b\cos(k\cdot i)\), with a nonzero allowed frequency
\(|k|\asymp b^{-1}\). The normalized smoothing operator
\(T=(1-\beta q)S_0\) multiplies this field by
\((1-\beta q)/(1-\widehat J(k))\). The box reference has
second moment comparable to \(R^2\), so for \(b/R\) large,
\(1-\widehat T(k)^4\asymp (R/b)^2\). Thus the mean-field
quartic difference on this field, per block, is of order
\(b^4\ell_b^4(R/b)^2\). With fixed \(b/R\) the
localized activity need not have size \(O(g^3)\).

\section{An exact four-difference quadratic remainder}

For a hypercubic-symmetric translation-invariant kernel \(E\) with
\(\sum_x|E(x)|e^{\eta|x|/R}=O(g)\), the quadratic form with symbol
\(\widehat E\) can be localized to its constant and nearest-neighbor
kinetic terms, with a four-difference remainder. Here is an
explicit finite-difference representation. Put
\(d(k)=1-\cos k\) and \(S_n(k)=\sum_{a=0}^{n-1}e^{iak}\) for
\(n\ge1\). Then
\[
 1-\cos(nk)=d(k)|S_n(k)|^2,\qquad
 \cos(nk)-1+n^2d(k)
   =2d(k)^2\sum_{m=1}^{n-1}(n-m)|S_m(k)|^2 .
\]
The first identity follows by factoring \(1-e^{ink}\). The second
uses \(|S_n|^2=n^2-2\sum_{m=1}^{n-1}(n-m)(1-\cos(mk))\).
The multiplier after \(d(k)^2\) is a finite-range convolution
kernel of absolute sum \(O(n^4)\). Coordinate reflections allow
us to replace \(\cos(k\cdot x)\) in \(\widehat E(k)\) by
\(\prod_i\cos(k_i x_i)\). Expand this product as a product of
\(1+(\cos(k_i x_i)-1)\). Its linear terms produce
\(-\sum_i x_i^2d(k_i)\) and the one-dimensional remainders above.
Every term with at least two factors \(\cos(k_i x_i)-1\)
contains a product of two \(d(k_i)\)'s by the first identity,
with a finite-range kernel of absolute sum \(O(|x|^4)\)
after choosing any two of its difference factors and bounding
the remaining factors by their absolute kernel sums.

Consequently the exact residual multiplier is a sum of
terms \(d(k_i)d(k_j)\widehat R_{ij}(k)\) after subtracting
\(\widehat E(0)\) and its kinetic coefficient. The kernels
\(R_{ij}\) have absolute sums \(O(g R^4)\) and retain an
exponential moment at a smaller rate. At the scale \(b\),
four discrete derivatives give the zero-field quadratic
prefactor
\[
 b^4\ell_b^2\,b^{-4}gR^4=g(R/b)^2 .
\]
This is an exact localization of the translation-invariant
kernel, not a quadratic localization at arbitrary noninvariant
forest interpolation points.


\section{A local quartic coordinate and exact quadratic absorption}

In the box preprocessing, put \(B=I-J\), \(S_0=B^{-1}\), and
\(E_\beta=S-S_0\), where \(S=D^2F(0)\). Make the exact linear field
change \(h=Bu\). The Gaussian precision in the new variable and the
quadratic coefficient of the residual interaction are
\begin{align*}
 B C_m^{-1} B&=A^{-1}B M_m,\\
 H&=(c-1)B-BE_\beta B .
\end{align*}
Indeed the residual interaction is
\(\frac c2\langle u,Bu\rangle-F(Bu)\). Write \(g=R^{-4}\), and
work in \(|\beta-1|+m^2\le C_0g\). The cavity estimate and the
finite-range bounds on \(B\) give
\begin{equation}\label{eq:fc_Hnorm}
 \sum_x |H(x)|e^{\eta'|x|/R}\le Cg
\end{equation}
for a fixed \(\eta'>0\). Set \(r=\widehat H(0)\) and \(K=H-rI\).
For a dimensionless parameter \(z\) and \(\mathcal D=-\Delta_{\rm nn}\),
absorb \(K+zR^2\mathcal D\) in the Gaussian precision. The new covariance
is exactly
\begin{equation}\label{eq:fc_newcov}
 \mathscr C=A T^{-1},\qquad
 T=BM_m+A(K+zR^2\mathcal D).
\end{equation}
The remaining quadratic interaction is
\[
 V^{(2)}_{r,z}(u)=\tfrac12r\sum_i u_i^2
       -\tfrac12zR^2\sum_i u_i(\mathcal D u)_i,
\]
so its Hessian is \(rI-zR^2\mathcal D\). All operators in these formulas commute.

\begin{proposition}[Uniform finite-range decomposition after Hessian absorption]\label{prop:fc_newcovbound}
Fix finite derivative and mass-decay orders \(k_*,N_*\). For fixed
sufficiently small \(|z|\), the covariance \eqref{eq:fc_newcov} has
positive finite-range pieces \(\mathscr C_j\), \(j\ge0\), with range
at most \(C R L^j\), such that for \(0\le k\le k_*\),
\[
 \|\nabla^k\mathscr C_j\|_\infty
 \le C R^{-4-k}L^{-(2+k)j}(1+m^2L^{2j})^{-N_*}.
\]
The constants are independent of \(R,\beta,m^2,z\) in the stated
parameter domain, and the pieces are continuous in the covariance
parameters. There is no logarithmic growth of the initial range.
\end{proposition}

\begin{proof}
Use \(k\) for the frequency variable in this proof. From
\eqref{eq:fc_Hnorm} and symmetry,
\(|\widehat K(k)|\le Cg\min(R^2|k|^2,1)\). Also
\(R^2\widehat A(k)\widehat{\mathcal D}(k)\le C\min(R^2|k|^2,1)\),
by rapid Fourier decay of the sampled smooth kernel \(A\).
Since \(B\) is uniformly positive and \(M_m\) has the box multiplier
comparison, \(\widehat T(k)\asymp m^2+\min(R^2|k|^2,1)\).

Choose a fixed integer \(P\), large compared to \(k_*+2N_*\).
There is a hypercubic-symmetric finite-range kernel \(K_f\) supported
in a fixed multiple of the range \(R\), with norm \(O_P(g)\) in
\eqref{eq:fc_Hnorm}, whose moments through degree \(P\) equal those
of \(K\). To construct it, choose a smooth compactly supported bump
positive on an open box. The Gram matrix of the monomials of degrees
at most \(P\) against this bump is positive definite. Linear combinations
of the monomials times the bump therefore give smooth compactly
supported functions dual to the moment functionals. Sample them at
scale \(R\) with the factor \(R^{-4}\). The matrix of the sampled,
normalized moments converges to the identity. Solving the moment
equations costs bounded coefficients times \(g\), since the target
moments divided by their powers of \(R\) are \(O_P(g)\).
Average the result over the hypercubic symmetry group. This preserves
the target moments and the norm bound.

Let \(T_f=BM_m+A(K_f+zR^2\mathcal D)\). It is finite range and has
the same multiplier comparison as \(T\), since \(K_f\) has zero row
sum and uniformly bounded scaled second moments of order \(g\).
The positive polynomial decomposition of the inverse of \(T_f\),
multiplied by \(A\), gives positive pieces \(\Gamma_i^f\) of range
\(O_P(RL^i)\). The Fourier integrals in Proposition~\ref{prop:frd}
give the bounds claimed here for \(\Gamma_i^f\). They also give the
same upper bound on
\(\int |k|^k |\widehat\Gamma_i^f(k)|\,dk\), with the exponent in
\(|k|^k\) understood as the requested derivative order.

Consider the positive multiplier
\[
 M(k)=\frac{\widehat T_f(k)}{\widehat T(k)}
 =1-\frac{\widehat A(k)(\widehat K(k)-\widehat K_f(k))}{\widehat T(k)}.
\]
At \(m^2=0,k=0\) use its continuous extension. For \(|\alpha|\le P-2\),
\begin{equation}\label{eq:fc_L1ratio}
 \int |\partial_k^\alpha(M(k)-1)|\,dk\le C_P g R^{|\alpha|-4}.
\end{equation}
To verify this near the origin, use \(\xi=Rk\). The rescaled
denominator has uniformly bounded derivatives, its first derivative
is \(O(|\xi|)\) by symmetry, and it is comparable to
\(m^2+|\xi|^2\) on \(|\xi|\le1\). Differentiating its reciprocal
therefore costs \((m^2+|\xi|^2)^{-1-|\alpha|/2}\). Moment matching
and the weighted norms give the bound
\(C_Pg|\xi|^{P+1-|\alpha|}\) for the corresponding numerator
derivatives. The ratio derivatives are consequently bounded by
\(C_Pg|\xi|^{P-1-|\alpha|}\), also in the zero-mass limit.
For \(|\xi|\ge1\), the denominator is bounded below. Its
scaled derivatives are uniformly bounded by the moment estimates,
and the numerator derivatives have the rapid decay of the sampled
kernel \(A\) times \(O_P(g)\). Integrating in \(\xi\) and returning
to \(k\) proves \eqref{eq:fc_L1ratio}.

Fourier coefficient integration by parts in a coordinate of maximal
absolute displacement, using \eqref{eq:fc_L1ratio}, now shows that
truncating the real-even Fourier series of \(M-1\) at range \(RL^n\)
has uniform error at most \(\varepsilon_n=C_PgL^{-s n}\), where
\(s=P-6\). Denote the truncation plus one by \(M^{[n]}\). Set
\[
 F_n=M^{[n]}-3\varepsilon_n,\qquad
 M_0=F_0,\qquad M_n=F_n-F_{n-1}\quad(n\ge1).
\]
For small \(g\), \(F_0\ge1/2\). Also \(F_n\ge F_{n-1}\), since
\(M^{[n]}-M^{[n-1]}\ge-\varepsilon_n-\varepsilon_{n-1}\) and
\(3(\varepsilon_{n-1}-\varepsilon_n)\ge
\varepsilon_n+\varepsilon_{n-1}\) for \(L^{-s}\le1/2\).
Thus the \(M_n\) are positive trigonometric polynomials of the
required ranges. Their sum is exactly \(M\), and for \(n\ge1\),
\(\|M_n\|_\infty\le C_{P,L}gL^{-s n}\).

Define \(\mathscr C_j\) by the Fourier multiplier sum of
\(\widehat\Gamma_i^f M_n\) over \(\max(i,n)=j\). These products
preserve positivity and have range \(O_P(RL^j)\). For \(i=j\),
sum over \(n\le j\) and use \(\sum_{n\le j} M_n\le M\lesssim1\)
and the weighted Fourier integral bound on \(\Gamma_j^f\).
For \(n=j>i\), use \(\|M_j\|_\infty\lesssim gL^{-s j}\) and
the sum of the Fourier integral bounds on the \(\Gamma_i^f\),
which is at most \(CR^{-4-k}\) at derivative order \(k\).
Choose \(P\) with \(s>2+k_*+2N_*\). Since \(m^2\le1\), these
last terms also obey the asserted mass-decay estimate. Absolute
uniform bounds and the fixed truncations give parameter continuity.
The sum of all products is \(A T_f^{-1}M=A T^{-1}\), as required.
\end{proof}

\begin{proposition}[Size of the bubble increments]\label{prop:fc_bubble}
Let \(A_{\mathrm{sc}}\ge1\) be a field rescaling factor, put
\(\lambda=A_{\mathrm{sc}}^{-2}\), and let
\(C_j=\lambda R^2\mathscr C_j\) be the covariance pieces for
\(\psi=(R/A_{\mathrm{sc}})u\). Write \(w_j=\sum_{i\le j}C_i\) and
\(w_{-1}=0\). For any fixed \(a_0>0\) define
\[
 B_j=a_0\sum_x \big(w_j(x)^2-w_{j-1}(x)^2\big).
\]
Then \(B_j\ge0\), and for every fixed decay order allowed by the
decomposition,
\[
 B_j\le C\lambda^2(1+m^2L^{2j})^{-N_*}.
\]
The base positive polynomial decomposition can be chosen so that
\(B_j\ge c\lambda^2\) whenever \(m^2L^{2j}\le c_0\) for a fixed
\(c_0>0\). The constants are independent of \(R,\lambda\), and the
covariance parameters in Proposition~\ref{prop:fc_newcovbound}.
\end{proposition}

\begin{proof}
The range and pointwise bounds give
\(\|C_j\|_1\le C\lambda R^2L^{2j}\) and
\(\|C_j\|_\infty\le C\lambda R^{-2}L^{-2j}(1+m^2L^{2j})^{-N_*}\).
Expanding the difference of squares, the mixed products have nonnegative
sums by Plancherel and positive semidefiniteness. Their upper bound is
\(\|C_j\|_\infty\sum_{i\le j}\|C_i\|_1\), times a fixed constant.

For the lower bound it suffices to retain the square of \(C_j\).
Use the construction of the positive polynomials in BPR, Lemma~3.8,
with \(f\) the square of the Fourier transform of a nonnegative bump,
normalized as there. Thus \(f(0)>0\). For \(t>1\) the nonnegative
periodization in (3.36)--(3.37) is bounded below by
\(f(t\arccos(1-s/2))\), where \(s\) is the scalar argument of the
polynomial. Rescale \(T_f\) by a fixed factor into the spectral
interval in that lemma. Its multiplier satisfies
\(s(k)\le C(m^2+R^2|k|^2)\) for \(|k|\le R^{-1}\). On a subinterval
of the scale-\(j\) time interval of length comparable to \(L^j\),
this lower bound is uniformly positive if
\(|k|\le c_1 R^{-1}L^{-j}\) and \(m^2L^{2j}\le c_0\), taking the
constants small. The time integral with measure \(t\,dt\) and
\(\widehat A(k)\ge c\) there give
\(\widehat\Gamma_j^f(k)\ge cL^{2j}\). At the initial scale the
constant positive polynomials for \(t\le1\) also give the bound.
Since \(M_0\ge1/2\) and all other terms in the product decomposition
are nonnegative, \(\widehat C_j(k)\ge c\lambda R^2 L^{2j}\) on
the same ball. Its volume is comparable to \(R^{-4}L^{-4j}\).
Plancherel now gives \(\sum_x C_j(x)^2\ge c\lambda^2\).
\end{proof}

The residual nonlinear interaction in the \(u\) variable has the
on-site mean-field quartic \(\sum_i u_i^4/12\). Its exact quartic
defect has anchored weighted row norm \(O(g)\) by the cavity comparison
and the finite-range field change. At field scale \(R^{-2}\) this
defect costs \(O(g^2)\) per range-scale block. The degree-six term
\(\sum_i u_i^6/30-\langle u^3,S_0u^3\rangle/18\) also costs \(O(g^2)\)
there. The exact-minus-mean-field sextic defect costs \(O(g^3)\).
Thus this change removes the order-\(g\) nonlocal quadratic and
quartic contributions. The next section constructs the remaining
order-\(g^2\) activity in a norm retaining large-field decay.

% Uses the cavity and forest sections and the exact absorbed covariance.
\section{Connected initialization in the Hessian-subtracted field}

For the induced reference graph and its real link array \(t\), let
\(B_t=I-J_X(t)\), \(S_t^{\rm ex}=D^2F_t(0)\), and put
\[
 \mathcal U_X(t,u)=\tfrac12\langle B_tu,S_t^{\rm ex}B_tu\rangle
                         -F_t(B_tu).
\]
Here \(F_t\) is the normalized reference pressure. At the full
volume and all links one, exact Hessian absorption leaves this
interaction plus the local mass and nearest-neighbor kinetic terms.
For a proper induced block set, \(\mathcal U_X\) is defined using
its own \(B_t\) and pressure; it is not obtained by restricting the
full-volume transformed interaction. The forest identity below is
applied with \(X\) equal to the entire ambient torus. Its induced
subvolume weights are auxiliary factors in that exact identity.
Set \(b=MR\), \(a\asymp b^{-1}\), \(g=R^{-4}\), and
\(\ell'\asymp A_{\rm sc} R^{-2}\), where \(A_{\rm sc}\ge1\) and \(A_{\rm sc}/R\to0\).
The comparability constants are fixed. Use the affine-quotient
regulator \(\widetilde G\) at fluctuation scale \(\ell'\), and
\(H_X=\exp(-\kappa(b^4a^2)^{-1}\sum_{i\in\Lambda_X}u_i^2)\).

\begin{proposition}[Exact quartic reference with an order-\(g^2\) seed]
\label{prop:transformed-init}
For every fixed \(\omega\in(0,1)\), one can choose a fixed aspect
\(M\) sufficiently large that the Hessian-subtracted weight at
all real fields has the representation
\[
 e^{-\mathcal U_X(\boldsymbol 1,u)}=(I_4\circ K)(X),
 \qquad I_4(B)=\exp\bigl(-\tfrac1{12}\sum_{i\in B}u_i^4\bigr).
\]
The polymer activity factors over geometric components. For every
nonempty \(Y\), uniformly for large \(R\),
\[
 |K(Y)|_{\ell',u}\le C_\omega A_{\rm sc}^6g^2\omega^{|Y|}G_Y(u),\qquad
 |K(Y)|_{a,u}\le C_\omega g^{1/2}\omega^{|Y|}H_Y(u)\widetilde G_Y(u).
\]
The estimates hold with up to two field-linear nilpotent sources
at sites, provided the source weights times the corresponding
field radius are bounded. If the physical local mass and kinetic
terms are included, with normalized coefficients \(|r|+|z|=O(g)\)
and kinetic term \(-zR^2\mathcal D\), the same estimates hold in
the derivative-sensitive \(T\) norms after multiplying the
reference and activity by the appropriate block factors. The
kinetic term may enlarge their field support to the one-step
block neighborhoods.
\end{proposition}

\begin{proof}
We first describe the analytic bounds for the transformed pressure.
At any real reference array, \(B_t\) and its inverse are bounded
in row and column absolute sum, uniformly in the induced volume.
Consequently the coercivity estimate in Corollary~\ref{cor:stable}
and the convexity and fixed-scale bounds for \(\Psi\) give
\(\mathcal U_X(t,u)\ge c\sum_i\Psi(u_i)\). Indeed,
\(u=B_t^{-1}(B_tu)\) and the weighted averaging inequality for
\(\Psi\) gives the needed lower bound. Its field derivatives
through order three vanish at zero, and its degree-four Taylor
polynomial is \(\sum_i u_i^4/12\) plus a tensor with absolute
weighted rows of order \(g\). The leading polynomial is independent
of all links, because the mean-field quartic is
\(\sum_i((I-J_X(t))^{-1}B_tu)_i^4/12\).

Fix a real field \(u\). At a link \(e=\{B,C\}\), put
\(r_e^0=(1+\sup_{B\cup C}|u|)^{-1}\). When taking Taylor
coefficients, regard these numbers as constants. Initially use
link increments \(r_e^0 v_e\). For any real field \(u'\) with
\(|u'_i|\le|u_i|\) and field increment \(y\), write
\[
 B_{t+r^0v}(u'+y)=B_tu'+B_ty-(J_{t+r^0v}-J_t)u'
                                      -(J_{t+r^0v}-J_t)y .
\]
At fixed small Taylor radii in \(v,y\), the nonconstant part
of this expression has a uniformly small anchored coefficient
norm. For the third term this follows from
\(r_e^0|u'_i|\le1\) on the support of link \(e\); for the other
terms use the absolute row bound and range of \(J\).
In the weighted tree norm, a field increment on a neighboring
site costs at most \(e^\eta\) to change its anchor, and a link
parameter incident on that site already has an endpoint in its
support there. Thus the cavity proof, including comparison with
mean field, applies in the same joint coefficient algebra with
these external fields. The real backgrounds are arbitrary, but
the size of their complex increments is uniformly small.
Differentiating the transformed pressure in \(u\) applies \(B_t\)
to the magnetization, and likewise uses the finite-range
coefficient matrices for link increments. This gives uniform
weighted absolute derivative rows of order at least two for
\(\mathcal U_X\), and analytic sums of all higher orders, uniformly
along every field segment with \(|u'_i|\le|u_i|\).
The covariance in its subtracted quadratic term has the same
analytic estimates by differentiating the reference magnetization
at zero. The exact-minus-mean-field quartic estimate holds in
this joint algebra, by the corresponding magnetization comparison.

We need a small link variation of the exponential in an algebra
with a large coefficient radius. The following choice avoids a
boundary trace estimate involving the coefficient of
\(\widetilde G\). Choose a large square-free link coefficient
radius \(D\). The order of choices in this proof is
\(D\), then \(\delta=\delta(D)\), then
\(\theta=\theta(D,\delta)\), then
\(M=M(D,\delta,\theta)\), and finally \(R\).
Constants with a subscript \(D\) may grow with \(D\); this growth
is absorbed by the final lower bound on \(R\). Call a block bad if its supremum of \(|u|\) exceeds
a fixed small \(\delta>0\). Mark a link if one of its endpoint
blocks is within two coarse blocks of a bad block, using the
periodic block metric. Use increments \(r_e v_e\), where
\(r_e=\theta r_e^0\) on marked links and \(r_e=r_e^0\) otherwise,
for a fixed \(0<\theta<1\) chosen below. The square-free link
algebra and anchored tree argument in
Proposition~\ref{prop:forest_zero} allow the change from the
small Taylor radius to \(D\), at the cost of a factor \(C_D\)
and using half the exponential tree weight, for \(M\) sufficiently
large. Coefficients involving a marked link gain a factor
\(\theta\) from the scaling. For coefficients with only unmarked
links and nonempty link support, every anchor in a bad block is
at distance at least \(cb\) from every link support involved.
The remaining tree weight bounds such rows at bad anchors by
\(C_D e^{-cM}\). There are at most a fixed number of marked
links per bad block, because the block graph has bounded degree.

Let \(\Delta\mathcal U=\mathcal U_X(t+rv,\cdot)-\mathcal U_X(t,\cdot)\).
The radii are kept fixed under all field and link derivatives at
the evaluation field \(u\).
We claim that in the joint field Taylor seminorm of radius \(a\)
and square-free link norm of radius \(D\),
\begin{equation}\label{eq:transformed-link-stability}
 \|\Delta\mathcal U\|_{a,u;D}
       \le \epsilon_0\sum_i\Psi(u_i)+C|X|,
\end{equation}
with \(\epsilon_0>0\) arbitrarily small and \(C\) independent of
\(D\) after choosing \(\delta,\theta,M\) in that order and then
\(R\) sufficiently large. To see this, split \(u=s+l\) into
coordinates with magnitude at most \(\delta\) and greater than
\(\delta\). Hold the adaptive radii fixed throughout the field
interpolation. At \(s\), Taylor expansion from zero through
order five, with integrated sixth derivative for the rest,
bounds the field Taylor terms of orders at most five by
\[
 C_D\sum_i\bigl(g(|s_i|+a)^4+(|s_i|+a)^6\bigr).
\]
This uses the absence of a link-dependent mean-field quartic.
For each product of field coordinates, symmetry of the
derivative tensor allows the row bound to be anchored at the
index carrying the largest power in the arithmetic--geometric
mean estimate. This gives at most
\(C_D(g+\delta^2)\sum_i\Psi(s_i)\) and
\(C_D|\Lambda_X|(ga^4+a^6)\). Choose \(\delta\) small first.
For increments from \(s\) to \(s+l\), apply the derivative rows
anchored at the index of the increment. For field Taylor order
zero, first estimate the first derivative at \(s\) by integrating
the Hessian from zero and then use the second-order formula for
the increment. For orders one through five, integrate the next
derivative directly. The resulting terms on the large coordinates
are bounded by \(C_{D,\delta}(\theta+e^{-cM})\sum_i\Psi(l_i)\),
since these coordinates belong to bad blocks. The conversion
uses \(|l_i|+l_i^2\le C_\delta\Psi(l_i)\) on \(|l_i|>\delta\).
The field Taylor orders at least six contribute
\(C_D|\Lambda_X|a^6\) at any field by analytic summability.
Now choose \(\theta\) small and \(M\) large. Since
\(|B|a^4\asymp1\), the normalized additive terms are bounded by
\(|X|\) for large \(R\), proving \eqref{eq:transformed-link-stability}.

The real weight has Taylor damping
\(\exp(C|X|-c\sum\Psi(u_i))\). This follows by the field-split
proof of Corollary~\ref{cor:stablerestr}, using the transformed
derivative rows and real coercivity; here there is no quadratic
Taylor term. Applying \eqref{eq:transformed-link-stability} in
the coefficient algebra shows that the same bound, with a
smaller damping constant, holds jointly in the links. For the
pure quartic \(V_X=\sum u_i^4/12\) it holds directly, and hence
it holds for exponentials of all convex interpolations between
\(V_X\) and \(\mathcal U_X(t+rv,\cdot)\). Constants are uniform
for the interpolation parameter in \([0,1]\).

Write the difference to the local quartic reference by the exact
identity
\[
 e^{-\mathcal U_X}-e^{-V_X}
  =(V_X-\mathcal U_X)\int_0^1e^{-(1-q)V_X-q\mathcal U_X}\,dq .
\]
At field Taylor radius \(a\), the first factor has joint norm
at most \(C_D\sum_i[g(|u_i|+a)^4+(|u_i|+a)^6]\). This follows by
subtracting its order-\(g\) quartic and integrating its sixth
derivative for the remaining lower Taylor orders, with the higher
orders again bounded by joint analyticity. On each block bound
these terms by the corresponding polynomial in
\(1+\sup_B|u|/a\) times \(|B|(ga^4+a^6)\lesssim_M g^{1/2}\).
The polynomial is absorbed by spare Gaussian damping and the
affine-quotient regulator, since \(\ell'\le a\) for large \(R\).
At radius \(\ell'\), the exponentials obey the same joint
stability by monotonicity, since \(\ell'\le a\). The first
factor is bounded by
\[
 C_D\sum_i[g(|u_i|+\ell')^4+(|u_i|+\ell')^6]
 \le C_{D,M}A_{\rm sc}^6g^2\sum_{B\in X}(1+\sup_B|u|/\ell')^6.
\]
The growing regulator satisfies
\(\log G_B(u)\ge c(\sup_B|u|/\ell')^2\), so any fixed
fraction of its exponent absorbs these polynomials, leaving
the other fractions for the forest extraction and reweighting.

The forest formula for \(e^{-\mathcal U_X}\) factors across forest
components because \(B_t\), the reference pressure, and its Hessian
all factor there. For each spanning tree extract the square-free
coefficient in the normalized link variables \(v\), paying
\(D^{-(|X|-1)}\), and then multiply by \(\prod_e r_e^{-1}\)
for the original link derivatives. At zero field all these
radii equal one. At a real field the extra factors are bounded by
\[
 \theta^{-C\#\{\text{bad blocks}\}}
             \prod_{B\in X}(1+\sup_B|u|)^{\Delta}.
\]
The powers of the suprema are absorbed per block at a fixed
exponential cost, using spare damping and the affine-quotient
regulator at radius \(a\), or a further fraction of the growing
regulator at radius \(\ell'\). For the constant factor on a
bad block at radius \(a\), use
\[
 \frac{\sup_B|u|}{a}\le C\left(a\|u_B\|_2
              +\frac{\ell'}a\sqrt{\log\widetilde G_B}\right).
\]
A bad block therefore supplies at least \(c\delta^2/a^2\) in the
sum of the Gaussian and quotient exponents, which absorbs the
fixed power of \(\theta^{-1}\) for large \(R\). There is still a
fraction of both exponents left. At radius \(\ell'\) the bad
block instead gives \(\log G_B\ge c\delta^2/(\ell')^2\);
the same absorption follows. These estimates also absorb up
to two source factors when the radius times their source weight
is bounded. The source factors are independent of the links and
factor over the forest components.

On a singleton let \(\delta A_B\) be the difference of the two
weights above. For a larger connected set let \(\delta A_X\) be
its exact tree activity, which is unchanged by subtracting the
link-independent quartic weight under a nonempty spanning-tree
derivative. Define \(K(Y)=\sum_{\pi\in\mathcal P(Y)}\prod_{C\in\pi}
\delta A_C\), with zero contribution from disconnected parts.
Expanding the exact forest partition by the local singleton
reference and the remaining activities gives \((I_4\circ K)(X)\).
For the bounds, retain the factor \(D^{-(|C|-1)}\) for every
connected part. At large field the product of the other small
factors \((C_Dg^{1/2})^k\) is at most
\(C_Dg^{1/2}D^{-(k-1)}\) for large \(R\); at radius \(\ell'\) use
\((C_DA_{\rm sc}^6g^2)^k\) in its place. The remaining exponential
and the count of connected partitions cost a fixed exponential
in \(|Y|\), which is absorbed in the original choice of \(D\).
This proves the two bounds.

Finally, the local mass, nearest-neighbor kinetic term, and
field-linear sources at full links are block-multiplicative
factors. Use the local kinetic density \(u_i(-\Delta_{\rm nn}u)_i\)
at site \(i\), so the reference is in Park's \(\mathcal V^{(0)}\)
coordinates. Multiplying by their block product replaces both
\(I_4(B)\) and \(K(Y)\) by the corresponding reweighted factors,
with the exact same circle identity. The mass costs
\(C|r|\sum_i(|u_i|+a)^2\) in the large-field norm and is
absorbed because \(|r|\ll a^2\). For the kinetic term the
two differences in \(-\Delta_{\rm nn}u\) act on the field or on
the corresponding test coordinate. The \(T\) seminorm at
radius \(a\) therefore costs at most
\(C|z|R^2 b^2a^2(1+N_{B^\square}(u;a))^2\) per block. Its
coefficient tends to zero for \(|z|=O(g)\), and the affine
bound with the block neighborhood absorbs the field factor
using the remaining damping and quotient slack. At radius
\(\ell'\), the kinetic estimate similarly costs
\[
 C|z|R^2 b^2(\ell')^2(1+N_{B^\square}(u;\ell'))^2
 \le C_M A_{\rm sc}^2g(1+\log G_B(u)).
\]
The mass costs the same order, from
\(C|r|b^4(\ell')^2(1+\sup_B|u|/\ell')^2\).
Since \(A_{\rm sc}^2g\to0\), both are absorbed in the remaining growing
regulator. The source factors have the previously stated
product bound. This proves the reweighted assertion.
\end{proof}

At natural fluctuation scale \(A_{\rm sc}=1\) the seed has order
\(g^2\). The rescaling and transient below put this exact seed in the
domain of the iterative map.


\section{Polymer norm and the enlarged analytic ball}\label{sec:contraction}
Use the test-function and source Taylor seminorms of
\cite[Sections 2.2--2.6]{ParkMap}, with the four sectors
\(1,\sigma_a,\sigma_c,\sigma_a\sigma_c\), and its field and source
weights. For block side \(s_j\), the bulk radii are
\(\ell_j=\ell_0s_j^{-1}\) and \(h_j=k_0G^{-1/4}s_j^{-1}\).
If \(k\) is the coalescence scale, the fluctuation source weights
are \(\ell_{\sigma,j}=G s_{j\wedge k}2^{(j-k)_+}\) and
\(\ell_{\sigma\sigma,j}=\ell_{\sigma,j}^2\).
The large-field single-source weight satisfies
\(h_{\sigma,j}h_j=k_0\), and its mixed weight is no larger than
\(h_{\sigma,j}^2\). Use the latter mixed weight precisely as in
(2.20) of that reference when applying the map estimates.
For connected \(X\), set
\(\mathcal A_j(X)=\varrho^{(|X|_j-16)_+}\), and put
\[
 \mathcal G_j(X;\ell)=G_j(X),\qquad
 \mathcal G_j(X;h)=H_j(X)\widetilde G_j(X).
\]
Here \(G_j(X)\) denotes the growing field regulator, not a coupling.
Define
\[
 \|K\|_{\mathcal W_j}
 =\sum_{\mathfrak h\in\{\ell,h\}}\omega(\mathfrak h)
  \sup_{\substack{X\ {
  \rm connected}\\\phi}}
  \frac{\|K(X)\|_{T_j(\phi;\mathfrak h)}}
       {\mathcal A_j(X)\mathcal G_j(X;\mathfrak h)},
 \qquad \omega(\ell)=1,\quad\omega(h)=G^{9/4}.
\]
Strengthened large-set and field-regulator powers are used in
intermediate maps as in the cited construction. Activities on
disconnected sets are the products on their components; their
empty-set value is one. Bounds concern nonempty components.
The condition on the mixed source sector before coalescence applies
only to small connected polymers, of at most 16 blocks. Large
polymers may contain both marks. All analytic activity variations
belong to the same linear support, parity and symmetry subspace.


\begin{lemma}[Enlarged activity ball]\label{lem:enlarged}
Fix the scale, covariance and source-norm hypotheses in dimension
four, with uniform constants, and work in the closed linear subspace
with the prescribed support, symmetry, parity and pre-coalescence
conditions.  The norm weights use a fixed positive reference coupling
\(G\), also at complex potential inputs.  Suppose that the exact
finite-volume maps \(\Psi=K_+\) and
\(R=V_+-\Phi^{\rm pt}(V)\) have a holomorphic extension to
\[
 \|V-V_0\|<cG,\qquad \|K\|_{\mathcal W}<cG^{9/4},
\]
with constants independent of scale, spread radius, volume and mass
cutoff, and that on this \emph{whole} joint neighborhood
\[
 \|\Psi(V,K)\|\le C G^{9/4},\qquad
 \|\Psi(V,0)\|+\|R(V,0)\|\le C\chi^{3/2}G^3,
 \qquad \|D_KR(V,K)\|\le C .
\]
Suppose also, at all real potential inputs in a fixed smaller domain,
that \(\|D_K\Psi(V,0)\|\le\kappa_0<1\).  Then, for
\(0<\epsilon<c/4\), on fixed smaller potential neighborhoods and
\(\|K\|_{\mathcal W}\le\epsilon G^{9/4}\),
\begin{align*}
 \|D_K\Psi(V,K)\|&\le\kappa_0+C\epsilon
                  &&\text{for real }V,\\
 \|\Psi(V,K)\|+\|R(V,K)\|
       &\le C\bigl(\chi^{3/2}G^3+\|K\|\bigr),\\
 \|D_V\Psi(V,K)\|+\|D_VR(V,K)\|
       &\le C G^{-1}\bigl(\chi^{3/2}G^3+\|K\|\bigr).
\end{align*}
The last two estimates hold also for complex \(V\) on smaller
neighborhoods.  The hypotheses on the whole enlarged neighborhood
are independent inputs: bounds or contraction only at \(K=0\)
do not imply them.  Their verification for the present map is given
in Proposition~\ref{prop:map_matching} below.
\end{lemma}

\begin{proof}
Banach-space Cauchy estimates on nested balls of radii comparable
to \(r=cG^{9/4}\) give
\[
 \|D_K\Psi\|\le C,
 \qquad \|D_K^2\Psi\|\le C G^{-9/4}.
\]
For the second estimate apply the one-variable Cauchy estimate twice
in two unit directions, keeping both discs inside the original ball.
The constants depend only on the fixed losses in radius.  Integrating
\(D_K^2\Psi(V,tK)[K,\cdot]\) from zero to one proves the first
assertion, including its uniform constant.  Integrating \(D_K\Psi\)
and \(D_KR\) along the same segment proves the second assertion,
using the stated bounds at zero.  Apply Cauchy's estimate in a
potential direction on a disc of radius \(c'G\), keeping \(K\)
fixed and using this second bound throughout the disc.  This proves
the third assertion.  No analytic continuation from a single real
activity, or from its derivative there, is used.
\end{proof}

For fixed real potential inputs in the stated domains, suppose
\(G_{j+1}/G_j\ge1-CG_j\), \(G_{j+1}\le G_j\), and
\(\chi_j/\chi_{j+1}\le\Omega\).  Choose the scale factor first so
that \(\kappa_0<\Omega^{-3/2}\), then choose
\[
 \kappa_0<\kappa<\vartheta<\Omega^{-3/2}.
\]
Choose \(\epsilon\) so that \(\kappa_0+C\epsilon<\kappa\),
and only afterwards make \(G_0\) small.  Thus
\[
 \|K_{j+1}\|\le C\chi_j^{3/2}G_j^3+\kappa\|K_j\|,
 \qquad
 \|K_j\|\le C'\chi_j^{3/2}G_j^3+\vartheta^j\|K_0\|.
\]
Indeed, the ratio of consecutive weights \(\chi_j^{3/2}G_j^3\)
is at least \(\Omega^{-3/2}(1-CG_0)^3>\vartheta\), so the
inhomogeneous sum is geometric after division by the output weight.
If also \(G_j\ge cG_0/(1+G_0j)\), then
\[
 \sup_j\frac{\vartheta^j\|K_0\|}{G_j^{9/4}}
 \le C\frac{\|K_0\|}{G_0^{9/4}}
       \sup_{j\ge0}\vartheta^j(1+G_0j)^{9/4}.
\]
The last supremum is bounded uniformly for small \(G_0\).  Hence
\(\|K_0\|=o(G_0^{9/4})\) makes both terms uniformly
\(o(G_j^{9/4})\), closing induction in the enlarged ball.
If \(\chi_j\ge\Omega^{-j}\) and
\(\|K_0\|/G_0^3\) grows at most as a fixed inverse power of
\(g=R^{-4}\), the transient is bounded by the first term after
\(O(|\log g|)\) steps: divide by \(\chi_j^{3/2}G_j^3\)
and use \(\vartheta\Omega^{3/2}<1\).  This estimate is uniform
in the mass cutoff.


\section{Finite covariance order and the symmetry class at interior scales}

Take the physical block side to be \(L^{j_0+j}\), where
\(L^{j_0}\in[MR,LMR]\), and start at \(j=0\). Put all previous
covariance pieces equal to zero. These are allowed degenerate
Gaussian covariances. Thus the cumulative covariance at the initial
scale is zero, so both \(W\) and the BBS normal-form change are
zero there. Choose \(L>96\) and then \(M\) large enough that each positive piece
assigned to the step from \(j\) to \(j+1\) has graph range at
most \(L^{j_0+j+1}/4\). This leaves the fixed collar of
fine-block neighborhoods needed for the independence of
disconnected coarse components when \(L\) is large. The
hypercubic and translation symmetries
hold for the numerator, both denominator kernels, the
symmetrized moment replacement, and the Fourier truncations.
They are preserved by every polynomial, positive product, and
periodization used in the decomposition.

At physical scale \(s=j_0+j\) the derivative estimates obtained
from the box decomposition, after normalizing the field, read
\[
 \|\nabla_x^\alpha\nabla_y^\gamma\Gamma_{s+1}\|_\infty
 \le C_{L,M,p}\lambda L^{-(2+|\alpha|+|\gamma|)s}
                   (1+m^2L^{2j})^{-p}
 \qquad (|\alpha|+|\gamma|\le 32).
\]
The factor \(M\) can be chosen fixed. The Park reference mass
is comparable to \(m^2 R^{-2}\), so the corresponding cutoff
indices differ by at most a fixed number. The finite-order version
of the interior-scale estimates uses exactly the displayed bound
for a fixed sufficiently large \(p\). Here is the check of the
covariance inputs.
\begin{itemize}
\item The algebra in Park Sections~4--5 uses translation and
hypercubic symmetries, positive Gaussian integration, finite
range for factorization, and Wick contractions of polynomials.
The latter involve only differences in the field monomials
in \(\mathcal V\); in four dimensions their order is at most
two in each monomial. The same applies to Lemma~4.7 and the
moment estimates for the BBS normal-form coefficients, which
use two spatial differences at most.
\item The bilinear estimates in Park Section~7.2 use the
\(\Phi_j\) norm of the covariance in (7.37), and at most
the mixed derivatives of total order \(p_\Phi=12\) from
(2.10). The Gaussian moments in (7.123) and the product
estimates in Sections~9--10 use Lemma~C.2. More precisely,
the Sobolev step (B.2) bounds the field seminorm by square
integrals of differences of order at most \(p_\Phi+d=16\).
The Gaussian quadratic form bound (C.5) therefore needs
at most 16 differences in each covariance variable.
Taking powers of these regulator bounds increases the
constants in (C.5) but does not increase the derivative order.
\item The \(H\) expectation uses only the covariance operator
norm, bounded by its absolute row sum in (C.14). This follows
from the order-zero estimate and finite range. For the last
direct Gaussian expectation there are only \(L^4\) blocks,
and the same Sobolev quadratic form estimate follows from
the periodized tail derivative bound; no factorization is
used at that step.
\end{itemize}
This is the finite collection of covariance estimates summarized
in Park (1.44) and in the sentence following it. No
differentiability in the covariance parameters is used in
these map bounds. Continuity suffices for the Gaussian
integrals, by the uniform regulator bound and positive
Gaussian coupling, and for the finite Wick polynomials.
The factor \(\lambda\to0\) makes the fixed powers of \(L,M\)
in these upper estimates harmless in the fluctuation
smallness conditions.

For spatial symmetry, let \(\mathsf G_s\) be the subgroup of
lattice isometries preserving the scale-\(s\) partition.
Use \(\mathsf G_s\)-equivariance for the bulk activity, and
the usual field parity in all source sectors. This is sufficient
also on the torus. The only uses of the activity symmetry
in the estimates are as follows. In Lemma~4.15, the localization
sum uses small polymers and the equivariance in (4.67) is
explicitly restricted to \(F(\mathcal B)=\mathcal B\).
Localization assigns the same monomial coefficients to all
sites of its domain. Thus the sum on each block has a constant
site density, translations by a whole block make its
coefficients identical on all blocks, and the signed
coordinate permutations and grid-preserving reflections give
the hypercubic constraints on those coefficients. This yields
the usual translation-invariant effective potential.
In Lemma~3.7 the required extra symmetry is exactly the
vanishing of the odd-field derivatives in (3.28).
Proposition~3.8 invokes it on small polymers only. If a
spatial isometry maps a small polymer to another polymer,
it preserves the grid: for \(L>2^4\), the projection of
a small polymer on each coordinate circle is a proper
nonempty union of whole block intervals, so a boundary
endpoint in each coordinate fixes the translation residue,
taking account of the reversal of endpoints under reflection.
There is no use of invariance under a shift that preserves
only a large wrapped set. In Section~10 the small-set
linearization is Lemma~10.5, invoking Proposition~3.8, and
the large-set linearization is Lemma~10.6, which instead
uses only \(1-\operatorname{Loc}_Y=1\) and the geometry
of the weights. The rough estimates in Sections~7--9
apply unchanged.

Maps~1--5 are equivariant under the next-scale grid
group, which is a subgroup of the input grid group.
Indeed the products, component decompositions, small-set
neighborhoods, localization, Gaussian expectation,
and the covariance redistribution are all equivariant
under this group. The forest initial activities are
equivariant under the initial grid group for the same
reason. The argument applies to connected and disconnected
activities, with the latter defined by the component
product. Field parity, support and the source grading
are preserved throughout. In particular, a small polymer's
neighborhood has diameter less than \(3\cdot2^4 L^s\),
so a two-source activity there vanishes before coalescence.

One may make the final chart change by a componentwise
circle product without using a singleton cancellation.
In \(d=4\) the perturbative output \(U\) and its canonical
potential \(V_+\) differ by the vacuum and mixed constants
and by a quadratic summation-by-parts density. Let
\[
 D(B)=e^{-V_+(B)}\bigl(W_{\rm pt}(B)-W_+(B)\bigr),\qquad
 K_+=D\circ_+\bigl(e^{U-V_+}K_{(5)}\bigr)
\]
on connected polymers, with the corresponding component
extension. Pull out the field-independent exponent using
\(U(\Lambda)-V_+(\Lambda)=\delta u(\Lambda)\).
The remaining equality is Park Lemma~E.1,
since the two block reference factors differ by exactly
\(D\). This also keeps the full \(D(B)\) on a singleton.

For this change, let \(\tau\) be the extended \(K\) radius,
\(G\) the reference coupling, and \(\chi\) the mass factor.
The redundant stiffness coefficient in \(U\) is
\(O_L(\chi G^2+\tau+\|K\|)\): the input potential is
canonical and its Gaussian expectation creates no such
coefficient, while \(Q\) is linear in \(K\) and the
perturbative correction is quadratic with covariance
factor \(\chi\). The local summation-by-parts polynomial
therefore has \(h\)-norm at zero of size
\(O_L(G^{-1/2}(\chi G^2+\tau+\|K\|))\).
Its exponential is a bounded multiplier from the
Map~5 regulator with exponent \(1/2\) to that with exponent
\(3/4\). In fact the block Sobolev bound (B.1) bounds
its field contribution by this coefficient times
\(C(1+\| \varphi/h_\varnothing\|_{L^2(B)}^2+
\|\varphi\|_{h,\widetilde\Phi(B^\square)}^2)\).
This is absorbed by the spare \(H\) and \(\widetilde G\)
powers. The \(\ell\) part can be checked at zero and
then in the equivalent mixed norm \(\mathcal Y\).
The mixed constant of \(Q\) has bounded \(h\)-norm
\(O_L(G^{-9/4}(\|K\|+\tau))\) by applying localization
directly at \(h\). Its exponential is a uniformly
bounded field-independent nilpotent factor.
The one-source linear coefficient is also bounded directly
at \(h\), where \(h_\sigma h_\varnothing\) is constant.
No factor \((h_\varnothing/\ell_\varnothing)^4\) is
used to compare source sectors.

Since \(W\) ignores field-independent terms and is
bilinear, Lemmas~7.3--7.4 bound \(D(B)\), after
polynomial absorption, by
\(O_L(\chi G(\chi G^2+\tau+\|K\|))\) at \(\ell\)
and \(O_L(\chi G^{-1/2}(\chi G^2+\tau+\|K\|))\)
at \(h\). These sizes are admissible in the weighted output
norm. To keep the precise large-set margin, use the
following count on a connected \(X\). Set
\(a=1+\xi/10\), \(a'=1+\xi/16\), and \(N_s=2^4\).
After multiplication by \(e^{U-V_+}\) the activity
on a component \(C\) has unweighted size at most
\(\delta A^a(C)\) times its regulator. The cost of the
multiplier is at most \(2^{|C|}\), absorbed by the
margin from \(1+\xi/8\) to \(a\) and a dimension-dependent
constant. Write \(d_1\) for the unweighted singleton
bound on \(D\). Both \(\delta\le\varrho^{aN_s}\) and
\(d_1\le\varrho^a\) hold by the choices of the extended
radius and small coupling. For \(Y\subseteq X\) with
\(m\ge1\) components and \(q=|X\setminus Y|\), the
unweighted bound on its convolution term is
\[
 \delta^m d_1^q A^a(Y)
 \le \delta\,\varrho^{a(|X|-N_s)_+},
\]
since \(\sum_{C\in\operatorname{Comp}(Y)}(|C|-N_s)_+
  +N_s(m-1)+q\ge (|X|-N_s)_+\).
For \(Y=\varnothing\) the same bound uses \(d_1\) in
place of \(\delta\). The at most \(2^{|X|}\) terms are
absorbed by \(a-a'\), with a dimension-dependent
constant. In the derivative in \(K_{(5)}\), one of the
\(\delta\)'s is replaced by the variation norm and the
number of choices is at most \(|X|\); the same margin
absorbs this factor as well. This proves the rough
output estimate. For the derivative in the intermediate
activity at original input \(K=0\), separate \(Y=X\).
There the derivative is just multiplication by
\(e^{(U-V_+)(X)}\). At \(K=0\) its block norm is at most
\((1+\epsilon_L(G))\mathcal G^{(3/4)}/\mathcal G^{(1/2)}\)
with \(\epsilon_L(G)\to0\). Indeed the vacuum and mixed
coefficients tend to zero in both norms, and the
field-dependent quadratic coefficient uses the spare
regulator power by the same block Sobolev bound. For a
connected polymer of size \(n\), the multiplier in the
output norm therefore costs at most
\[
 (1+\epsilon_L(G))^n
    \varrho^{(1+\xi/8-a')(n-N_s)_+}
 \le (1+\epsilon_L(G))^{N_s}.
\]
Every other term of the intermediate derivative has
at least one \(D\) block. Extract its size \(d_1=o_G(1)\)
and use the preceding count on the remaining factors.
This loses at most one factor \(\varrho^{-a}\),
independent of \(X\) and \(G\), since the exponent in
the count decreases by at most one. The sum of all
these terms is \(o_G(1)\) in operator norm for fixed
\(L,\varrho\). Thus the intermediate derivative has
norm at most \(1+o_G(1)\) in the stronger output norm.
With \(K_{(5)}\) held
fixed at its \(K=0\) value, its partial derivative in the
original \(K\) is \(O_L(G^{3/4})\) by Cauchy on
\(\tau=c_LG^{9/4}\), because the full weighted output
on that disc is
\(O_L(\chi^{3/2}G^3+\chi G\tau)\).
Use Park (10.74) through \(i=5\), before the last change,
then the chain rule. This preserves the arbitrarily
small contraction at \(K=0\), with \(L\) chosen large
and \(G\) small. The enlarged-domain Cauchy argument
then transfers it to \(o(G^{9/4})\) activities.

\begin{proposition}[Interior map at large spread radius]
\label{prop:map_matching}
Use the covariance of Proposition~\ref{prop:fc_newcovbound}
on the finite tori and the initial activity of
Proposition~\ref{prop:transformed-init}. Choose
\(1/12<a_*<1/10\), set \(A_{\rm sc}=g^{-a_*}\),
\(\lambda=A_{\rm sc}^{-2}\), and let
\(G_0\asymp gA_{\rm sc}^4\) denote the initial quartic
coefficient in the field \(\psi=Ru/A_{\rm sc}\).
There are constants \(\epsilon>0\), \(\kappa<1/8\), and \(C\),
independent of the spread radius, volume and mass cutoff, for which
the exact graded RG map obeys, on \(\|K\|_{\mathcal W_j}\le
\epsilon G^{9/4}\),
\[
 \|D_K\Phi^K_{j+1}\|\le\kappa,\qquad
 \|\Phi^K_{j+1}(V,K)\|+|R_j(V,K)|
 \le C(\chi_j^{3/2}G^3+\|K\|_{\mathcal W_j}).
\]
Here \(R_j\) is the bulk remainder in the normal-form coordinates;
the source estimates use the four-sector norm specified above.
This interior-scale map satisfies
the bounds in Proposition~\ref{prop:scaled_flow},
with constants uniform in the spread radius, torus size and
mass cutoff.
\end{proposition}

\begin{proof}
First check the initialization in the precise Park norms.
In physical \(u\) units the Park radii on the initial block
are
\(\ell_\varnothing=\ell_0A_{\rm sc}/(R L^{j_0})\)
and \(h_\varnothing\asymp k_0L^{-j_0}\).
Choose \(M\) large compared to \(\ell_0\) and \(k_0\) small.
The radii in Proposition~\ref{prop:transformed-init}
dominate these. Its affine regulator uses a physical
small-field radius no smaller than that of Park, so is
bounded by the corresponding Park regulator. Its
\(H\) coefficient can be any sufficiently large fixed
constant: by
\(\Psi(u)\ge c_\kappa L^{-2j_0}u^2-C_\kappa L^{-4j_0}\),
one may reserve the necessary \(H\) power in the real
damping and absorb the fixed loss per block into the
choice of the forest coefficient radius. This matches
\(H_{\rm Park}\) also after the change of field amplitude.
The small-field source weight times its field radius is
\(O(G_0)\), and the large-field product is \(O(1)\).
The two-source weight is no larger than the product
of the one-source weights. Thus the source part of the
same initial bound holds for one-source and two-source
flows. The absolute Taylor norm used in the initialization
majorizes the \(T\) norm by its derivative test bounds.
The kinetic factor is treated by the derivative-sensitive
estimate at the end of that proposition. Choosing the
forest rate smaller than the Park large-set rate and
using \(\mathcal Y\)-to-\(\mathcal W\) equivalence yields
\[
 U:=C\left(g^2A_{\rm sc}^6+G_0^{9/4}g^{1/2}\right),
 \qquad \|K_0\|_{\mathcal W}\le U .
\]
Here \(U=o(G_0^{9/4})\), \(U=o(\lambda G_0)\), and
\(G_0=o(\lambda^3)\) by the choice of \(a_*\).

We give a source-norm detail for the application of the
rough map bounds on the enlarged ball. Use the extended
Taylor radii \(\tau=c_LG^{9/4}\) in \(K\) and \(c_LG\)
in the bulk potential. The space of extended directions
uses the same grid symmetry, parity and mixed support
conditions as its real centers. In (9.12)--(9.14) of Park,
localize \(K/I\) separately in each radius using
Proposition~3.4. This gives \(O_L(\tau)\) at \(\ell\)
and \(O_L(\omega(h)^{-1}\tau)\) at \(h\). The latter
is uniformly as small as required by E.7 on choosing
\(c_L\) small. The stable product \(I^X J_B(X)\)
obeys the same estimates, using the polynomial bound and
absorbing its field polynomial in the quartic damping.
This suffices for Map~1. In Map~2,
\[
 e^{-V}-e^{-(V-Q)}=-\int_0^1 e^{-(V-rQ)}Q\,dr .
\]
The sign plays no role in the bound. At both radii,
the integrand is bounded by the corresponding stable
factor times the localized coefficient estimate just given.
The difference of the two \(W\)'s is bilinear and has
at least one \(Q\). The scalar source coefficient in
\(Q\) has absolute size \(O_L(\omega(h)^{-1}\tau)\),
since \(h_\sigma h_\varnothing\) is constant. The
mixed constant is field-independent and does not enter
\(W\). This gives the inputs for the circle estimate
in Lemma~9.5.

In Map~3 the field-independent part of \(\widehat V\)
drops out of
\(\theta\widehat V-\mathbb E_+\theta\widehat V\)
and may be factored exactly from the exponentials.
The other source terms are linear and their fluctuation
derivatives gain \(\ell_\varnothing/h_\varnothing\) at \(h\).
The full mixed covariance contribution at \(h\) is at
most
\[
 h_\sigma^2|\lambda_a\lambda_b|\,|\Gamma_+(a,b)|
   \le C_L\chi G^{1/2}.
\]
The corresponding \(\ell\) bound is \(C_L\chi G^2\).
The vacuum and mixed constants of the perturbative
correction thus satisfy the same bounds as its other
sectors. This gives the bounds of Lemmas~9.10--9.13 on
the one-block difference, of size
\(O_L(\chi^{1/2}G)\) at \(\ell\) and
\(O_L(\chi^{1/2}G^{1/4})\) at \(h\).
For the Gaussian estimate on products of segment
bounds, each block factor uses its own scalar multiple
of the same full fluctuation field. Pointwise,
\[
 H_B(\varphi+r\zeta)\le H_B(\varphi)^{1/2}
   e^{\kappa L^{-4s}h_\varnothing^{-2}\sum_B\zeta_x^2},
 \qquad
 \widetilde G_B(\varphi+r\zeta)
 \le \widetilde G_B(\varphi)^2\widetilde G_B(\zeta)^2 .
\]
These bounds hold uniformly in \(r\in[0,1]\), as does
the identical growing-regulator bound for \(G_B\).
The extra Gaussian factor and fixed fluctuation
polynomials are bounded by a fixed power of \(G_B(\zeta)\)
times a constant. Lemma~C.2 gives the product estimate
of Lemma~9.7 with the suprema inside the expectation.
The covariance reapportionment in Map~4 cancels the
one- and two-block second-order terms as in Lemma~5.5.
Using the printed covariance of \(V\), instead of
\(\widehat V\), changes this cancellation by at most
\(C_L\chi G\tau\) at \(\ell\) and
\(C_L\chi G^{1/2}\omega(h)^{-1}\tau\) at \(h\), by
bilinearity, centering, and the direct bound on \(Q\).
Both fit the input size of Lemma~9.17.
Map~5 and the corrected final change use the same
circle estimate as above. At \(K=0\), the rough output
is \(O_L(\chi^{3/2}G^3)\); on the enlarged complex ball
it is \(O_L(\chi^{3/2}G^3+\tau)\).
These are bounds on the full joint complex neighborhood, with
constants depending only on the fixed scale and regulator choices
and the stated finite collection of covariance upper bounds. In
particular no lower bound on the covariance amplitude, and no
negative power of \(\lambda\), is used in these analytic estimates.
Here is the holomorphy check, including the order of constants.
On the extended ball the extracted \(Q\) has \(\ell\)-norm
\(O_L(c_LG^{9/4})\) and \(h\)-norm \(O_L(c_L)\).
The reference quartic has fixed positive normalized \(h\)-size
proportional to \(k_0^4\). Choose \(c_L\) small relative to that
size and to the spare regulator margins. The quartic part of
\(\operatorname{Re}(V-tQ)\) then retains a fixed positive
coefficient for every \(t\in[0,1]\); lower-degree terms are
absorbed by the same stability and regulator estimates. The
source algebra is finite and nilpotent, so bounded source
coefficients add only fixed-degree field polynomials. The body
of every zero-field factor inverted by localization is bounded
away from zero, and its nilpotent inverse is a finite sum.
Finite-volume circle products are finite sums. Gaussian expectation
is a bounded linear map between the indicated regulator spaces,
with a common integrable majorant on smaller complex balls.
These facts prove Fr\'echet holomorphy of the composed maps and
justify differentiation under the Gaussian integrals there. This is the whole-neighborhood
input to Lemma~\ref{lem:enlarged}; it is logically separate from
the derivative contraction at \(K=0\) verified above. Choose the
scale so that this last contraction has the desired strict margin,
then choose the enlarged radius \(\epsilon\), and only then make
\(R\) large. Small \(G\) alone would not absorb the term
\(C\epsilon\) in that lemma.
The bulk polynomial remainder in the original chart
has the explicit formula (8.7), with \(Q\) linear in \(K\);
its norm is bounded by \(C_L\|K\|\), analytically in
the bulk potential. The argument of
Lemma~\ref{lem:enlarged} and the small derivative
through Park Map~5 now give the derivatives and
contraction assumed in Proposition~\ref{prop:scaled_flow}.

For the perturbative bulk part use the scalar Wick formulas and normal-form change in
\cite[Section 3.2, especially (3.16)--(3.26) and (3.33) onwards]{BBSModel},
together with the covariance-moment estimates of
\cite[Sections 3--4]{BBSIII}. These are Wick-polynomial identities
for the finite-range covariance, with the scalar
combinatorial coefficients, and make no use of the
nearest-neighbor resolvent in the computation. For
clarity, the estimates on their coefficients can
be read directly from the covariance. At physical
scale \(s\), with \(w\) the cumulative covariance
before the step, finite range and the slice estimates
give
\[
 \|w\|_1\le C\lambda L^{2s},\qquad
 \sum_x |x|^2|w(x)|\le C\lambda L^{4s},\qquad
 \sum_x |x|^2w(x)^2\le C\lambda^2L^{2s}.
\]
For the last bound sum the pointwise inequality
\(|w(x)|\le C\lambda(1\vee|x|)^{-2}\) over its support,
with the origin-scale bound on the first block.
The unscaled mass coefficient \(\eta'_j\) is a
constant times \(\Gamma_+(0)\); its dimensionless
size is \(O(\lambda\chi)\).
The \(\beta,\omega,\pi\) coefficients use
differences of \(w^2\) or \(w\Delta w\), hence have
size \(O(\lambda^2\chi)\). For \(w\Delta w\)
move \(\Delta\) to \(\Gamma_+\) in the mixed term
by summation by parts, then use the \(\ell^1\)
bound on \(w\). The \(\theta,\xi\) coefficients
have three covariance factors and have size
\(O(\lambda^3\chi)\). For the potentially
logarithmic part of \(\xi\), combine
\(\delta\sum w^3\) with
\(-3\Gamma_+(0)\sum w^2\). The mixed term is then
\(3\sum w(x)^2[\Gamma_+(x)-\Gamma_+(0)]\).
Symmetry and the second differences bound the
bracket by \(C\lambda\chi L^{-4s}|x|^2\);
the last displayed moment bound applies.
The other terms use the same three moments
and the pointwise \(\Gamma_+\) bound.
Positivity and the lower bubble bound come from
Proposition~\ref{prop:fc_bubble}; its cutoff may
differ from Park's by a bounded number of scales,
which changes only fixed constants.

The coordinate transformations in BBS have
quadratic coefficients bounded by the normalized
moments \(L^{-2s}\|w\|_1\) and
\(L^{-4s}\sum |x|^2|w(x)|\), hence \(O(\lambda)\).
Their inverse is analytic on a fixed small
neighborhood by uniform smallness of the quadratic
derivative. The normal-form error is analytic and starts at
degree three. All its coefficients vanish when
the covariance moments entering the perturbative map
and the normal-form changes are zero. The fresh
covariance coefficients decay faster than
\(\chi^{3/2}\) past the mass scale by the slice bounds.
The normalized cumulative moments also do: summing
the slices gives
\[
 L^{-2s}\|w\|_1\lesssim \lambda L^{-2(j-J)_+},\qquad
 L^{-4s}\sum_x |x|^2|w(x)|
   \lesssim\lambda L^{-4(j-J)_+}.
\]
Here one chooses the mass-decay power larger than
two. The coefficients of the perturbative map, the
quadratic changes, and their inverses are analytic
functions of these finitely many normalized moments
and fresh covariance coefficients, on a uniform
small potential disc. Their degree-three-and-higher
part vanishes at zero covariance data, so it is
bounded on that disc by a constant times the maximum
size of these data. This supplies the
\(\chi^{3/2}\) factor on a fixed disc by taking \(L\)
large. Taylor and Cauchy bounds on the smaller
potential discs give the claimed cubic and first
derivative estimates. This establishes the full
normal-form bounds.
\end{proof}

\section{A weighted recursion with a small bubble coefficient}

The following lemma concerns the recursion once the required one-step
bounds have been proved. It separates those bounds from the issue of
uniformity as the bubble coefficient tends to zero.

\begin{proposition}[Uniform critical tuning for the normal-form recursion]\label{prop:scaled_flow}
Let \(0<\lambda\le1\), \(x_0>0\), and \(\chi_j=2^{-(j-J)_+}\) with
\(J\in\mathbb N\cup\{\infty\}\). Suppose
\[
 x_{j+1}=x_j-b_jx_j^2,\qquad
 0\le b_j\le C\lambda^2\chi_j,\qquad
 b_j\ge c\lambda^2\quad(j\le J).
\]
Fix \(d_0=1/4\), and put
\(f_j=\chi_j^{3/2}x_j^3+U d_0^j\).
Suppose the controlled recursion, in normal-form coordinates, is
\begin{align*}
 g_{j+1}&=g_j-b_jg_j^2+e_j^g(v_j,K_j),\\
 z_{j+1}&=z_j-\theta_jg_j^2+e_j^z(v_j,K_j),\\
 \mu_{j+1}&=L^2\mu_j+\eta_jg_j-\xi_jg_j^2
       -\omega_jg_j\mu_j-\pi_jg_jz_j+e_j^\mu(v_j,K_j),\\
 K_{j+1}&=P_j(v_j,K_j),\qquad v_j=(g_j,z_j,\mu_j),
\end{align*}
with
\[
 |\eta_j|\le C\lambda\chi_j,\qquad
 |\theta_j|+|\xi_j|\le C\lambda^3\chi_j,\qquad
 |\omega_j|+|\pi_j|\le C\lambda^2\chi_j.
\]
Assume the following estimates hold on
\(|g_j-x_j|\le\varepsilon x_j\),
\(|z_j|+|\mu_j|\le M\lambda x_j\),
\(\|K_j\|\le M_K f_j\), for arbitrary fixed \(M,M_K\) and all
sufficiently small parameters. The constant in the last two estimates
may depend on these fixed constants:
\begin{align*}
 |e_j(v,K)|&\le C(\chi_j^{3/2}x_j^3+\|K\|),\\
 \|P_j(v,0)\|&\le C\chi_j^{3/2}x_j^3,\qquad
 \|P_j(v,K)-P_j(v,K')\|\le\kappa\|K-K'\|,\quad\kappa<1/8,\\
 |e_j(v,K)-e_j(v',K')|&\le C\big(f_j|v-v'|/x_j+\|K-K'\|\big),\\
 \|P_j(v,K)-P_j(v',K)\|&\le C f_j|v-v'|/x_j .
\end{align*}
Here the constants in the first two lines are uniform over the
indicated coordinate boxes. If
\[
 \lambda\longrightarrow0,\qquad
 \frac{x_0}{\lambda^3}\longrightarrow0,\qquad
 \frac{U}{\lambda x_0}\longrightarrow0,
\]
then for every \(K_0\) of norm at most \(U\) the recursion has a
solution with \(g_0=x_0\), \(z_j,\mu_j\to0\), and
\[
 |g_j-x_j|\le\varepsilon x_j,\qquad
 |z_j|+|\mu_j|\le C\lambda\chi_j x_j,\qquad
 \|K_j\|\le C f_j.
\]
The same statement, with both terminal coordinates set to zero at the
last local scale, holds on a finite horizon with the same constants.
If the inputs are continuous in external parameters in a fixed
Banach-space chart, so are the two initial tuned couplings.
\end{proposition}

\begin{proof}
All constants below are independent of \(J\) and of a finite terminal
scale.  Shrink \(\lambda^2x_0\) so that \(b_jx_j<1/4\).
Telescoping \(x_j^{-1}\) gives
\[
 \frac{x_0}{1+C\lambda^2x_0j}\le x_j\le x_0,
 \qquad
 \sum_{i\ge j}\chi_i x_i^p\le C_p\lambda^{-2}x_j^{p-1}
 \quad(p=2,3).
\]
Up to \(J\), the second estimate follows by telescoping
\(x_i^{p-1}\), using \(b_i\ge c\lambda^2\); after \(J\),
\(x_i\asymp x_{J+1}\) and the remaining sum is geometric.
In particular,
\begin{equation}\label{eq:recursion-tail-estimates}
 \sum_{i\ge j}f_i\le C(\lambda^{-2}x_j^2+Ud_0^j),
 \qquad \frac{Ud_0^j}{x_j}\le C\frac U{x_0}.
\end{equation}
Put \(A_{j,i}=\prod_{i<l<j}(1-2b_lx_l)\) for \(i<j\).
Since \(1-2t\le(1-t)^2\),
\[
 A_{j,i}\le(x_j/x_{i+1})^2,
 \qquad
 \sum_{i<j}A_{j,i}b_ix_i^2\le Cx_j,
 \qquad
 \sum_{i<j}A_{j,i}f_i
 \le Cx_j\left(\frac{x_0}{\lambda^2}+\frac U{x_0}\right).
\]
For the middle estimate use
\(\sum_{i<j}b_i\le x_j^{-1}-x_0^{-1}\).
For the last estimate, before the mass scale the summands from
\(\chi_i^{3/2}x_i^3\), divided by \(x_j^2\), are bounded by
\(Cx_i\).  The inequality
\(\log(x_i/x_{i+1})\ge b_ix_i\) and
\(t\log(x_0/t)\le Cx_0\) give the displayed bound.
After the mass scale use the geometric \(\chi_i\)-sum.  For the
transient, the lower bound on \(x_i\) and
\(\sum_i d_0^i(1+i)^2<\infty\) suffice.  Also,
\begin{equation}\label{eq:recursion-geometric-estimate}
 \sum_{i\ge j}L^{-2(i-j+1)}f_i
 \le C(x_j^3+Ud_0^j),
 \qquad f_j/f_{j+1}\le4.
\end{equation}
The last inequality holds for small parameters because the adjacent
ratio of \(\chi_j^{3/2}x_j^3\) is at least
\(2^{-3/2}(1-C\lambda^2x_0)^3>1/4\), whereas the transient
has ratio exactly \(d_0=1/4\).

Write \(s_j=g_j-x_j\).  On sequences with \(s_0=0\) and the
prescribed \(K_0\), define a fixed-point map by the following Green
formulas.  The output \(s\) solves the forward linear equation with
coefficient \(1-2b_jx_j\) and forcing
\(-b_js_j^2+e_j^g\).  The output \(z\) is
\[
 z_j=\sum_{i\ge j}(\theta_i g_i^2-e_i^z).
\]
Writing
\(F_i^\mu=\eta_i g_i-\xi_i g_i^2-
\omega_i g_i\mu_i-\pi_i g_i z_i+e_i^\mu\), the output mass is
\[
 \mu_j=-\sum_{i\ge j}L^{-2(i-j+1)}F_i^\mu .
\]
All right-hand sides use the input sequence.  The output activity is
\(K_{j+1}=P_j(v_j,K_j)\).  For a finite horizon replace both
backward sums by sums over the local update indices; then the two
terminal coordinates are exactly zero.  Every estimate above remains
valid for these truncated sums.

Use the maximum of the four weighted suprema with denominators
\[
 \varepsilon x_j,\quad M\lambda x_j,\quad M\lambda x_j,
 \quad M_K f_j.
\]
The closed product box is complete; using separate bounds for \(z\)
and \(\mu\) is legitimate by replacing \(M\) in the hypotheses
by \(2M\).  Choose \(M_K\) large enough that the activity box is
preserved: its normalized output is at most
\(4C/M_K+4\kappa<1\).  Choose \(M\) large compared with the
uniform coefficients in the leading \(\theta\) and \(\eta\)
forcings.  These forcings have size at most \(C\lambda x_j\).
All other backward forcings have size \(o(\lambda x_j)\) for
these fixed boxes.  In fact, set
\[
 a=\frac{x_0}{\lambda^2}+\frac U{x_0},\qquad
 b=\frac{x_0}{\lambda^3}+\frac U{\lambda x_0}.
\]
Then \(a=\lambda b\to0\) and \(b\to0\).  The activity and
cubic errors in the \(z\) sum, divided by \(\lambda x_j\),
are \(O(b)\) by \eqref{eq:recursion-tail-estimates}; those in
the mass sum are \(O(x_0^2/\lambda+U/(\lambda x_0))=o(1)\).
The remaining perturbative mass terms are smaller as well.
The normalized forward output is at most
\(C\varepsilon+Ca/\varepsilon\).  Thus a sufficiently small
fixed \(\varepsilon\), followed by small parameters, preserves
all four boxes.

Here are the difference estimates that give uniform contraction.
Let \(D_s,D_z,D_\mu,D_K\) be the four normalized input distances.
Constants in an \(o(1)\) below may depend on the now fixed box
sizes and on \(\varepsilon\), but are independent of \(J\)
and the terminal scale.  The forward Green estimate and the
Lipschitz hypothesis on \(e\) give
\[
 D_s^{\rm out}\le(C\varepsilon+o(1))D_s
                  +o(1)(D_z+D_\mu+D_K).
\]
Indeed the error contributions are bounded by fixed multiples of
\(a\), \(M\lambda a/\varepsilon\), and
\(M_Ka/\varepsilon\); all tend to zero.
For the backward coordinates the leading perturbative dependence on
\(s\) is \(C\varepsilon/M\), whereas the error dependence
uses \(b\to0\).  Hence
\[
 \max(D_z^{\rm out},D_\mu^{\rm out})
 \le(C\varepsilon/M+o(1))D_s
       +o(1)(D_z+D_\mu+D_K).
\]
For example the \(K\)-to-\(z\) coefficient is at most
\(C(M_K/M)b\), which is the potentially delicate backward
sum when the bubble coefficient is small.
Finally, the activity estimates and \(f_j/f_{j+1}\le4\) give
\[
 D_K^{\rm out}\le4\kappa D_K+
 \frac{4C_{M,M_K}}{M_K}
       \bigl(\varepsilon D_s+M\lambda(D_z+D_\mu)\bigr).
\]
Choose \(\varepsilon\) after \(M,M_K\) so that the fixed
terms in all these row sums are strictly less than one; then make
\(\lambda,b\) small.  The maximum row sum is bounded by some
\(q<1\) independent of the mass scale and terminal scale.
This proves the claimed uniform contraction, rather than relying
on a contraction in any unweighted coordinate norm.  Its unique
fixed point in the box solves the recursion.

For \(j>J\), the backward forcing bounds improve by \(\chi_j\).
Here \(x_i\asymp x_j\) for \(i\ge j\),
\(\sum_{i\ge j}\chi_i x_i^p\le C\chi_j x_j^p\), and
\[
 Ud_0^j\le C(U/x_0)\chi_j x_j.
\]
For the last estimate use \(\chi_j\ge2^{-j}\),
\(2d_0<1\), and the lower bound on \(x_j\).
The same statements for \(\chi^{3/2}\) and the geometrically
weighted mass sums follow immediately.  Inserting these bounds in
the Green formulas gives
\( |z_j|+|\mu_j|\le C\lambda\chi_j x_j\).
For \(j\le J\) the previously proved bound suffices.
Consequently both terminal coordinates tend to zero, including
when \(J=\infty\), since then \(x_j\to0\).

For continuity, iterate this uniformly contracting map starting
from the reference bulk sequence and an admissible activity
sequence with the prescribed initial value.  At each fixed index,
finite-coordinate operations are continuous by hypothesis.
The infinite backward mass sum has a geometric majorant.
The backward \(z\) sum has a geometric tail near a fixed finite
mass scale.  As the mass scale tends to infinity, its tail starting
at \(i\) is bounded by
\(C(\lambda x_i+Ud_0^i)\), and for \(J\ge i\)
\(x_i\le(x_0^{-1}+c\lambda^2i)^{-1}\).  This tends to zero
locally uniformly in the external parameters with \(\lambda,x_0>0\).
Thus fixed coordinates of every finite iterate are continuous.
The common contraction constant makes their convergence uniform
at each fixed index, proving continuity of the two initial tuned
couplings.  The same reasoning is immediate on a finite horizon.
\end{proof}


\section{Matching the tuned trajectory to the Ising parameters}
\label{sec:matching-physical}

Use the notation of the exact field change and write
\(A_{\mathrm{sc}}=g^{-a_*}\), with
\(1/12<a_*<1/10\), for the amplitude factor. It is distinct
from the kernel \(A=\beta P\). In the variable
\(\psi=(R/A_{\mathrm{sc}})u\), the initial quartic coefficient,
in the convention \(g_0\psi^4/4\), and the covariance prefactor are
\[
 G_0=\frac{A_{\mathrm{sc}}^4g}{3},\qquad
 \lambda=A_{\mathrm{sc}}^{-2}.
\]
Choose the initial block side \(b\) to be a power of \(L\) with
\(M R\le b<MLR\), using the fixed sufficiently large aspect
from the connected initialization. The physical fluctuation
scale in \(u\) is a fixed multiple of \(A_{\mathrm{sc}} R^{-2}\),
and its large-field scale is a fixed multiple of \(b^{-1}\).
Indeed the large-field scale in \(\psi\) is
\(G_0^{-1/4}b^{-1}\), and
\(A_{\mathrm{sc}}G_0^{-1/4}/R\) is fixed.
Thus Proposition~\ref{prop:transformed-init}, with the large-set
factor chosen for the Park norm, gives the initial bound
\[
 U=C\big(A_{\mathrm{sc}}^6g^2+G_0^{9/4}g^{1/2}\big)
\]
on the full polymer norm. The corresponding smallness ratios are
\begin{align*}
 U/G_0^{9/4}
    &\le C(g^{3a_*-1/4}+g^{1/2})=o(1),\\
 G_0/\lambda^3&=\tfrac13g^{1-10a_*}=o(1),\
 U/(\lambda G_0)&\le
 C(g^{1-4a_*}+g^{7/4-7a_*})=o(1).
\end{align*}
The estimates for the field range and large-set factor also
hold for the one- and two-source initializations. In the
fluctuation sector the source weight times the corresponding
field radius is \(O(G_0)\), and in the large-field sector it
is bounded. One may take the two marked sites far enough apart
that coalescence occurs after the fixed transient.

\begin{proposition}[Physical matching at positive auxiliary mass]
\label{prop:physical-matching}
For sufficiently large \(R\), fix \(0<m^2\le g\) and a sufficiently
large torus of side \(bL^N\). There are parameters
\(\beta_{m,N},z_{m,N}\) in the exact Ising representation such that
\[
 |\beta_{m,N}-(1-m^2)|+|z_{m,N}|\le Cg
\]
and the bulk flow through scale \(N-1\) satisfies the controlled
bounds of Proposition~\ref{prop:scaled_flow}, with terminal
normal-form \(z,\mu\) equal to zero. The constants are uniform
in \(m^2\) and \(N\).
\end{proposition}

\begin{proof}
Lift the exact torus Hessian as in Section~\ref{sec:finitevolume}
and construct the covariance slices for this lift. Use the
prescribed torus periodization for the local steps. Choose
\(b/R\) large enough that the support of each such slice is
less than \(s_{j+1}/4\) at the step from \(j\) to \(j+1\) and its support
does not wrap before taking periodization. Proposition~\ref{prop:map_matching} verifies the exact map and
normal-form hypotheses on these local steps, uniformly in the lift. In particular, Proposition~\ref{prop:fc_bubble}
gives a bubble coefficient comparable to \(\lambda^2\) up to the
mass scale. The fixed multiplicative constant in the mass cutoff
can be chosen to match the lower bound in that proposition.
The finite-range bounds give the required decay afterwards.

For each covariance parameter pair in a physical box, fix
\(g_0=G_0\) and the exact initial activity and solve the finite
horizon recursion of Proposition~\ref{prop:scaled_flow} through
scale \(N-1\). In that proposition truncate both backward
Green sums at the last local update, so \(z_{N-1}=\mu_{N-1}=0\).
The sum estimates and the contraction are unchanged. The two
returned initial couplings \((\mu_0^*,z_0^*)\) obey
\[
 |\mu_0^*|+|z_0^*|\le C\lambda G_0
\]
by the smallness ratios above. They are continuous in the
external parameters at fixed \(N\). For this, iterate the
uniform contraction starting with the reference sequence.
At each finite iteration, the coordinates and the relevant
field jets are continuous by the local finite sums and
Gaussian integrals in the step, using the regulator bounds
to dominate the integrals. The initial activities have this
continuity from their finite exact formulas. At a fixed scale
the contraction makes these iterates converge uniformly in the
parameters for every coordinate and field-jet evaluation
under consideration. This also proves continuity when the
reference couplings, and hence the norms of the activity
spaces, vary with the covariance parameters.

Let \(\widehat E_N(0)\) denote the row sum of the reference
Hessian difference on the torus. The physical residual mass is
\[
 r_N(\beta,m^2)
 =\frac{(1-\beta q)(1-m^2-\beta)}{\beta p}
     -(1-\beta q)^2\widehat E_N(0).
\]
The cavity estimate is uniform for \(\beta\) in a fixed small
complex neighborhood of the real reference interval. This
follows by treating its increment as a single coupling
variable in the same recursion, with coupling change
\((\delta\beta)Q\). Cauchy's estimate therefore gives
\(\partial_\beta\widehat E_N(0)=O(g)\) uniformly. Consequently
\(\partial_\beta r_N=-1+O(g)\) when
\(|\beta-1|+m^2=O(g)\). For fixed \(m^2\), one can
use \(r\) in place of \(\beta\) on an interval
\(|r|\le C_1 g\), taking a large enough constant in the allowed
\(\beta\) interval about \(1-m^2\).

Before the first Gaussian integration the cumulative covariance
is zero, so the quadratic normal-form change of coordinates
is the identity. The bare local mass and kinetic couplings
in that chart are
\[
 \mu_0^{\rm bare}=\frac{b^2}{R^2\lambda}r,
 \qquad
 z_0^{\rm bare}=-\frac{z}{\lambda}.
\]
Here \(z\) is the coefficient used in the covariance absorption
of \(zR^2\mathcal D\), so the residual interaction has
kinetic coefficient \(-zR^2\) in \(u\).
On the physical box \(|r|,|z|\le C_1g\) consider the map
\[
 (r,z)\longmapsto
 \left(\frac{R^2}{b^2}\lambda\mu_0^*,\ -\lambda z_0^*\right),
\]
where the starred quantities are computed at the covariance and
initial activity of that external pair. Its image has size
\(O(\lambda^2G_0)=O(g)\). All constants in this estimate
can be chosen before \(C_1\): the covariance estimates need only
a fixed small upper bound on the weighted Hessian norm and on
\(|z|\), and the large-field initialization absorbs the local
quadratic factors when \(C_1g^{1/2}\) is small. Choose \(C_1\)
large, then take \(R\) sufficiently large. The map is a continuous
self-map of the closed box and has a fixed point by Brouwer's
theorem. At that point both bare local couplings agree with
the tuned initial values, while the quartic and activity
were already matched. This proves the proposition.
\end{proof}


\section{Finite tori, the last covariance, and the physical endpoint}\label{sec:finitevolume}

The exact Hessian absorption can be used on every sufficiently large
torus. Let \(D_N\) be its side and let \(E_N=S_N-S_{0,N}\) be the
periodic high-temperature Hessian difference. Lift the torus convolution
kernel to \(\mathbb Z^4\) by using shortest coordinate representatives
in \([-D_N/2,D_N/2]^4\), splitting the coefficient equally among the
representatives when there is a tie. This lift is hypercubic symmetric,
has the same row sum, and has exponential weighted \(\ell^1\) norm
\(O(R^{-4})\) by the same torus-distance bound for \(E_N\). In
particular the moment matching and positive product decomposition
of Proposition~\ref{prop:fc_newcovbound} apply to the lift with
uniform constants. Periodize the resulting covariance pieces.
Their sum is the exact torus covariance for \(m^2>0\), because
periodization samples the Fourier multipliers at the torus
frequencies. The lift converges to the infinite-volume \(E\) in
a possibly weaker exponential weighted \(\ell^1\) norm. Indeed,
each fixed entry converges by high-temperature locality, while
the stronger exponential bound controls the tails uniformly.
This argument also gives continuity jointly along converging
\(\beta\) and \(z\) at every fixed scale.

The last Gaussian integral can be handled directly when the
torus side is much larger than the mass length. Let its side be
\(D_N=bL^N\) and stop the local RG at scale \(N-1\), with block side
\(D_N/L\). There are then exactly \(L^4\) blocks. Denote the remaining
periodized covariance by \(C_{\rm tail}\). For a piece with spatial
scale \(RL^j\), each periodic point receives at most
\(C(1+RL^j/D_N)^4\) translates. Thus the derivative bounds for the
piece, summed over the tail, give the same bounds at the last block
scale, with arbitrary postmass decay after decreasing the exponent
in the original bound. In particular all the rescaled derivative
bounds of \(C_{\rm tail}\) divided by the square of the fluctuation
field scale tend to zero for each fixed \(m^2>0\) as \(N\to\infty\).
The Gaussian regulator-moment argument of Park's Lemma C.2
applies to this fixed number of blocks. More explicitly, the block
Sobolev inequality bounds the sum of their regulator exponents by
a positive quadratic form \(Q(\zeta)\). The normalized derivative
bound gives \(\mathbb E Q(\zeta)=o(1)\). The largest eigenvalue
of the corresponding Gaussian quadratic form is at most its trace,
so its determinant formula bounds every fixed exponential moment
uniformly for large \(N\). It uses the derivative
bounds rather than finite range at this last integration.
Consequently the shifted Gaussian expectation of any last-scale
regulated activity has its zero-field Taylor seminorm bounded
by a fixed multiple of its regulated input norm.

More precisely, assume the local controlled RG flow through scale
\(N-1\) is constructed with the uniform envelopes of
Proposition~\ref{prop:scaled_flow}, including any indicated
observable sector. Write \(\mathcal Z_{N-1}=I_{N-1}\circ K_{N-1}\)
for the weight after removal of its field-independent factor.
Since there are only \(L^4\) blocks, the local stability estimate
\(\|I_{N-1}-1\|_{\ell,G}=O_L(x_{N-1})\), the polymer bound,
and the last Gaussian bound give
\[
 Z_{\rm tail}(0)=1+O_L(x_{N-1}+f_{N-1}),\qquad
 \|Z_{\rm tail}\|_{T_\ell(0)}\le C_L,
 \qquad Z_{\rm tail}=\mathbb E_{C_{\rm tail}}\theta\mathcal Z_{N-1}.
\]
The first estimate here refers to the bulk component. In particular
that component is bounded below by \(1/2\) for a small initial
coupling. The estimates follow directly by expanding the finite
circle product around the product of unit singleton factors,
using its product bound and the Gaussian moment estimate for
the product of regulators. With two fixed source sites, the
last contribution to the source-pair coefficient of
\(\log Z_{\rm tail}(0)\) is \(O_L(\ell_{N-1,\sigma\sigma}^{-1})\),
which tends to zero at fixed \(m^2>0\) because \(x_{N-1}\) stays
bounded below and the source-pair scale increases exponentially
after coalescence.

For the susceptibility it suffices to use the bulk Taylor bound.
If \(C^\Lambda\) is the full auxiliary covariance of the rescaled
field \(\psi\) and \(C_0^\Lambda=\widehat C^\Lambda(0)\) its row
sum, completing the square in a uniform source gives the exact
identity
\[
 \chi_\psi^\Lambda=C_0^\Lambda+
       (C_0^\Lambda)^2 |\Lambda|^{-1}
       D^2_{\boldsymbol1,\boldsymbol1}\log Z_{\rm tail}(0).
\]
The discarded bulk normalization is independent of the auxiliary
field and cancels in the derivative. The Taylor bound and lower
bound at zero give
\(|\Lambda|^{-1}|D^2_{\boldsymbol1,\boldsymbol1}\log Z_{\rm tail}(0)|
 \le C_L D_N^{-2}\).
By the row-sum calculation in the absorbed covariance,
\[
 C_0^\Lambda
 = \frac{\lambda R^2\beta p}{(1-\beta q)m^2},
 \qquad
 \chi_\psi^\Lambda
 =\lambda R^2\frac{\beta p+(\beta p)^2\chi_\sigma^\Lambda}
                    {(1-\beta q)^2}.
\]
Both identities are exact for the same finite torus parameters.
For the first, the row sum of \(K+zR^2\mathcal D\) is zero and
\(\widehat M_m(0)=m^2\). For the second, conditional on the spins,
the original auxiliary field \(h\) has Gaussian mean \(A\sigma\)
and covariance \(A\). Apply \(S_0\) on both sides of its covariance
and then multiply by \(\lambda R^2\).

We record explicitly the error in these identities, since the tuned
parameters depend on the volume. For each fixed \(R\) and \(m^2>0\),
the preceding Taylor bound is uniform over the compact physical
parameter box. Consequently, at the tuned parameters,
\begin{equation}\label{eq:finite-suscept-error}
 \chi_\sigma^{\Lambda_N}(\beta_{m,N})
 =\frac{1-\beta_{m,N}q}{\beta_{m,N}p\,m^2}
       -\frac1{\beta_{m,N}p}
       +O_{R,m}(D_N^{-2}).
\end{equation}
Indeed the error before converting from \(\psi\) to \(\sigma\)
is bounded by \(C_L(C_0^{\Lambda_N})^2D_N^{-2}\); the conversion
factor is
\((1-\beta_{m,N}q)^2/[\lambda R^2(\beta_{m,N}p)^2]\).
Thus one may bound the displayed error by
\(C_L\lambda R^2(m^2)^{-2}D_N^{-2}\), up to the fixed
field-normalization constants. This estimate is independent of the
choice of the tuned kinetic parameter. In particular the
susceptibilities are bounded at each fixed positive mass, and their
limit along any convergent temperature subsequence is prescribed.

Recall that \(\beta_c^D\)
denotes the critical inverse temperature of the modified kernel
\(D\) used in the Hubbard--Stratonovich representation. The
conversion to the original box kernel is made in the proof of
Theorem~\ref{thm:main}.

Fix \(m^2>0\), and choose a subsequence with
\(\beta_{m,N}\to\beta_m\), possible by the compact physical
parameter box. Let \(C_m<\infty\) bound the susceptibilities in
\eqref{eq:finite-suscept-error} along this subsequence. If
\(\beta<\beta_m\), each fixed finite free box can eventually be
embedded in \(\Lambda_N\) with \(\beta\le\beta_{m,N}\).
The Griffiths inequalities, first for the added bonds and then
for the temperature, bound its partial susceptibility by \(C_m\).
Taking increasing free boxes gives
\(\chi_\sigma^{\mathrm f}(\beta)\le C_m\). Moreover
\[
 \chi_\sigma^{\mathrm f}(\beta_m)
   =\sup_{\beta<\beta_m}\chi_\sigma^{\mathrm f}(\beta).
\]
To justify this equality without assuming uniqueness at \(\beta_m\),
write the free susceptibility as the supremum of finite sums of
finite-free-volume correlations. Each such sum is continuous and
increasing in \(\beta\), and these suprema commute. It follows
that \(\chi_\sigma^{\mathrm f}(\beta_m)<\infty\).

We next obtain a neighborhood with uniform decay; this step is
needed before passing a susceptibility through the thermodynamic
limit. For a finite box \(S\) containing the origin define
\[
 \varphi_\beta(S)=
 \sum_{x\in S,\,y\notin S}
  \tanh(\beta D(x-y))
       \langle\sigma_0\sigma_x\rangle^{\mathrm f}_{S,\beta}.
\]
Since \(D\) has fixed finite range at fixed \(R\),
\(\tanh(\beta_mD)\le\beta_mD\), and
\(\sum_yD(x-y)=1\), finiteness of the free susceptibility
implies \(\varphi_{\beta_m}(S)<1\) for a sufficiently large
box: the sum is bounded by \(\beta_m\) times a tail of that
susceptibility, because only sites within the interaction range
of \(\partial S\) contribute. Finite-volume continuity therefore
gives an interval \(I\) about \(\beta_m\), a number
\(\beta_+>\beta_m\) in \(I\), and \(a<1\) such that
\(\sup_{\beta\in I}\varphi_\beta(S)\le a\).
The finite-set criterion of \cite[Section 2]{DCT} identifies
\(\sup\{\beta:\varphi_\beta(S')<1\text{ for some finite }S'\}\)
with the usual spontaneous-magnetization threshold. Hence
\[
 \beta_m<\beta_+\le\beta_c^D.
\]

The same criterion gives a bound uniform in the periodic volume,
not merely a bound in the infinite-volume free state. Apply the
Simon--Lieb inequality on the torus using translates of the
embedded box \(S\). Each exit from such a box costs a total
weight at most \(a\), and moves the base point by at most a
fixed distance \(r_S\), depending only on \(S\) and the
interaction range. Iterating until the target is reached, and
using that every correlation is at most one, gives constants
\(C,c>0\) for which
\begin{equation}\label{eq:uniform-periodic-decay}
 0\le\langle\sigma_x\sigma_y\rangle_{\Lambda_N,\beta}
   \le C e^{-c\,\operatorname{dist}_{\Lambda_N}(x,y)},
 \qquad \beta\in I,
\end{equation}
for all sufficiently large tori. The constants do not depend on
\(N\) or on \(\beta\in I\). The finite box used inside the
Simon--Lieb inequality has its free correlations, exactly as in
\(\varphi_\beta(S)\); no finite-volume uniqueness assumption is
being made here.

Since \(\beta_m<\beta_c^D\), the infinite-volume Gibbs state at
\(\beta_m\) is unique. Every local subsequential limit of the
periodic laws at the varying temperatures \(\beta_{m,N}\) is a
Gibbs law at \(\beta_m\), by the finite-range DLR equations and
\(\beta_{m,N}\to\beta_m\). Compactness and uniqueness therefore
give convergence of all local correlations along the chosen
subsequence. The uniform exponential bound
\eqref{eq:uniform-periodic-decay} allows their sums to converge:
use shortest torus representatives, sum first over a fixed ball,
and then bound its complement by the summable exponential tail.
Combining this with \eqref{eq:finite-suscept-error} proves
\begin{equation}\label{eq:physical_suscept}
 \chi_\sigma^{\mathrm f}(\beta_m)
   =\frac{1-\beta_m q}{\beta_m p\,m^2}
        -\frac1{\beta_m p}.
\end{equation}

The physical parameter box keeps \(\beta_m p\) bounded above
and away from zero, and \(1-\beta_mq\) bounded away from zero.
Thus \(\chi_\sigma^{\mathrm f}(\beta_m)\ge c/m^2-C\), and
it diverges as \(m^2\downarrow0\). Each \(\beta_m\) lies below
\(\beta_c^D\). If an accumulation point were strictly smaller
than \(\beta_c^D\), choose a fixed temperature strictly between
them and \(\beta_c^D\). Monotonicity and subcritical finiteness
would bound those susceptibilities by a fixed finite number,
a contradiction. Therefore
\begin{equation}\label{eq:tuned-critical-endpoint}
 \beta_m\longrightarrow\beta_c^D\qquad(m^2\downarrow0).
\end{equation}
No monotonic choice of the tuned parameters or uniqueness of the
Brouwer fixed point is needed. Finally, free correlations are
left-continuous in inverse temperature: for each pair, write the
infinite-volume free correlation as the supremum of its
finite-free-volume values, and commute that supremum with
\(\sup_{\beta<\beta_c^D}\). Monotonicity then implies that the
correlations at any of the above \(\beta_m\) tend to the critical
free correlations.


\section{The two-source flow and the massive thermodynamic limit}

Let the block length at the \(j\)-th integration be \(s_j=bL^j\).
Use the reference quartic \(x_j\) of
Proposition~\ref{prop:scaled_flow}, the covariance prefactor
\(\lambda_{\rm cov}=A_{\rm sc}^{-2}\) from the field rescaling,
and the activity envelope
\[
 f_j=\chi_j^{3/2}x_j^3+Ud_0^j,\qquad
 x_j\asymp \frac{x_0}{1+\lambda_{\rm cov}^2x_0\min(j,J)} .
\]
Assume the controlled bulk trajectory and the extended one-step
estimate with the source norms before coalescence. We use the
extended estimate during the transient. Afterwards we use the
ordinary-size activity bound for the same exact map. We use pairs
whose coalescence scale exceeds the transient length \(J_*+2\).
To obtain a uniform \(J_*=O(|\log g|)\), iterate the contraction
inequality for \(K\). At \(K=0\) the forcing is at most
\((C_{\rm RG}/2)\chi_{j+1}^{3/2}x_{j+1}^3\), taking
\(C_{\rm RG}\) large enough for the actual chart change, as
in the proof of Park Proposition~9.1. For small \(x_0\) the consecutive
ratios of \(\chi_j^{3/2}x_j^3\) exceed \(q=1/3\).
Taking \(\kappa<1/8\) gives the bound
\[
 \|K_j\|\le \tfrac45 C_{\rm RG}\chi_j^{3/2}x_j^3+U d_0^j .
\]
The ratio of the last term to \(\chi_j^{3/2}x_j^3\) is at most
\((U/x_0^3)(d_0/q)^j\), and \(U/x_0^3\) is at most a fixed
power of \(1/g\). Thus the trajectory is in the ordinary
Park activity ball after \(J_*\), uniformly in the mass cutoff.
This argument applies in the full source norm as long as the
source couplings stay in a fixed bounded domain, which is
verified below.

Write \(\nu_j,g_j^{\rm old}\) for the mass and quartic in the
original polynomial coordinates. They obey
\[
 s_j^2|\nu_j|\le C\chi_jx_j,\qquad
 s_j^2|\nu_{j+1}-\nu_j-\eta'_jg_j^{\rm old}|
       \le C(\chi_jx_j^2+f_j),
 \qquad \eta'_j=3\Gamma_{j+1}(0).
\]
For the first bound, solve the normal-form mass equation
backward. Using the estimates of Proposition~\ref{prop:scaled_flow},
its forcing is bounded by \(C(\lambda_{\rm cov}\chi_i x_i+f_i)\).
Consequently the geometric backward sum bounds \(|\mu_j|\) by
\(C(\lambda_{\rm cov}\chi_j x_j+\chi_j^{3/2}x_j^3+Ud_0^j)\).
Since \(\chi_j\ge 2^{-j}\) and
\(x_j\ge x_0/(1+C\lambda_{\rm cov}^2x_0j)\), this is
\(O(\chi_jx_j)\). The inverse normal-form mass change is
\(s_j^2\nu_j=\mu_j+O(\mu_j^2)\), by its explicit formula
involving \(s_j^{-2}w_j^{(1)}\). For the second bound, subtract
the linear Wick mass increment in the original perturbative
formula. Each remaining quadratic term contains an increment
with a fresh covariance slice, giving \(O(\chi_jx_j^2)\);
the exact \(R^U\) contributes \(O(f_j)\) in the normalized
mass coordinate. For the only term with a possibly logarithmic
bound, the mass coefficient contains
\(\sum_y[(w_j+\Gamma_{j+1})^3-w_j^3](y)
-3\Gamma_{j+1}(0)\sum_y w_j(y)^2\).
The first term linear in \(\Gamma_{j+1}\) now has
\(\Gamma_{j+1}(y)-\Gamma_{j+1}(0)\), whose derivative bound
pays the extra factor \(|y|/s_j\).
For the slices assume
\[
 \|\Gamma_{j+1}\|_\infty\le C\lambda_{\rm cov}\chi_j s_j^{-2},
 \qquad \left|\sum_y\sum_{i\le j}\Gamma_i(y)\right|
       \le C\lambda_{\rm cov}s_j^2 .
\]
The second bound follows from the first and the support radius
of each slice. The smallness hypotheses in
Proposition~\ref{prop:scaled_flow} imply
\(x_0/\lambda_{\rm cov}^2+U/x_0=o(1)\).

\begin{proposition}[Observable consequence of the controlled trajectory]
Let \(m>0\) and let the source pair coalesce after \(J_*+2\).
On tori with side tending to infinity, the two-point
function of the rescaled field tends to the infinite-volume
source sum \(q_\infty(a,c)\) if the finite-scale maps converge
locally to their infinite-volume versions. This sum tends to zero
as \(|a-c|\) tends to infinity. If the trajectory and covariance
data tend locally along a massive sequence to massless data, then
the limits \(q_\infty(a,c)\) converge at each fixed pair to the
massless source sum. When the latter is a translation-invariant
two-point function and the massless Gaussian covariance \(C\)
satisfies \(C(x)\sim C_0|x|^{-2}\) for \(C_0>0\), it has the
asymptotic form \(q_\infty(a,c)\sim \lambda^2 C_0|a-c|^{-2}\)
for a constant \(\lambda>0\).
\end{proposition}

\begin{proof}
Use two sources linear in \(\psi\), both initially of coefficient
one. In the physical \(u\) coordinate their insertion is
\(\exp[\sigma_a(R/A_{\rm sc})u_a+\sigma_c(R/A_{\rm sc})u_c]\).
Write \(\lambda_{a,j}\) for the one-source coefficient, and
\(\widehat\lambda_{a,j}\) for its value after the small-set
extraction \(Q_j\). The graded structure of the exact RG maps
means the projection to \(\mathbb R[\sigma_a]/(\sigma_a^2)\)
commutes with the maps. To estimate that projection at any
finite scale \(j\), place the second source at distance more
than \(48s_j\) and set its coefficient zero. Thus the
fluctuation source weight in that sector is \(x_j s_j\) for
every step under consideration. Because a linear field has
fluctuation radius comparable to \(s_j^{-1}\), the localization
bound at zero gives
\[
 |\widehat\lambda_{a,j}-\lambda_{a,j}|
       \le C f_j/x_j.
\]
The same bound holds for the actual pair of sources, since the
one-source map is independent of the second source. This use of
an unused source avoids the decreased fluctuation weight for a
pair after its coalescence.

Put \(w_j^{(1)}=\sum_y\sum_{i\le j}\Gamma_i(y)\) and define
\[
 d_j=(\nu_j+\eta'_jg_j^{\rm old})w_{j+1}^{(1)}
                                      -\nu_j w_j^{(1)}.
\]
The one-source perturbative formula
\cite[(A.7)--(A.10)]{ParkTorus}, applied to \(V_j-Q_j\),
gives
\[
 \lambda_{a,j+1}=(1-d_j)\lambda_{a,j}+e_{a,j},
 \qquad |e_{a,j}|\le Cf_j/x_j
\]
as long as \(|\lambda_{a,j}|\) stays in a fixed bounded
interval. In fact the bulk extraction changes \(\nu_j,g_j\)
by at most \(Cf_j/s_j^2,Cf_j\),
respectively, so it changes \(d_j\) by at most \(Cf_j\).
One can also obtain the one-source perturbative formula by
contracting the linear source with the quadratic bulk term
after the quartic tadpole has been added.

The definition of \(d_j\) telescopes:
\[
 \sum_{j=k}^{n-1}d_j
  =\nu_nw_n^{(1)}-\nu_kw_k^{(1)}
    +\sum_{j=k}^{n-1}(\nu_j+\eta'_jg_j^{\rm old}-\nu_{j+1})
       w_{j+1}^{(1)}.
\]
The right-hand side has a limit as \(n\) tends to infinity.
Indeed the boundary term at \(n\) vanishes, and the remaining
sum has absolute bound
\(C\sum_{j\ge k}(\chi_jx_j^2+f_j)\).
The support and slice bounds, together with the mass estimate,
also give \(|d_j|\le C\chi_jx_j\), hence
\(\sum_{j\ge k}d_j^2\le C\sum_{j\ge k}\chi_jx_j^2\).
Finally
\(\sum_{j\ge k}|e_{a,j}|\le C\sum_{j\ge k}f_j/x_j\).
The reciprocal comparison for \(x_j\) bounds the sums at \(k=0\)
by \(C(x_0/\lambda_{\rm cov}^2+U/x_0)\). They vanish as \(k\) tends to
infinity, uniformly in \(a\), and uniformly as the mass cutoff
tends to infinity. For example,
\(\sum_{j\ge k}\chi_jx_j^2\le C\chi_k x_k/\lambda_{\rm cov}^2\), by summing
\(x_j^2\) before the cutoff and the geometric tail after it.
The factors \(\prod_{i<j}(1-d_i)\) are uniformly close to one
and converge, by the expansion of \(\log(1-d_i)\) and the
telescoping estimate. Variation of constants gives the same
conclusions for \(\lambda_{a,j}\), with a limit
\(\lambda_{a,\infty}\) uniformly bounded above and bounded away
from zero. This also closes the bounded-interval assumption by
induction. In particular \(\widehat\lambda_{a,j}\) converges to
the same limit, uniformly in \(a\).

Let \(k\) be the coalescence scale of \(a,c\). The mixed constant
in the effective potential is extracted from small polymers. Its
localized density is supported at the marked sites, so summing it
over the volume introduces no volume factor.
The part from \(Q_j\) is zero for \(j<k\) and has absolute
value at most \(Cf_j/(x_j^2s_k^2 4^{j-k})\) for
\(j\ge k\), by the two-source fluctuation weight in (2.19) of
\cite{ParkMap}. The remaining part is exactly the Gaussian
contraction of the two linear sources. The mixed part of the
interior \(W_j\) is zero because this contraction is constant
and is removed by localization. These claims also follow
directly from the mixed specialization in \cite[Lemma~8.4]{ParkTorus}.
With the sign convention that \(q_\infty\) is the coefficient
in the log partition function, we obtain
\[
 q_\infty(a,c)=\sum_{j\ge0}\widehat\lambda_{a,j}
            \widehat\lambda_{c,j}\Gamma_{j+1}(a-c)+E(a,c),
 \qquad
 |E(a,c)|\le \frac{C}{s_k^2}
 \left(\chi_k^{3/2}x_k+\frac{Ud_0^k}{x_k^2}\right).
\]
For the last inequality use the monotonicity of \(\chi_j,x_j\)
and \(x_k/x_j\le C(1+\lambda_{\rm cov}^2x_0(j-k))\) in the sum. Covariance
support makes the first sum start no earlier than \(k-O(1)\).
Together with \(\|\Gamma_{j+1}\|_\infty\le C\lambda_{\rm cov}\chi_j s_j^{-2}\),
this proves decay at positive mass. At zero mass the error is
\(o(s_k^{-2})\) as \(k\) tends to infinity. Replacing the source
coefficients in the covariance sum by their limits also costs
\(o(s_k^{-2})\), uniformly in the sites.

For the thermodynamic limit, stop the exact local torus flow
at \(N-1\) and use the direct Gaussian integration of the
remaining covariance tail. The periodized tail derivative
bounds and the fixed number \(L^4\) of remaining blocks give
\[
 Z_{\rm tail,\varnothing}(0)\ge\tfrac12,\qquad
 \|Z_{\rm tail}\|_{T_{\ell_{N-1}}(0)}\le C_L .
\]
Here \(Z_{\rm tail}=\mathbb E_{C_{\rm tail}}\theta
(I_{N-1}\circ K_{N-1})\) is the tail integral after removal of
the accumulated field-independent factor. These bounds follow
from the finite circle product and the regulator moments at
the tail covariance as in the finite-volume reduction. They
do not require a finite-range property for that covariance.
The one-source coefficients at zero vanish by spin-flip
symmetry. Thus the additional mixed contribution to the
log partition function is \(O_L(\ell_{N-1,\sigma\sigma}^{-1})\),
which tends to zero at each fixed positive mass and fixed
pair, because \(x_{N-1}\) remains bounded below and the
mixed source weight grows like \(4^{N-1-k}\).
The finite-scale increments converge by locality. At each
fixed pair their tails are uniformly summable, including
as the mass goes to zero. For instance the mixed remainder
at step \(j\ge k\) is dominated by a constant times
\(4^{-(j-k)}(1+\lambda_{\rm cov}^2x_0j)^2\). This proves
both the thermodynamic limit and the claimed fixed-pair
continuity of the source sum. The near pairs are not needed
in these arguments. The susceptibility comparison uses the
bulk uniform-source derivative of the same last Gaussian
integral, independently of the two-source flow.

The massless estimate gives
\(q_\infty(a,c)=\lambda_{a,\infty}\lambda_{c,\infty}C(a-c)
+o(|a-c|^{-2})\), uniformly in the sites. Translation invariance
then gives
\(\lambda_{a+x,\infty}/\lambda_{x,\infty}
\to\lambda_{0,\infty}/\lambda_{a,\infty}\) as \(|x|\) tends
to infinity with \(a\) fixed. Iterating for any fixed positive
integer \(p\) gives the \(p\)-th power of this limit for the
ratio with \(a\) replaced by \(pa\). Both numerator and
denominator stay uniformly bounded and positive, so
\(\lambda_{a,\infty}=\lambda_{0,\infty}\) for every site.
\end{proof}


\section{Critical two-point function for the box interaction}
\label{sec:physical-conclusion}



\begin{proof}[Proof of Theorem~\ref{thm:main}]
At fixed \(m^2>0\) use Proposition~\ref{prop:physical-matching}
on tori tending to infinity. The lifted reference Hessians
converge locally and in the weaker weighted kernel norm
of Section~\ref{sec:finitevolume}. Choose a subsequence on
which the tuned \(\beta_{m,N},z_{m,N}\) converge. The
initial local activities converge on each fixed polymer
by the finite exact forest formulas. Since each fixed-scale
step depends only on finitely many local polymer inputs,
its coordinates converge by the continuity of the Gaussian
step in the covariance and the activities. Thus the
controlled coordinate bounds pass to the corresponding
infinite-volume local trajectory. The direct final Gaussian
estimate and the endpoint argument of
Section~\ref{sec:finitevolume} identify the limiting spin
law as the infinite-volume subcritical law at \(\beta_m\)
and give \eqref{eq:physical_suscept}. In particular
\(\beta_m\to\beta_c^D\) as \(m^2\downarrow0\).

Along a sequence of such positive masses tending to zero
choose a subsequence with \(z_m\to z_*\). The cavity Hessian,
the finite-range pieces at each fixed scale, and the local
initial activities converge again. Hence every fixed-scale
coordinate, including the sources, converges. The source-sum
estimates in the observable analysis allow this limit in the
two-point function at every fixed pair to which the pair
flow applies. This suffices for all sufficiently distant pairs
at the fixed \(R\). Independently, the critical free spin
correlations are the limits of the subcritical ones at
\(\beta_m\). The auxiliary Gaussian conditional law then
identifies the limiting source sum with the translation
invariant two-point function of \(\psi\) at those pairs.

At the massless endpoint the absorbed denominator has the
expansion
\[
 \widehat T(k)=t_2 R^2|k|^2+O_R(|k|^4),\qquad t_2>0.
\]
The constant is positive by the uniform quadratic multiplier
comparison. The scalar form of the quadratic part follows from
hypercubic symmetry. The kernels in \(T\) have exponential
moments at fixed \(R\), so the expansion and its differentiated
bounds follow from Taylor's theorem. Since \(\widehat A(0)=\beta_c^D p\),
the massless covariance in the \(\psi\) normalization satisfies
\[
 \widehat C(k)=\lambda R^2\widehat A(k)/\widehat T(k)
       =\frac{\lambda\beta_c^D p}{t_2|k|^2}+O_R(1)
\]
near zero. Subtract the first term with a smooth cutoff
equal to one near zero. The remainder has weak derivatives
of order at most two in \(L^1\) on the frequency torus:
a derivative of order \(r\le2\) is bounded near zero by
\(C_R|k|^{-r}\). This follows by differentiating the symbol
expansions. The boundary terms on spheres shrinking to zero
vanish, so the classical derivatives away from zero are
also the weak derivatives. Integration by parts in a coordinate
with maximal \(|x_i|\), followed by the Riemann--Lebesgue
lemma, gives a Fourier coefficient \(o(|x|^{-2})\) for the
remainder. The inverse transform of the cutoff singular term
has leading coefficient \(1/(4\pi^2)\) times its prefactor,
as follows from the fundamental solution of the Laplacian
in dimension four and convolution with the Schwartz kernel
of the cutoff. Therefore
\[
 C(x)\sim \frac{\lambda\beta_c^D p}{4\pi^2 t_2}|x|^{-2}.
\]
By the observable estimates, the interacting auxiliary
covariance has this leading term times \(\lambda_\infty^2\),
where \(\lambda_\infty>0\) is the limiting one-source
coefficient. The observable argument uses the translation
invariance of the physical covariance to make this coefficient
independent of the mark's location in the block grid.

Finally put \(P_*=BA^{-1}\), with both kernels evaluated at the
critical parameter. In the original joint representation
\(h\) conditional on \(\sigma\) is Gaussian of mean \(A\sigma\)
and covariance \(A\). Since \(\psi=\sqrt\lambda\,R B^{-1}h\),
it gives the exact relation
\[
 G_\sigma=(\lambda R^2)^{-1}P_* * G_\psi *P_* - A^{-1}.
\]
Here \(A^{-1}\) and \(P_*\) have exponentially summable kernels
at fixed \(R\). Indeed \(\widehat A\) is a trigonometric polynomial
strictly positive on the real frequency torus, so its reciprocal
is analytic in a complex neighborhood. Also
\(\sum_xP_*(x)=(1-\beta_c^D q)/(\beta_c^D p)\).
For convolution with either exponentially summable kernel,
the leading \(|x|^{-2}\) term is multiplied by the kernel's
sum. To verify this, split the sum at displacement \(|x|/2\),
use dominated convergence after multiplying by \(|x|^2\) on the
inner part, and use the exponential tail and boundedness of
the covariance on the outer part. The latter boundedness
also follows directly from the conditional Gaussian identity.
It follows that
\[
 C_R=\frac{\lambda_\infty^2(1-\beta_c^D q)^2}
             {4\pi^2 t_2\beta_c^D p R^2}>0 .
\]
Finally, the diagonal normalization gives, in every finite volume,
\[
 \langle\sigma_0\sigma_x\rangle^f_{D,\beta}
 =\langle\sigma_0\sigma_x\rangle^f_{D_R,\beta/(1+\delta_R)},
 \qquad \delta_R=P_0(0).
\]
The same identity holds after the free infinite-volume limit.
Since \(\beta_c^D=(1+\delta_R)\beta_c^{\rm box}\), the endpoint
just identified is exactly the critical free state in
Theorem~\ref{thm:main}. The positive constant above is therefore
the asserted \(C_R\) for the original box model.
\end{proof}


\begin{thebibliography}{99}

\bibitem{AR}
A. Abdesselam and V. Rivasseau.
\newblock Trees, forests and jungles: a botanical garden for cluster expansions.
\newblock In \emph{Constructive Physics}, Lecture Notes in Physics 446,
7--36 (1995).
\newblock \href{https://arxiv.org/abs/hep-th/9409094}{arXiv:hep-th/9409094}.


\bibitem{Bauerschmidt}
R. Bauerschmidt.
\newblock A simple method for finite range decomposition of quadratic forms
and Gaussian fields.
\newblock \emph{Probab. Theory Relat. Fields}, 157, 817--845 (2013).
\newblock \href{https://arxiv.org/abs/1206.2212}{arXiv:1206.2212}.

\bibitem{BPR}
R. Bauerschmidt, J. Park, and P.-F. Rodriguez.
\newblock The Discrete Gaussian model, I. Renormalisation group flow at high temperature.
\newblock \href{https://arxiv.org/abs/2202.02286}{arXiv:2202.02286}.

\bibitem{ParkMap}
J. Park.
\newblock A Renormalisation Group Map for Short- and Long-ranged Weakly
Coupled \(\lvert\varphi\rvert^4\) Models in \(d\ge4\) at and Above the Critical Point.
\newblock \href{https://arxiv.org/html/2511.03495v2}{arXiv:2511.03495v2} (2025).

\bibitem{ParkTorus}
J. Park.
\newblock Torus scaling limits and the plateau of the critical weakly
coupled \(\lvert\varphi\rvert^4\) model in \(d\ge4\).
\newblock \href{https://arxiv.org/html/2511.06321v1}{arXiv:2511.06321v1} (2025).

\bibitem{ST}
G. Slade and A. Tomberg.
\newblock Critical correlation functions for the 4-dimensional weakly
self-avoiding walk and \(n\)-component \(\lvert\varphi\rvert^4\) model.
\newblock \emph{Commun. Math. Phys.}, 342, 675--737 (2016).
\newblock \href{https://arxiv.org/abs/1412.2668}{arXiv:1412.2668}.


\bibitem{BBSModel}
R. Bauerschmidt, D. C. Brydges, and G. Slade.
\newblock Scaling limits and critical behaviour of the 4-dimensional
\(n\)-component \(\lvert\varphi\rvert^4\) spin model.
\newblock \emph{J. Stat. Phys.} 157, 692--742 (2014).
\newblock \href{https://arxiv.org/abs/1403.7424}{arXiv:1403.7424}.

\bibitem{BBSIII}
R. Bauerschmidt, D. C. Brydges, and G. Slade.
\newblock A renormalisation group method. III. Perturbative analysis.
\newblock \href{https://arxiv.org/abs/1403.7252}{arXiv:1403.7252}.

\bibitem{DCT}
H. Duminil-Copin and V. Tassion.
\newblock A new proof of the sharpness of the phase transition for
Bernoulli percolation and the Ising model.
\newblock \emph{Commun. Math. Phys.} 343, 725--745 (2016).
\newblock \href{https://arxiv.org/abs/1502.03050}{arXiv:1502.03050}.
\end{thebibliography}

\subsection*{AI Statement}

This manuscript is generated using ChatGPT 6.1 Sol.

\end{document}
